% =================================================
% 性学系列 · 性健康基础
% 共享基座见 common.tex
% =================================================
\documentclass[12pt,UTF8,openany]{ctexbook}
% =================================================
% common.tex —— 性学系列 共享基座
% 导言区 / 配色 / TikZ 库 / ctexset / C 列型 / 字体 / 书眉页码
% 各册 % =================================================
% common.tex —— 性学系列 共享基座
% 导言区 / 配色 / TikZ 库 / ctexset / C 列型 / 字体 / 书眉页码
% 各册 % =================================================
% common.tex —— 性学系列 共享基座
% 导言区 / 配色 / TikZ 库 / ctexset / C 列型 / 字体 / 书眉页码
% 各册 \input{common.tex} 引用
% =================================================
\usepackage[top=2.5cm,bottom=2.5cm,left=2.6cm,right=2.6cm,headheight=2.0cm]{geometry}
\usepackage{amsmath}
\usepackage{amssymb}
\usepackage{array}
\usepackage{tabularx}
\usepackage{multirow}
\usepackage{booktabs}
\usepackage{enumitem}
\usepackage{xcolor}
\usepackage{graphicx}
\usepackage{tikz}
\usetikzlibrary{positioning,calc,arrows.meta,shapes.geometric,matrix}
\definecolor{myblue}{HTML}{185FA5}
\definecolor{myblueL}{HTML}{E6F1FB}
\definecolor{mygreen}{HTML}{3B6D11}
\definecolor{mygreenL}{HTML}{C0DD97}
\definecolor{myteal}{HTML}{1D9E75}
\definecolor{myamber}{HTML}{854F0B}
\definecolor{myamberL}{HTML}{FAEEDA}
\definecolor{myred}{HTML}{A32D2D}
\definecolor{myredL}{HTML}{FCEBEB}
\definecolor{mypurple}{HTML}{534AB7}
\definecolor{mypurpleL}{HTML}{EEEDFE}
\definecolor{mycoral}{HTML}{D85A30}
\definecolor{mycoralL}{HTML}{FAECE7}
\definecolor{mypink}{HTML}{D4537E}
\definecolor{mypinkL}{HTML}{F4C0D1}
\definecolor{mypinkM}{HTML}{ED93B1}
\definecolor{mygray}{HTML}{888780}

% 字体
\setCJKfamilyfont{hei}{SimHei}
\newcommand{\myhei}{\CJKfamily{hei}}

% 书眉与页码（印刷友好：中文页眉 + 居中页码）
\usepackage{fancyhdr}
\pagestyle{fancy}
\fancyhf{}
\fancyhead[L]{\small\myhei 性学系列}
\fancyhead[R]{\small\leftmark}
\fancyfoot[C]{\small\thepage}
\renewcommand{\headrulewidth}{0.4pt}
\fancypagestyle{plain}{\fancyhf{}\fancyfoot[C]{\small\thepage}\renewcommand{\headrulewidth}{0pt}}
\fancypagestyle{empty}{\fancyhf{}\renewcommand{\headrulewidth}{0pt}}

% 行距（零基础易读，略松）
\linespread{1.3}

% 超链接：印刷友好（黑色，电子版仍可点击）
\usepackage[colorlinks,linkcolor=black,citecolor=black,urlcolor=black]{hyperref}

\ctexset{
  chapter/name={第,章},
  chapter/number={\arabic{chapter}},
  section/format=\Large\bfseries,
}

\newcolumntype{C}{>{\centering\arraybackslash}X}
 引用
% =================================================
\usepackage[top=2.5cm,bottom=2.5cm,left=2.6cm,right=2.6cm,headheight=2.0cm]{geometry}
\usepackage{amsmath}
\usepackage{amssymb}
\usepackage{array}
\usepackage{tabularx}
\usepackage{multirow}
\usepackage{booktabs}
\usepackage{enumitem}
\usepackage{xcolor}
\usepackage{graphicx}
\usepackage{tikz}
\usetikzlibrary{positioning,calc,arrows.meta,shapes.geometric,matrix}
\definecolor{myblue}{HTML}{185FA5}
\definecolor{myblueL}{HTML}{E6F1FB}
\definecolor{mygreen}{HTML}{3B6D11}
\definecolor{mygreenL}{HTML}{C0DD97}
\definecolor{myteal}{HTML}{1D9E75}
\definecolor{myamber}{HTML}{854F0B}
\definecolor{myamberL}{HTML}{FAEEDA}
\definecolor{myred}{HTML}{A32D2D}
\definecolor{myredL}{HTML}{FCEBEB}
\definecolor{mypurple}{HTML}{534AB7}
\definecolor{mypurpleL}{HTML}{EEEDFE}
\definecolor{mycoral}{HTML}{D85A30}
\definecolor{mycoralL}{HTML}{FAECE7}
\definecolor{mypink}{HTML}{D4537E}
\definecolor{mypinkL}{HTML}{F4C0D1}
\definecolor{mypinkM}{HTML}{ED93B1}
\definecolor{mygray}{HTML}{888780}

% 字体
\setCJKfamilyfont{hei}{SimHei}
\newcommand{\myhei}{\CJKfamily{hei}}

% 书眉与页码（印刷友好：中文页眉 + 居中页码）
\usepackage{fancyhdr}
\pagestyle{fancy}
\fancyhf{}
\fancyhead[L]{\small\myhei 性学系列}
\fancyhead[R]{\small\leftmark}
\fancyfoot[C]{\small\thepage}
\renewcommand{\headrulewidth}{0.4pt}
\fancypagestyle{plain}{\fancyhf{}\fancyfoot[C]{\small\thepage}\renewcommand{\headrulewidth}{0pt}}
\fancypagestyle{empty}{\fancyhf{}\renewcommand{\headrulewidth}{0pt}}

% 行距（零基础易读，略松）
\linespread{1.3}

% 超链接：印刷友好（黑色，电子版仍可点击）
\usepackage[colorlinks,linkcolor=black,citecolor=black,urlcolor=black]{hyperref}

\ctexset{
  chapter/name={第,章},
  chapter/number={\arabic{chapter}},
  section/format=\Large\bfseries,
}

\newcolumntype{C}{>{\centering\arraybackslash}X}
 引用
% =================================================
\usepackage[top=2.5cm,bottom=2.5cm,left=2.6cm,right=2.6cm,headheight=2.0cm]{geometry}
\usepackage{amsmath}
\usepackage{amssymb}
\usepackage{array}
\usepackage{tabularx}
\usepackage{multirow}
\usepackage{booktabs}
\usepackage{enumitem}
\usepackage{xcolor}
\usepackage{graphicx}
\usepackage{tikz}
\usetikzlibrary{positioning,calc,arrows.meta,shapes.geometric,matrix}
\definecolor{myblue}{HTML}{185FA5}
\definecolor{myblueL}{HTML}{E6F1FB}
\definecolor{mygreen}{HTML}{3B6D11}
\definecolor{mygreenL}{HTML}{C0DD97}
\definecolor{myteal}{HTML}{1D9E75}
\definecolor{myamber}{HTML}{854F0B}
\definecolor{myamberL}{HTML}{FAEEDA}
\definecolor{myred}{HTML}{A32D2D}
\definecolor{myredL}{HTML}{FCEBEB}
\definecolor{mypurple}{HTML}{534AB7}
\definecolor{mypurpleL}{HTML}{EEEDFE}
\definecolor{mycoral}{HTML}{D85A30}
\definecolor{mycoralL}{HTML}{FAECE7}
\definecolor{mypink}{HTML}{D4537E}
\definecolor{mypinkL}{HTML}{F4C0D1}
\definecolor{mypinkM}{HTML}{ED93B1}
\definecolor{mygray}{HTML}{888780}

% 字体
\setCJKfamilyfont{hei}{SimHei}
\newcommand{\myhei}{\CJKfamily{hei}}

% 书眉与页码（印刷友好：中文页眉 + 居中页码）
\usepackage{fancyhdr}
\pagestyle{fancy}
\fancyhf{}
\fancyhead[L]{\small\myhei 性学系列}
\fancyhead[R]{\small\leftmark}
\fancyfoot[C]{\small\thepage}
\renewcommand{\headrulewidth}{0.4pt}
\fancypagestyle{plain}{\fancyhf{}\fancyfoot[C]{\small\thepage}\renewcommand{\headrulewidth}{0pt}}
\fancypagestyle{empty}{\fancyhf{}\renewcommand{\headrulewidth}{0pt}}

% 行距（零基础易读，略松）
\linespread{1.3}

% 超链接：印刷友好（黑色，电子版仍可点击）
\usepackage[colorlinks,linkcolor=black,citecolor=black,urlcolor=black]{hyperref}

\ctexset{
  chapter/name={第,章},
  chapter/number={\arabic{chapter}},
  section/format=\Large\bfseries,
}

\newcolumntype{C}{>{\centering\arraybackslash}X}


\begin{document}

% ===== 封面 =====
\begin{titlepage}
\centering
\vspace*{1.5cm}
\begin{tikzpicture}
\draw[myblue,line width=2pt] (0,0) circle (1.2);
\node[text=white] at (0,0) {\myhei\zihao{3}性};
\end{tikzpicture}
\vspace{0.6cm}
\noindent\rule{\textwidth}{0.4pt}
\vspace{0.5cm}
\myhei\zihao{5} 性学系列\,·\,家庭健康科普\,·\,全五册
\vspace{2.5cm}
\myhei\zihao{0} 性健康基础\par
\vspace{1.2cm}
\itshape\zihao{3} 认知·同意·防护——先建立安全底线\par
\vspace{5.2cm}
\noindent\rule{\textwidth}{0.4pt}
\vspace{0.5cm}
\zihao{5} 性学系列\,·\,\today
\end{titlepage}

% ===== 免责声明 =====
\begin{center}\zihao{5}\itshape
本书为家庭健康科普读物，内容力求科学严谨，但\textbf{不替代专业医生的诊断与治疗}。\\
涉及个人健康问题，请以正规医疗机构的检查与建议为准。\\
全书以安全、自愿、相互尊重为底色，尊重多元与个体差异。\par\end{center}
\clearpage

% ===== 如何使用本书 =====
\chapter*{如何使用本书}
\addcontentsline{toc}{chapter}{如何使用本书}
\noindent\textbf{性学系列共五册：}
\begin{itemize}[leftmargin=2em]
\item \textbf{A 性健康基础} —— 认知、同意、防护、STI、避孕总论（建议先读）
\item \textbf{B 亲密实践与技巧} —— 前戏、体位、高潮引导、呼吸与盆底
\item \textbf{C 男性性健康} —— 男性解剖、ED/自慰、绝育、不育
\item \textbf{D 女性性健康} —— 女性解剖、妇科、不孕、人流、避孕、孕期/更年期
\item \textbf{E BDSM 与 Kink} —— 安全框架下的亲密表达
\end{itemize}
\vspace{0.6em}
\noindent\textbf{建议阅读路径（零基础）：} 先通读 A 建立认知与安全底线 → 按需读 C 或 D 了解自己的身体 → 再读 B 学技巧 → 兴趣到了再读 E。每章末附有本章小结，便于快速回顾。
\clearpage

% ===== 目录 =====
\tableofcontents
\clearpage
\chapter{为什么要读这本书}

对性好奇、想事先弄明白，是非常正常且负责任的表现。真正危险的不是"懂"，
而是"似懂非懂就行动"。性健康是整体健康的一部分，它关乎身体安全、
心理成熟与人际尊重。

\section{一个核心判断}

本书反复强调一条主线，可以浓缩为六个字：

\begin{quote}
\centering \myhei\zihao{3}\textcolor{myblue}{安全 · 自愿 · 知情}
\end{quote}

\begin{itemize}
  \item \textbf{安全}：采取科学措施，避免意外怀孕与性传播感染。
  \item \textbf{自愿}：任何行为都建立在双方清醒、自愿的前提之上。
  \item \textbf{知情}：了解足够的知识，在明白后果的前提下做决定。
\end{itemize}

\section{本书与专业教材的关系}

本书以"零基础可读"为基调。其中“性医学基础：医学维度的速查”一章偏专业，
系统速查生殖解剖、性激素、性功能障碍等医学内容，给"想了解医学底层"的读者一个入口；
其余各章则从"我该了解什么、该准备什么"的角度切入。若需更深的专业内容，
请参考正规性医学／男科／妇科教材。

\chapter{五个方面的准备}

面对"初次接触性，我该准备什么"这个问题，可以把它拆成五个相互支撑的方面。
先建立整体框架，再逐一展开。

\begin{table}[htbp]
  \centering
  \caption{性健康准备的五个方面}
  \begin{tabularx}{\textwidth}{@{}l l X@{}}
    \toprule
    \textbf{方面} & \textbf{关键词} & \textbf{要解决的核心问题} \\
    \midrule
    认知准备   & 了解       & 懂基本的生理结构与生殖常识，不被谣言误导 \\
    安全防护   & 保护       & 学会正确使用安全套，懂得避孕与防病 \\
    沟通与同意 & 尊重       & 明确意愿与边界，学会和伴侣沟通 \\
    健康维护   & 预防       & 疫苗、筛查、卫生，把风险挡在前面 \\
    心态准备   & 从容       & 不攀比、不勉强，尊重彼此节奏 \\
    \bottomrule
  \end{tabularx}
\end{table}

以下各章按这五个方面逐一讲述。

\chapter{认知准备：先懂基础}

\begin{center}
\begin{tikzpicture}[scale=0.85,
  organ/.style={fill=mypinkL,draw=mypink,thin},
  lead/.style={mygray,thin,dashed},
  lab/.style={font=\footnotesize}
]
  % 内生殖器（左）
  \draw[organ] (1,2.2) ellipse (0.35 and 0.5);
  \draw[organ] (4.5,2.2) ellipse (0.35 and 0.5);
  \draw[mypink,thick] (2.3,2.4) to[out=150,in=100] (1.4,2.3);
  \draw[mypink,thick] (3.7,2.4) to[out=30,in=80] (4.6,2.3);
  \fill[mypinkL,draw=mypink] (3,2.5) .. controls (2.6,2.2) and (2.6,1.7) .. (3,1.5) .. controls (3.4,1.7) and (3.4,2.2) .. (3,2.5);
  \fill[mypinkL,draw=mypink] (2.9,1.45) -- (2.85,0.5) -- (3.15,0.5) -- (3.1,1.45) -- cycle;
  \draw[lead] (1,2.6) -- (0.4,3) node[lab,left]{卵巢};
  \draw[lead] (2,2.5) -- (1.3,3.2) node[lab,above]{输卵管};
  \draw[lead] (3.4,2) -- (4.4,2.7) node[lab,right]{子宫};
  \draw[lead] (3.1,1) -- (4.4,0.8) node[lab,right]{阴道};
  % 外生殖器（右）
  \begin{scope}[shift={(7,0)}]
    \draw[mypink,thick] (0.5,2.6) to[out=-95,in=95] (0.5,0.6);
    \draw[mypink,thick] (2.5,2.6) to[out=-85,in=85] (2.5,0.6);
    \draw[mypinkM,thick] (0.9,2.5) to[out=-90,in=90] (0.9,0.8);
    \draw[mypinkM,thick] (2.1,2.5) to[out=-90,in=90] (2.1,0.8);
    \fill[mypinkM,draw=mypink] (1.5,2.7) circle (0.15);
    \fill[mypink] (1.5,2.1) circle (0.08);
    \draw[organ] (1.5,1.6) ellipse (0.3 and 0.2);
    \draw[lead] (1.5,2.85) -- (1.5,3.3) node[lab,above,mypink]{阴蒂};
    \draw[lead] (0.5,2) -- (-0.2,2.2) node[lab,left]{大阴唇};
    \draw[lead] (1.1,1.9) -- (0.2,1.9) node[lab,left]{小阴唇};
    \draw[lead] (1.6,2.1) -- (2.6,2.3) node[lab,right]{尿道口};
    \draw[lead] (1.8,1.6) -- (2.7,1.5) node[lab,right]{阴道口};
  \end{scope}
\end{tikzpicture}
\par\vspace{0.5em}\noindent\footnotesize 图：女性生殖系统示意（左：内生殖器；右：外生殖器）。科普示意，非精确解剖。
\end{center}

\section{了解身体与生命周期}

无论男女，了解自己的生殖系统结构与基本功能是第一步。建议通过正规医学
教材、权威健康网站或医院科普资料学习，理解以下内容：

\begin{itemize}
  \item 男性和女性外生殖器、内生殖器的名称与基本功能。
  \item 女性月经周期与排卵的大致规律，以及"安全期"并非绝对安全的原因。
  \item 精液、阴道分泌物等正常生理现象，以及哪些变化属于异常、需要就医。
\end{itemize}

\section{怀孕是怎么发生的}

怀孕需要精子与卵子结合（受精）并在子宫着床。理解这一点，就能理解各种
避孕方法的原理：有的是阻挡精子（屏障法），有的是抑制排卵（激素法），
有的是阻止着床。方法不同，但目标一致——在"不打算要孩子"时避免意外。

\section{常见误区}

\subsection*{误区一："体外排精"很安全}
失败率较高。据美国 CDC 数据，体外排精在典型使用下第一年失败率约 22\%，
完美使用（每次都正确、全程不遗漏）约 4\%——差距之大，正说明"时点控制"
极依赖当事人的经验与自觉。原因有二：射精前分泌物中可能已含有精子，
且撤退时点难以精确把握。

\subsection*{误区二："第一次不会怀孕"}
只要有排卵和无保护的行为，就有怀孕的可能，与第几次无关。

\subsection*{误区三："看起来健康就没病"}
许多性传播感染在早期没有明显症状，无法靠外表判断，只能通过检测确知。

\chapter{性医学基础：医学维度的速查}
\label{chap:medbase}

\begin{center}
\myhei\large 快速看懂性医学骨架：解剖、激素、常见障碍、检查路径——遇到问题时的"指南"
\end{center}

这一章把性医学的\textbf{医学框架}压成一页速查。它不逐病详解（那在 C 册、D 册和各专章），而是给你一张"遇到什么问题、大致是什么、该往哪走"的地图。

\section{两性生殖解剖速查}

\begin{tabularx}{\textwidth}{p{2.2cm}X}
\toprule
& 关键结构 \\
\midrule
男性 & 外生殖器（阴茎/阴囊）＋内生殖器：睾丸（产精＋分泌睾酮）、附睾、输精管、精囊（约 60\% 精液）、前列腺（约 30\%）、尿道球腺；精液中精子 $<$5\% \\
女性 & 外阴（大/小阴唇/阴蒂[唯一以快感为目的的器官]/阴道前庭）＋内生殖器：阴道、子宫、输卵管（受精场所）、卵巢（排卵＋分泌激素） \\
性腺轴 & 下丘脑（GnRH） $\to$ 垂体（FSH/LH） $\to$ 性腺（睾酮/雌孕激素）——调控整个生殖系统 \\
\bottomrule
\end{tabularx}

\noindent 示意图见本册“身体与性：男女构造、敏感点与荷尔蒙”一章；专科细节见 C 册“男性生殖系统”、D 册“女性生殖系统”。

\section{性激素速查}

\begin{tabularx}{\textwidth}{lX}
\toprule
激素 & 功能 \\
\midrule
睾酮 (Testosterone) & 精子发生、维持性欲、肌肉量、骨密度；来源睾丸（95\%）＋肾上腺 \\
双氢睾酮 (DHT) & 前列腺发育；过高与雄激素性脱发相关 \\
雌二醇 (E2) & 子宫内膜增生、第二性征、骨骼健康；来源卵巢颗粒细胞 \\
孕酮 (P4) & 维持妊娠、子宫内膜分泌期改变；来源黄体、胎盘 \\
\bottomrule
\end{tabularx}

\noindent 随年龄变化：\textbf{男性}睾酮约 12 岁起升、20--30 岁达峰、40 岁后每年约降 1\%（迟发性性腺功能减退 LOH）；\textbf{女性}雌激素青春期起升、育龄期周期波动、围绝经期（45--55 岁）急降、绝经后低水平维持。

\section{性反应周期（速查）}

Masters \& Johnson 四期模型：\textbf{兴奋 $\to$ 平台 $\to$ 高潮 $\to$ 消退}。要点：男性高潮后进入不应期（随年龄延长）；女性消退慢、可再次高潮。唤起是渐进过程，女性往往需要更久——\textbf{这是正常差异，不是"冷淡"}。完整机制详见本册“亲密过程中的身体反应：性反应周期与注意事项”一章。

\section{常见性功能障碍速查}

\subsection*{男性}

\begin{tabularx}{\textwidth}{lX}
\toprule
问题 & 要点 \\
\midrule
勃起功能障碍 (ED) & 持续 $\ge$3 个月难以获得/维持勃起；多为器质性（血管性 80--90\%），糖尿病/高血压/心血管病/吸烟是危险因素；详见 C 册 \\
早泄 (PE) & 插入后 $<$1 分钟、难以控制；行为训练＋药物（达泊西汀等） \\
迟发性性腺功能减退 (LOH) & 40 岁后睾酮下降＋症状（性欲降/乏力/情绪低）；晨间总睾酮 $<$300 ng/dL 需复查 \\
\bottomrule
\end{tabularx}

\subsection*{女性}

\begin{tabularx}{\textwidth}{lX}
\toprule
问题 & 常见原因 \\
\midrule
性欲障碍 & 性兴趣/幻想持续缺乏；关系、激素、药物（SSRIs） \\
性唤起障碍 & 不能维持润滑/兴奋；绝经后雌激素缺乏、血管因素 \\
性高潮障碍 & 高潮延迟/缺失；刺激不足、心理、药物 \\
性交疼痛 & 阴道干涩、子宫内膜异位、盆底痉挛——\textbf{需排查器质性病因} \\
\bottomrule
\end{tabularx}

\section{不孕不育（速查）}

\begin{tabularx}{\textwidth}{lX}
\toprule
& 说明 \\
\midrule
定义 & 规律无保护性生活 $\ge$12 个月未孕＝不孕；能孕但反复流产＝不育 \\
病因分布 & 女性因素 $\sim$40\%、男性因素 $\sim$40\%、双方/不明 $\sim$20\% \\
基本检查 & 男性：精液常规；女性：排卵监测、输卵管通畅、盆腔超声 \\
\bottomrule
\end{tabularx}

\noindent 详见 C 册“男性不育”、D 册“不孕不育治疗”。

\section{性传播感染与避孕（指向专章）}

\begin{itemize}[leftmargin=2em,itemsep=0.3em]
  \item \textbf{STI}：可治愈类（衣原体/淋病/梅毒/滴虫）与病毒类（HIV/HSV/HPV/乙肝）——完整检测、窗口期、PEP 详见本册“性传播感染详解”一章。
  \item \textbf{避孕}：按真实可靠度分层＋决策逻辑，详见本册“避孕方法详解”一章。
\end{itemize}

\section{性心理速查}

\noindent 健康性心理四要素：\textbf{性认同}（接纳自己的性别与取向）、\textbf{性知识}、\textbf{性沟通}（开放表达需求）、\textbf{性价值观}（尊重/自愿/安全/负责）。

\begin{tabularx}{\textwidth}{lX}
\toprule
常见困扰 & 处理 \\
\midrule
性表现焦虑 & 过度关注"表现"导致功能下降；认知行为疗法＋伴侣配合 \\
性欲不一致 & 伴侣需求差异；沟通、协商、必要时求助（见 B 册“欲望不匹配”） \\
性创伤后应激 & 性侵/虐待后对性的回避与恐惧；见 A 册“性侵犯与强奸” \\
\bottomrule
\end{tabularx}

\section{就医路径速查}

\begin{tabularx}{\textwidth}{lX}
\toprule
情况 & 建议科室/路径 \\
\midrule
勃起/射精困扰、男性功能 & 泌尿外科 / 男科 \\
月经/妇科/阴道不适、女性功能 & 妇科 / 生殖医学 \\
备孕 1 年未孕 & 生殖医学中心（双方同查） \\
疑似 STI & 皮肤性病科 / 泌尿妇科，定期复查 \\
性心理困扰（焦虑/创伤/关系） & 心理科 / 正规心理咨询（性治疗师） \\
\bottomrule
\end{tabularx}

\vspace{1em}
\noindent\textbf{要点：} 性医学是生理$\times$心理$\times$社会的交叉学科；\textbf{大多数性功能障碍可有效治疗}——放下羞耻、找对科室、按时复查，比"硬扛"或"乱试"都有效。

\chapter{安全防护：最实际的一章}

\begin{center}
\begin{tikzpicture}[
  st/.style={draw=myblue,fill=myblueL,rounded corners=2pt,minimum width=1.9cm,minimum height=1.1cm,align=center,font=\small}
]
\node[st] (s1) {1\\检查\\有效期/尺寸};
\node[st,right=0.5cm of s1] (s2) {2\\小心撕开\\(别咬)};
\node[st,right=0.5cm of s2] (s3) {3\\戴在\\勃起上};
\node[st,right=0.5cm of s3] (s4) {4\\全程使用\\不翻转};
\node[st,right=0.5cm of s4] (s5) {5\\射精后\\按住退出};
\draw[-Stealth] (s1) -- (s2);
\draw[-Stealth] (s2) -- (s3);
\draw[-Stealth] (s3) -- (s4);
\draw[-Stealth] (s4) -- (s5);
\end{tikzpicture}
\par\vspace{0.5em}
\begin{tikzpicture}[font=\small]
  \node[anchor=west,font=\normalsize\bfseries] at (0,1.6){紧急避孕时间窗（越早越好）};
  \fill[mygreenL] (0,0.9) rectangle (6,1.2);
  \node[anchor=west,mygreen] at (6.2,1.05){$\le$72h\ 左炔诺孕酮};
  \fill[myblueL] (0,0.3) rectangle (10,0.6);
  \node[anchor=west,myblue] at (10.2,0.45){$\le$120h\ 乌利司坦/铜IUD(最可靠)};
  \draw[mygray] (0,-0.2) -- (10,-0.2);
  \draw[mygray] (0,-0.3) -- (0,-0.1) node[below,mygray,font=\footnotesize] at (0,-0.4){0h};
  \draw[mygray] (6,-0.3) -- (6,-0.1) node[below,mygray,font=\footnotesize] at (6,-0.4){72h};
  \draw[mygray] (10,-0.3) -- (10,-0.1) node[below,mygray,font=\footnotesize] at (10,-0.4){120h};
\end{tikzpicture}
\end{center}

这是全书最具操作性的部分。请认真读完，并落实到位。

\section{安全套：一举两得的工具}

安全套（避孕套）是\textbf{目前唯一能同时避孕和降低性传播感染风险}的常用工具。
据美国 CDC 资料，在血清状态不一致（一方感染、一方未感染）的异性恋伴侣中，
坚持使用安全套者感染 HIV 的风险比不使用者的伴侣低约 71\%--80\%；
对衣原体、淋病、乙肝、滴虫等其他感染也有明确的降低风险作用。
世界卫生组织同样指出，安全套是唯一兼具避孕与预防包括 HIV 在内的性传播感染功能的方法。

\subsection*{关于有效率的真实数字}

\begin{table}[htbp]
  \centering
  \caption{男用乳胶安全套的第一年失败率（避孕）}
  \begin{tabularx}{\textwidth}{@{}l C C@{}}
    \toprule
    \textbf{使用方式} & \textbf{第一年失败率} & \textbf{说明} \\
    \midrule
    典型使用（日常真实情况） & 约 15\% & 常因遗漏、佩戴不当所致 \\
    坚持且正确使用           & 约 2\%  & 接近理想效果 \\
    \bottomrule
  \end{tabularx}
\end{table}

可见安全套本身可靠，问题几乎总是出在"没有每次、正确地使用"上。
CDC 也指出，安全套失效多源于使用不一致或不正确，而非破裂——
乳胶安全套在性交与退出过程中的破裂率约为每一百个中两个。

\subsection*{正确使用的要点}

\begin{enumerate}[leftmargin=2em,itemsep=0.3em]
  \item 每次行为都用\textbf{新的}安全套，检查包装完整与有效期。
  \item 在生殖器发生任何接触\textbf{之前}佩戴：阴茎勃起后，捏住储精囊排出空气，再展开至根部。
  \item 全程使用，结束后在阴茎仍勃起时按住套口边缘一同退出，避免精液外溢。
  \item 用后打结或用纸巾包好丢入垃圾桶，不可重复使用。
  \item 若中途破裂或滑脱，应立即停止并更换新的安全套。
  \item 润滑剂只用\textbf{水基或硅基}；凡士林、婴儿油、护肤乳等油基会削弱乳胶、导致破裂。
  \item 对乳胶过敏者可选聚氨酯等合成材质安全套。
\end{enumerate}

\section{其他避孕方法}

\begin{table}[htbp]
  \centering
  \caption{常见避孕方法概览（ efficacy 为典型使用下的大致水平）}
  \begin{tabularx}{\textwidth}{@{}l X C@{}}
    \toprule
    \textbf{方法} & \textbf{特点} & \textbf{能否防性病} \\
    \midrule
    男用安全套   & 易获取、即时可用；需每次正确使用 & 能 \\
    女用安全套   & 由使用者自主掌握；同样为屏障法   & 能 \\
    短效口服避孕药 & 需每日按时服用；适合固定伴侣   & 不能 \\
    宫内节育器   & 长效可逆，由医生放置           & 不能 \\
    皮下埋植     & 长效可逆，由医生放置           & 不能 \\
    绝育手术     & 永久性的，适用于确定不再生育者 & 不能 \\
    安全期 / 体外排精 & 可靠性低，不推荐单独依赖 & 不能 \\
    \bottomrule
  \end{tabularx}
\end{table}

选择哪种方法，应结合健康状况、伴侣情况与个人需求，并咨询医生。
注意：\textbf{除安全套外，其他方法基本不能预防性传播感染}，
必要时需与安全套联合使用（双重保护）。

\section{万一发生了无保护行为}

\begin{itemize}
  \item \textbf{紧急避孕}：在无保护行为后应尽快（越早越好，通常在 72 小时内，
        部分方案可延长至 120 小时）咨询医生或药师使用紧急避孕药，
        它不能作为常规避孕手段。
  \item \textbf{检测}：如有感染风险，应在医生建议的时间窗后做性传播感染筛查；
        很多感染早期无症状，检测是唯一确知方式。
  \item \textbf{就医}：出现疼痛、异常分泌物、出血、发热等症状，请及时到正规医院就诊。
\end{itemize}

\chapter{沟通与同意：不可逾越的底线}

\section{什么算"同意"}

同意必须同时满足几个条件：\textbf{清醒}（未受酒精或药物影响）、
\textbf{自愿}（没有被强迫、欺骗或施压）、\textbf{具体}（针对当下这件事）、
\textbf{可撤回}（任何时候都可以改变主意）。沉默、犹豫、畏惧都不等于同意。

需要特别强调：性行为涉及法律年龄界限，各地规定不同。若你或对方未成年，
请务必先了解当地法律，并听取可信成年人的意见。

\section{如何与伴侣沟通}

 Opening the conversation 常常让人觉得尴尬，但这是成熟与负责的表现。
可以在轻松的场合，坦诚地谈：

\begin{itemize}
  \item 双方对避孕与防病的打算：是否使用安全套、是否做过检测。
  \item 彼此的意愿、顾虑与边界：什么可以做、什么不接受。
  \item 可以随时喊停的"安全词"或约定，让任何一方安心表达不适。
\end{itemize}

谈这些不会破坏气氛，反而会让双方更放松、更信任。真正尊重你的人，
会认真对待你的边界。

% =================================================
% A 册补充专题：性侵犯与强奸（认知与应对）
% 现代循证撰写，紧跟「沟通与同意」章
% =================================================

\chapter{性侵犯与强奸：认知与应对}

\section{先正名：这不是"你的错"，也不是"她没拒绝够大声"}
性侵犯（sexual assault / rape）的核心是\textbf{同意的缺失}，而不是"用了什么方式""在哪里发生"。
只要一方没有\textbf{自愿、清醒、可撤回}的同意，就是侵犯——无论发生在婚姻内、恋爱中、职场、校园，还是陌生人之间。
把它单独成章，是因为它\textbf{不是女性专利，也不是"小事"，更不是某个性行为的变体}，而是一种以权力与暴力为底色的侵犯。
理解它的第一课，就是把"受害者有罪论"请出去。

\section{常见误解（请逐条拆开）}
\begin{tabularx}{\textwidth}{l X}
\toprule
\textbf{误解} & \textbf{事实} \\
\midrule
"他没穿裙子 / 没喝醉，就不会被侵犯" & 是否穿着、是否饮酒，都与"谁在施加、谁没同意"无关。 \\
"都结婚了，谈不上强奸" & 婚内强奸在多数国家（含中国刑法修订后）已被承认为犯罪；同意不是结婚证附带的。 \\
"她没大声喊'不'，就是默许" & 沉默、僵硬（freezing）、事后才说出"不"，都可能是同意缺失的表现；同意须是\textbf{主动且持续}的。 \\
"隔了很久才说，是编的" & 创伤后记忆常是非线性的；羞耻、恐惧、关系压力都会让人延迟求助，延迟不等于不真。 \\
"只有暴力才算" & 胁迫、利用职权/信息差、事后无法撤回的同意，都算侵犯，未必有淤青。 \\
\bottomrule
\end{tabularx}

\section{当下该怎么做（第一反应清单）}
目标是\textbf{先保证安全，再处理后续}，顺序可随情况灵活调整：
\begin{enumerate}
  \item \textbf{到安全的地方}：离开现场，去能信任的人身边，或拨打 110 / 当地求助热线。
  \item \textbf{不要急着"洗掉证据"}：在未就医取证前，尽量避免洗澡、刷牙、换衣、排便——这些会冲掉唾液、精液、皮屑、纤维等物证。若不得不清洗，尽量用干净容器保存衣物。
  \item \textbf{尽快就医（最好 72 小时内）}：到\textbf{急诊或妇科/男科}，说明"可能需要性侵取证（rape kit）"。医生会检查伤情、取物证、并评估 STI 与妊娠风险。
  \item \textbf{取证与报案可分开做}：可以先体检、保留证据，之后再决定是否走法律程序；证据保存期通常有限，但\textbf{身体安全永远排在第一}。
  \item \textbf{记录与见证}：条件允许时，记下时间、地点、人物特征、对话，保存聊天记录/监控线索；有人陪同更好。
\end{enumerate}

\section{健康后果：可查、可治、可防}
性侵不只是"那一次"，它带来一连串\textbf{可处理}的医学问题，别自己扛：
\begin{itemize}
  \item \textbf{性传播感染（STI）暴露}：HIV、梅毒、淋病、HPV 等风险上升。就医后可评估是否需要\textbf{HIV 暴露后预防（PEP，72 小时内开始最有效}）、抗生素预防及后续随访检测（通常 2--4 周、3 个月、6 个月复测）。
  \item \textbf{意外妊娠}：无保护性交后，紧急避孕（口服紧急避孕药 72 小时内，或铜宫内节育器 5 天内）成功率较高；后续可评估常规避孕方案。
  \item \textbf{生殖器 / 泌尿道损伤}：撕裂、血尿、疼痛，需检查处理；女性创伤细节见\textbf{本系列 D 册“女性性健康”}相关章节。
  \item \textbf{慢性盆腔痛 / 性交痛（dyspareunia）}：若持续，常与创伤后肌肉紧张、心理因素相关，可治疗（见 B 册"常见困惑深挖"、D 册相关专题）。
\end{itemize}

\section{长期心理影响：可以恢复，不必"想开点"}
\begin{itemize}
  \item 常见反应：闪回、失眠、回避、情绪麻木、信任感下降，重者出现\textbf{创伤后应激（PTSD}）。这些是\textbf{正常人对异常事件}的应激，不是"脆弱"。
  \item 求助通道：心理科 / 精神科、心理咨询（推荐\textbf{创伤聚焦认知行为 TF-CBT} 或\textbf{EMDR}），以及本地妇女维权 / 心理援助热线。
  \item 关系修复（如与伴侣）：以\textbf{自愿、可撤回、循序渐进}为前提，不强迫"复原到从前"；必要时接受伴侣共同咨询。
  \item 一句给当事人的话：\textbf{你不需要为别人的暴力负责；你有权按自己的节奏处理，而不是被别人催着"放下"。}
\end{itemize}

\section{男性也是受害者（常被沉默的那一半）}
\begin{itemize}
  \item 统计上女性受害率更高，但\textbf{男性、非顺性别者同样会发生}；男性受害者常被"男人该扛住"的刻板印象耽误求助。
  \item 男性同样面临：STI / 妊娠（对伴侣）风险、生殖道损伤、PTSD，以及更强的羞耻与"不被相信"感。
  \item 处理路径与上文一致：安全→就医取证→STI/妊娠评估→心理支持；男性可到\textbf{男科或泌尿外科}说明情况，不必因"面子"硬扛。详见\textbf{本系列 C 册“男性性健康”}。
\end{itemize}

\section{预防与保护：把"同意"做成习惯}
\begin{itemize}
  \item \textbf{事前}：了解并尊重彼此的边界；饮酒后决策力下降，重要决定尽量在清醒时沟通好。
  \item \textbf{事中}：把"可以吗？还想继续吗？"当成自然的询问，而不是"质问"；\textbf{同意可以撤回}，任何时刻说"停"都有效。
  \item \textbf{环境}：公共/半公共场合的性行为注意隐私与照明；夜间出行尽量结伴、保持联络；职场/师生等权力不对等关系中，对"难以拒绝"的邀请保持警觉。
  \item \textbf{对他人}：被求助时，第一反应是"信 + 陪伴 + 不追问细节"，而不是"你怎么不早点说 / 你当时怎么没跑"。
\end{itemize}

\section{一页要点}
性侵犯 = 同意缺失，\textbf{不是受害者的错}。当下先保安全、再谈取证与就医（72 小时内就医 + PEP 窗口）；
健康后果（STI/妊娠/损伤）可查可治，别硬扛；心理影响可以恢复，按自己的节奏来。
男性同样是受害者，路径相同。预防的核心是\textbf{把"同意"做成日常习惯}，而不是事后追查。
相关女性专科细节见 D 册，男性专科细节见 C 册。


\chapter{附录：关键术语对照与权威来源}

\section{关键术语首现对照（便于跨册速查）}
以下缩写/术语在五册中反复出现，首次出现时请按下表理解，之后不再重复解释：

\begin{tabularx}{\textwidth}{l X X}
\toprule
\textbf{缩写/术语} & \textbf{中文全称} & \textbf{一句话说明} \\
\midrule
ED & 勃起功能障碍（erectile dysfunction） & 持续无法获得或维持足以性交的勃起，$\ge$3 个月；常与心血管、内分泌、心理相关。 \\
PEP & HIV 暴露后预防（post-exposure prophylaxis） & 疑似 HIV 暴露后 72 小时内开始、连服 28 天抗逆转录病毒药，可大幅降低感染率。 \\
PrEP & HIV 暴露前预防 & 高危人群每日服抗逆转录病毒药，事前降低感染风险。 \\
STI & 性传播感染（sexually transmitted infection） & 泛指 HIV/梅毒/淋病/HPV/衣原体等，感染（infection）比"病（disease）"更准确。 \\
HPV & 人乳头瘤病毒 & 分高危型（致癌：16/18 等）与低危型（疣：6/11）；疫苗可防主流高危型。 \\
SSRI & 选择性 5-羟色胺再摄取抑制剂 & 一线抗抑郁药；常见副作用为性欲下降、射精延迟、女性高潮困难。 \\
PDE5i & 5 型磷酸二酯酶抑制剂（万艾可/希爱力等） & 一线 ED 药物；需性刺激才起效；禁与硝酸酯类同服。 \\
TF-CBT & 创伤聚焦认知行为治疗 & PTSD/性侵后抑郁的一线心理治疗，聚焦"创伤记忆 + 认知 + 回避"。 \\
EMDR & 眼动脱敏与再加工 & 通过双侧刺激处理创伤记忆，对 PTSD 有较强证据。 \\
Kegel & 凯格尔运动 & 盆底肌收缩训练；男女皆适用，需正确"上提"方向（见 B 册“呼吸与盆底肌（Kegel）技巧”）。 \\
edging & 边缘控制 / start--stop & 接近高潮时退后、停留再升，延时长 + 增强度（见 B 册“高潮引导与边缘控制”）。 \\
G-spot & "G 点" & 阴道前壁 2--3cm 处的海绵体组织，个体差异大，并非人人有明显反应。 \\
SSC & 安全·理性·知情同意（Safe, Sane, Consensual） & BDSM 圈通行原则，所有玩法的前提。 \\
S\&M & 施虐（sadism）·受虐（masochism） & BDSM 中的两大核心动力，可单独或结合。 \\
D/s & Dominant / submissive & 支配者 / 服从者角色，与性别无关。 \\
aftercare & 事后关怀 & 高强度玩法（含 orgasm）后的情绪/身体照护，BDSM 与一般性都适用。 \\
\bottomrule
\end{tabularx}

\section{权威来源清单（正文引用与延伸阅读）}
本系列正文的事实与数据主要来自以下机构/共识/教材，按主题分类：

\begin{itemize}
  \item \textbf{全球性教育标准（ICSE）}：\href{https://www.wego.org}{wego.org}（WHO/UNESCO 联合维护），用于本系列"性是什么、性反应周期、同意与沟通"章节。
  \item \textbf{WHO / 联合国人口基金（UNFPA）}：生殖健康与性健康定义、全球 STI 指南、HPV 疫苗接种策略。
  \item \textbf{CDC（美国疾控中心）}：STI 诊疗指南（HIV/梅毒/淋病/衣原体/HPV）、PEP/PrEP 推荐、紧急避孕窗口。
  \item \textbf{EAU（欧洲泌尿外科学会）}：ED 诊疗指南、早泄（PE）共识、前列腺健康。
  \item \textbf{AUA / NICE（英国国家卫生与临床优化研究所）}：ED 与 PE 的一线治疗路径、盆底肌康复。
  \item \textbf{ACOG（美国妇产科医师学会）}：经期/孕期/更年期性生活、避孕方法可靠性。
  \item \textbf{DSM-5（精神障碍诊断与统计手册 第5版）}：性功能障碍分类（欲/兴/高潮/痛等）。
  \item \textbf{NHS（英国国家医疗服务体系）}：面向大众的性健康、PEP/PrEP、避孕选择。
  \item \textbf{中国}：《中国人类精子库技术规范》、国家卫健委 STI 诊疗指南、《中国女性健康指南》相关章节。
  \item \textbf{经典教材}：Human Sexual Response Cycle（Masters \& Johnson, 1966）、Koopman \& Berman（性反应模型更新）、Kaplan（三阶段模型）。
\end{itemize}

\textbf{使用提示}：本系列面向大众，非处方或诊断依据；具体诊疗请结合\textbf{主诊医生}建议。正文中出现"建议就医/复测"处，均为循证节点，不是营销话术。




\chapter{健康维护：把风险挡在前面}

\section{疫苗}

\subsection*{HPV 疫苗}
人乳头瘤病毒（HPV）感染与宫颈癌等多种疾病相关，疫苗是有效的预防手段。
在我国，二价、四价、九价 HPV 疫苗的获批接种年龄已扩展至 9--45 岁
（具体以最新说明书和当地接种门诊为准）。CDC 建议 11--12 岁常规接种，
可从 9 岁开始；26 岁前未接种者建议补种；27--45 岁者可与医生共同决策。
总体原则是：\textbf{越早接种，保护效果越好}。

\subsection*{乙肝疫苗}
乙肝可通过性接触传播。CDC 建议所有未接种、未感染且有风险者接种乙肝疫苗。

\section{定期筛查}

\begin{itemize}
  \item 有固定性伴侣前，双方做一次性传播感染筛查，是负责任的做法。
  \item 女性按年龄与医生建议定期进行宫颈癌筛查（如 HPV 检测、TCT）。
  \item 很多感染无症状，"没感觉"不等于"没感染"。
\end{itemize}

\section{日常卫生}

保持基本清洁即可，不必过度。重点是：行为前后适当清洁、
不共用可能接触体液的个人物品、出现不适及时就医而不是自行乱用药。

\chapter{心态准备：从容比技巧更重要}

\begin{itemize}
  \item \textbf{不必比较}：没有"标准剧本"，每个人的节奏都不同。
  \item \textbf{允许笨拙}：紧张、不熟练都很正常，沟通与体谅远比表现重要。
  \item \textbf{可以不做}："不想"本身就是充分的理由，任何人都不应被勉强。
  \item \textbf{尊重彼此}：把对方当作有感受、有边界的人，而不是目标。
\end{itemize}

健康的性观念，本质是把自己和对方都当作值得尊重的人。

\chapter{去哪里获取可靠信息}

信息源决定了观念的质量。请优先选择权威渠道，远离以流量为目的的
"经验分享"和以表演为目的的成人内容——后者不是教材，会建立错误的期待。

\begin{table}[htbp]
  \centering
  \caption{推荐的可信信息渠道}
  \begin{tabularx}{\textwidth}{@{}l X@{}}
    \toprule
    \textbf{渠道} & \textbf{说明} \\
    \midrule
    世界卫生组织（WHO）   & 全球权威的性与生殖健康指南与数据 \\
    美国疾控中心（CDC）   & 性病预防、安全套使用、疫苗接种的循证资料 \\
    国家卫健委 / 各地疾控 & 国内政策、疫苗接种点与健康教育信息 \\
    正规医院妇科 / 泌尿外科 / 皮肤性病科 & 个体化的诊断、检测与咨询 \\
    权威医学工具书         & 如默沙东诊疗手册（MSD Manual）等 \\
    \bottomrule
  \end{tabularx}
\end{table}

\chapter{深入篇：高频疑问解答}
\label{chap:faq}

前面各章给出的是框架，这一章回答初学者最常问的几个具体问题。
顺序从"工具怎么用"到"方案怎么选"，再到"话怎么讲""事怎么办"。

\section{安全套的正确使用}

多数失败并非产品质量问题，而是使用环节出错。对照下表自查：

\begin{table}[htbp]
  \centering
  \caption{安全套使用的常见错误与正确做法}
  \begin{tabularx}{\textwidth}{@{}X X@{}}
    \toprule
    \textbf{常见错误} & \textbf{正确做法} \\
    \midrule
    中途才戴、先接触后戴 & 生殖器接触之前就戴好 \\
    用牙齿/指甲撕开包装   & 沿锯齿边小心撕开，避免划破 \\
    没排空气               & 捏住顶端储精囊排出空气再展开 \\
    戴反了翻过来再用       & 戴反即换新，不翻面重复使用 \\
    用油基润滑（凡士林等） & 只用水基或硅基润滑剂 \\
    一个用到底、重复使用   & 每次每个行为都用新的 \\
    射精后不按住就退出     & 按住套口、趁勃起时退出 \\
    存放在高温/暴晒处      & 阴凉干燥处存放，看有效期 \\
    \bottomrule
  \end{tabularx}
\end{table}

\section{避孕方式怎么选}

没有"最好"的方法，只有"最适合当下情况"的方法。可以分三步判断。

\subsection*{第一步：看关系状态}

\begin{itemize}
  \item \textbf{新伴侣、非固定关系，或对方健康状况不明}：\textbf{必须用安全套}。
        这是唯一能降低性传播感染风险的工具，不能省。
  \item \textbf{稳定单一伴侣且双方已检测}：可在安全套之外，叠加更高效的方法，
        如短效口服避孕药、宫内节育器（IUD）、皮下埋植，形成"双保险"。
  \item \textbf{还没决定或只是偶尔}：安全套通常已足够，不必把问题复杂化。
\end{itemize}

\subsection*{第二步：看有效率需求}

若把"意外怀孕"的风险降到最低，长效可逆避孕（IUD、埋植）与绝育术的有效率
最高，短效药次之，安全套与行为类方法相对较低。但请记住：
\textbf{高效避孕 $\neq$ 防病}，只有安全套兼顾后者。

\subsection*{第三步：看健康与可及性}

激素类方法并非人人适合（如有某些血栓风险、偏病史等情况需医生评估）；
IUD、埋植需由医生放置；短效药需每日按时服用。请结合身体状况咨询医生，
不要自行长期服用来源不明的药物。

\section{如何与伴侣谈检测与边界}

很多人觉得难以启齿，但这恰恰是成熟的表现。可以事先想好一句开场白，例如：

\begin{quote}
"我想认真聊一下我们的健康和安全——不是不信任你，
而是我们都该对自己负责。我们要不要一起去做个检测，
以及确认一下彼此都能接受的边界？"
\end{quote}

要点：\textbf{对事不对人}，强调"保护我们俩"而非"怀疑你"；
把话题落在具体行动（检测、安全套、可以喊停）上，而不是抽象审问。
对方若尊重你，会认真对待；若回避、施压或嘲讽，这本身就是重要信号。

\section{无保护行为后的应对}

\begin{enumerate}[leftmargin=2em,itemsep=0.3em]
  \item \textbf{尽快求助}：紧急避孕越早越好，通常在 72 小时内，部分方案可延至 120 小时，
        请咨询医生或药师，不要当作常规避孕手段。
  \item \textbf{安排检测}：很多感染早期无症状，需在医生建议的时间窗后筛查；
        过早检测可能出现假阴性。
  \item \textbf{留意症状}：出现异常分泌物、疼痛、出血、发热等，及时就医。
  \item \textbf{下次改进}：把这次经历转化为经验，落实"每次、全程、正确"使用安全套。
\end{enumerate}

\section{必须就医的情形}

出现以下任一情况，请尽快到正规医院妇科、泌尿外科或皮肤性病科就诊，
不要自行买药或拖延：

\begin{itemize}
  \item 疑似无保护行为后需要紧急避孕或暴露后预防（如 HIV 暴露）。
  \item 出现疼痛、异常出血、分泌物异常、溃疡、皮疹、发热等症状。
  \item 检测呈阳性，需要规范治疗（多数性传播感染可治愈或可有效控制）。
  \item 对避孕方法有疑问，或需要 HPV、乙肝等疫苗接种建议。
\end{itemize}

\chapter{性传播感染详解}
\label{chap:sti}

性传播感染（STI）是初学者最容易恐慌、也最容易被谣言误导的话题。
这一章用循证医学的视角把它讲清楚：哪些能治、哪些不能、如何发现、出了问题怎么办。

\section{三个基本事实}

\begin{enumerate}[leftmargin=2em,itemsep=0.3em]
  \item \textbf{多数 STI 没有症状。} 世界卫生组织指出，大多数性传播感染没有症状，
        女性以及咽部、直肠等非尿道部位尤其如此。所以"没感觉"不等于"没感染"，
        唯一确知方式是检测。
  \item \textbf{分为"可治愈"和"可控"两大类。} 细菌性和寄生虫类通常能根治；
        病毒类多为终身感染，但可用药物抑制，并不影响正常寿命。
  \item \textbf{早发现、早治疗能防并发症。} 未经治疗可导致不孕、宫外孕、死胎、
        神经系统与心血管疾病，并增加感染 HIV 的风险。
\end{enumerate}

\section{可治愈类：细菌与寄生虫}

\begin{table}[htbp]
  \centering
  \caption{常见的可治愈性传播感染}
  \begin{tabularx}{\textwidth}{@{}l X X@{}}
    \toprule
    \textbf{感染} & \textbf{常见表现} & \textbf{要点} \\
    \midrule
    衣原体 & 常无症状；可有分泌物增多、排尿灼痛 & 可治愈；不治可致女性生殖系统永久损伤、不孕或宫外孕 \\
    淋病 & 男性可有尿痛、黄绿色分泌物；女性常无症状或很轻微 & 可治愈，但耐药菌株增多，须完成全程，不可自行买药或分给别人 \\
    梅毒 & 一期为单个无痛溃疡；二期可全身皮疹（含手掌足底） & 可治愈；潜伏期无症状但血清学阳性 \\
    滴虫 & 女性分泌物异常、外阴瘙痒、性交痛 & 可治愈；性伴需同时治疗 \\
    \bottomrule
  \end{tabularx}
\end{table}

\noindent 共性：一旦确诊，\textbf{你和伴侣必须同时接受治疗}，否则容易反复交叉感染。

\section{不可治愈但可控类：病毒}

\begin{table}[htbp]
  \centering
  \caption{常见的病毒类性传播感染}
  \begin{tabularx}{\textwidth}{@{}l X X@{}}
    \toprule
    \textbf{感染} & \textbf{特点} & \textbf{应对} \\
    \midrule
    HIV & 攻击免疫细胞；感染头几个月可能无症状但传染性强 & 无法治愈，但抗逆转录病毒治疗（ART）可长期抑制，病毒量测不出时不会传染性伴 \\
    生殖器疱疹（HSV） & 疼痛性水疱或溃疡，易复发；很多人无症状或极轻微 & 终身感染，抗病毒药可减少复发；不会经马桶、泳池、毛巾等日常物品传播 \\
    乙肝（HBV） & 可经性接触传播 & \textbf{有疫苗}；慢性感染需长期管理 \\
    HPV & 极常见，约九成感染两年内自行消退；少数持续感染可致宫颈等癌症 & \textbf{有疫苗}；宫颈癌筛查（HPV 检测、TCT）可早发现癌前病变 \\
    \bottomrule
  \end{tabularx}
\end{table}

\noindent 关键区别：病毒类"治不好"不等于"没救"。ART 抑制 HIV、抗病毒药控制疱疹、
疫苗预防乙肝和 HPV，都是成熟手段。

\section{中国法定管理的 8 种性病 + 常被忽略的病种}

\noindent 我国卫生部门将其列为法定管理的性病共 8 种：\textbf{艾滋病、淋病、梅毒、尖锐湿疣（生殖器疣/HPV）、生殖器疱疹、非淋菌性尿道炎（NOU）、软下疳、性病性淋巴肉芽肿（LGV）}。其中发病较多的是淋病、非淋菌性尿道炎、尖锐湿疣与梅毒。

\noindent 除上表已述者外，还有几个常被漏讲、但确实可经性接触传播的病：

\begin{tabularx}{\textwidth}{lX}
\toprule
病种 & 要点 \\
\midrule
非淋菌性尿道炎(NOU) & 多由衣原体/支原体引起，尿道刺痒/分泌物，症状轻于淋病但易漏诊 \\
软下疳 & 痛性溃疡 + 腹股沟淋巴结肿大；国内已较少见 \\
性病性淋巴肉芽肿(LGV) & 沙眼衣原体 L 血清型，先小溃疡后腹股沟肿痛 \\
阴虱 & 寄生在阴毛的虱，剧烈瘙痒，属寄生虫而非细菌/病毒 \\
生殖器念珠菌病 & 念珠菌（真菌）过度繁殖，见下节，\textbf{不算性病} \\
\bottomrule
\end{tabularx}

\section{非性病但与性相关的常见感染}

\noindent 这些常被误当"性病"，实际多为菌群失调或普通感染，可经性接触加重，但\textbf{不属于法定 STI}，处理也不同——混淆会耽误就医：

\begin{tabularx}{\textwidth}{l X X}
\toprule
感染 & 典型表现 & 说明 \\
\midrule
外阴阴道念珠菌病(鹅口疮) & 剧烈瘙痒 + 凝乳/豆渣样分泌物 & 念珠菌过度繁殖，与抗生素使用、潮湿有关；用抗真菌药，非抗生素 \\
细菌性阴道病(BV) & 灰白稀薄分泌物 + 鱼腥样异味 & 乳酸杆菌减少、厌氧菌增多；\textbf{不是念珠菌也不是滴虫}，勿自行乱用药 \\
滴虫性阴道炎 & 黄绿色泡沫分泌物 + 异味瘙痒 & \textbf{这个才算 STI}，须性伴同治 \\
膀胱炎 & 尿频、尿急、尿痛 & 多为肠道菌上行，非性传播；但性活动后易发，可"性交后及时排尿"预防 \\
\bottomrule
\end{tabularx}

\noindent\textbf{要点：}分泌物/瘙痒/异味\textbf{不等于性病}，但也别自行当念珠菌或 BV 处理——\textbf{拿不准就去确诊}（滴虫、衣原体、淋病都会表现为类似症状）。

\section{传播途径：不止性接触}

\begin{itemize}[leftmargin=2em,itemsep=0.3em]
  \item \textbf{性接触}（阴道/肛/口）：绝大多数 STI 的主要途径。
  \item \textbf{血液}：共用注射器、未严格消毒的采血/器械、输血（梅毒、HIV、乙肝）。
  \item \textbf{母婴}：孕期/产道/哺乳可传梅毒、HIV、乙肝、淋病、衣原体等——\textbf{孕前查清并治愈是关键}。
  \item \textbf{日常接触}：绝大多数 STI 不靠毛巾、马桶、泳池、坐椅传播；真正需要留意的是分泌物直接污染的物品（如共用毛巾、剃须刀），保持清洁即可，不必过度恐慌。
\end{itemize}

\noindent 参考量级（因个体而异，仅供理解"多数 STI 传染性不低"）：与梅毒一期患者单次性接触传染概率很高（早期资料估计数十至上百次中的大多数）；淋病随接触次数累积上升；与疱疹/疣患者的传染概率也偏高——\textbf{具体数据随病种与阶段差异很大，以医生评估为准}。

\section{症状与窗口期}

出现以下任一信号应尽快就医：异常分泌物、生殖器溃疡或肿物、排尿灼痛、下腹痛、
性交痛、发热皮疹。但更要记住——\textbf{没有任何信号也可能已经感染}。

以 HIV 为例，多数人在感染后 28 天内产生抗体；WHO 建议高危暴露后检测阴性者，
\textbf{28 天后再复查一次}，因为过早检测可能出现假阴性。其他感染的检测时机各有不同，
请遵医嘱安排。

\section{检测：怎么做、多久一次}

\begin{itemize}[leftmargin=2em]
  \item \textbf{方式}：多为尿液或拭子采样（尿道、宫颈、咽部、直肠），无创快捷。
        主动告诉医生你的性行为方式，便于选对采样部位。
  \item \textbf{频率}：CDC 建议有性行为的男男性行为者至少每年查 HIV、梅毒、乙肝、丙肝（依风险）；
        25 岁以下或有新性伴、多性伴、性伴已感染的女性每年查衣原体与淋病。
        普通人确定固定伴侣前双方各查一次，是负责任的做法。
  \item \textbf{阳性不是终点}：多数 STI 可治愈或可控，规范治疗、通知伴侣、按期复查即可。
\end{itemize}

\section{暴露后怎么办：PEP}

HIV 暴露后预防（PEP）用于可能接触 HIV 的紧急情况（如安全套破裂、共用针具、性侵）：
\textbf{越早越好，最好 24 小时内、最晚不超过 72 小时}到急诊或医疗机构启动，
\textbf{每天服药、连服 28 天}。PEP 不能替代常规防护；若反复暴露，
应咨询医生改用暴露前预防（PrEP）。

在国内，各地疾控中心、传染病医院、三甲医院感染科或急诊可提供咨询与 PEP，
也可拨打当地 12320 卫生热线询问。

\section{预防小结}

\begin{enumerate}[leftmargin=2em,itemsep=0.3em]
  \item \textbf{安全套}：每次、全程、正确使用。对体液传播的感染（淋病、衣原体、HIV 等）效果好；
        对皮肤接触传播的感染（疱疹、HPV、梅毒）只能部分降低风险。
  \item \textbf{疫苗}：HPV、乙肝（及甲肝）疫苗。
  \item \textbf{固定性伴与双方检测}：双方均未感染的单一性伴关系风险最低。
  \item \textbf{定期筛查}：无症状感染只能靠筛查发现。
  \item \textbf{避免酒精与药物影响下决策}：会降低判断力，导致不戴套等风险行为。
\end{enumerate}

\chapter{关系与现实：小白最该知道的几件事}
\label{chap:reality}

前面各章偏重"知识"，这一章偏重"观念"。对初学者来说，
破除迷思、学会沟通、保护自己，比记住任何技巧都重要。

\section{破除对"第一次"的迷思}

\begin{itemize}[leftmargin=2em]
  \item \textbf{"出血 = 第一次"是过时观念}：阴道瓣（旧称处女膜）形态各异，
        很多人第一次并不出血；剧烈运动、卫生棉条等都可能影响它。
        用它评判别人既不科学也不尊重。
  \item \textbf{疼痛不是必须的}：放松、沟通、充分的前戏与润滑能大幅减少不适。
        如果很痛就停下来——不是"不够好"，而是方法或时机需要调整。
  \item \textbf{别追求"完美表现"}：第一次更重要的是彼此信任和沟通，不是表演。
  \item \textbf{影视夸大了戏剧性}：现实往往平淡得多，甚至有点失望，这完全正常。
\end{itemize}

\section{不要把色情内容当作参照}

如果你对性的"首次印象"来自网络视频，请意识到\textbf{那是表演，不是教育}：

\begin{itemize}[leftmargin=2em]
  \item 时长、体型、反应、场景都经过设计和剪辑，和真实生活差距很大。
  \item 它几乎不涉及\textbf{安全、同意、情感}——而这恰恰是现实中最重要的三件事。
  \item 长期拿它当标准，容易对自己的身体和伴侣产生不必要的不满与焦虑。
\end{itemize}

\section{性欲与节奏：没有标准答案}

\begin{itemize}[leftmargin=2em]
  \item 每个人的性欲高低、唤起方式、节奏都不同，还会随年龄、压力、健康、情绪波动。
  \item 有人浪漫先行，有人身体先行——都正常。
  \item \textbf{差异不等于问题}。真正值得关注的信号是：欲望突然、显著地变化
        并带来困扰，那可以咨询医生或心理咨询师。
\end{itemize}

\section{学会说"不"，也学会接受"不"}

这是初学者最需要、却最少被教的一课：

\begin{itemize}[leftmargin=2em]
  \item \textbf{拒绝不需要理由}："我不想"就是完整理由，不必解释到对方满意。
  \item \textbf{具体说法}："我现在不想""我还没准备好""换个方式吧""我累了想休息"。
  \item \textbf{对方说"不"时}：尊重、不追问、不施压、不冷战。那不是拒绝你这个人。
  \item \textbf{危险信号}：用"如果你爱我你就会""大家都这样""你不配合我就生气"
        来施压——这是操控，不是亲密。
\end{itemize}

\section{身体意象与自我接纳}

\begin{itemize}[leftmargin=2em]
  \item 身体的形状、大小、肤色、体毛千差万别，\textbf{没有"标准款"}。
  \item 图片和视频里的光影、角度、剪辑都会改变呈现，真实的人不是那样。
  \item 好的亲密关系建立在\textbf{接纳}上，而不是把自己或伴侣和某个图像比较。
\end{itemize}

\section{几个容易被忽视的实际问题}

\begin{itemize}[leftmargin=2em]
  \item \textbf{酒后 / 药后}：判断力和自控力下降。提前想好界限，不单独赴不明场合的约。
  \item \textbf{线上交友}：不轻易透露住址、单位、经济状况；初次见面选公共场所并告知朋友；
        \textbf{不发送私密照片}——一旦发出难以收回，可能被用于敲诈。
  \item \textbf{隐私}：任何亲密行为都属私密。\textbf{未经同意拍摄或传播他人私密影像，
        是严重伤害，且违法。}
  \item \textbf{金钱与礼物往来}：保持清醒，别被物质裹挟做出不愿的决定。
\end{itemize}

\section{如果感到困惑或受伤}

\begin{table}[htbp]
  \centering
  \caption{遇到困惑或伤害时的求助渠道}
  \begin{tabularx}{\textwidth}{@{}l X@{}}
    \toprule
    \textbf{情况} & \textbf{求助渠道} \\
    \midrule
    性健康知识、检测咨询 & 正规医院、当地疾控、12320 卫生热线 \\
    情绪困扰、心理压力 & 心理咨询师、心理援助热线 \\
    性骚扰、性侵 & 立即离开现场、拨打 110、保留证据、尽快就医并联系可信任的人 \\
    \bottomrule
  \end{tabularx}
\end{table}

\section{常见误解：男女对彼此需求的迷思}

以下十项是伴侣之间（尤其长期关系中）最常见的性误解。
\textbf{核心原则：没有"标准男女"，只有"你我"。}
下表列出误解、事实依据、以及可操作的沟通建议。

\begin{tabularx}{\textwidth}{@{}l X X@{}}
\toprule
\textbf{误解} & \textbf{事实} & \textbf{怎么办} \\
\midrule
"女性最重视阴道性交" & 多数女性需\textbf{直接阴蒂刺激}才达高潮（Masters \& Johnson 及后续大样本调查）；阴道插入只是方式之一 & 共同探索；不必把"插入"当唯一目标 \\
"女性喜欢快而粗" & 多数人偏好\textbf{慢而渐进}的唤起；节奏应视双方感受调整 & 用语言/身体反馈告知偏好；"轻一点/慢一点"不算矫情 \\
"女性应该'自然而然'到高潮" & 阴道性交中女性高潮率显著低于男性；\textbf{需要技巧与沟通}，不是自动发生 & 告知对方你喜欢的刺激方式；可引导、示范、用器具辅助 \\
"快 = 好，性爱应速战速决" & 女性通常需要更长的唤起时间（前戏）；\textbf{时长因人而异}，不追求"效率" & 把前戏视为"正文"而非"序章"；关注双方满意度 \\
"亲吻/拥抱 = 性要求" & 很多亲密触碰表达的是\textbf{爱与亲近}，不一定指向性交 & 拿不准就问："你想做爱，还是只想抱一会儿？" \\
"我的身体 = 你的身体" & 敏感区域因人而异，\textbf{同一部位对他人不一定有同样快感} & 观察对方反应（表情、呼吸、主动靠近）；"喜欢请说、不喜欢也请说" \\
"女性以性换取爱" & 女性同样会为了\textbf{乐趣和亲密}主动要求性，不一定要"交换" & 主动表达需要不丢人；被拒绝也不代表不够爱 \\
"性是很严肃的事" & 适度的\textbf{幽默、轻松氛围}反而增强亲密感和唤起 & 笑一笑、说句俏皮话，让性不只是"任务" \\
"女性性欲范围窄" & 女性性幻想和兴趣的多样性与男性相当，甚至更广 & 尊重对方探索欲望的权利；不必替对方定义"正常范围" \\
"勃起消失 = 不爱了" & 勃起受\textbf{疲劳、压力、注意力、药物}等多因素影响；波动是正常生理 & 不把每次勃起失败都当"感情问题"；必要时就医排查 \\
\bottomrule
\end{tabularx}

\vspace{1em}
\noindent\textbf{关于"性欲 Q\&A"的补充}（依据既往家庭性教育文献，已用现代循证改写）：

\begin{itemize}[leftmargin=2em,itemsep=0.3em]
  \item \textbf{性欲不是纯本能}——它存在于大脑中，受情绪、关系质量、压力、睡眠、文化等多因素影响；没有"标准水平"，两人差不多才最重要。
  \item \textbf{长期关系中性欲不"消失"，但会"变形"}——从"新鲜好奇驱动"转为"亲密安全感驱动"；这是正常演变，不是感情破裂。
  \item \textbf{绝经后性欲可保持甚至增强}——激素变化影响润滑和频率，但不等于"不需要了"。
  \item \textbf{主动表达需求不是"不知羞"，被拒绝也不是"不爱你"}——健康关系里双方都有说"想"和"不想"的权利。
  \item \textbf{性欲低 $\ne$ 爱不够}——多数性欲障碍与压力、药物、疲劳、关系动态有关，而非"感情没了"。
\end{itemize}

\section{长期关系：从激情到日常}

"恋爱时火热如漆、结婚后平淡如水"——这不是你的错，而是\textbf{关系阶段的正常演变}。
关键是\textbf{从"本能驱动"转向"主动经营"}。

\textbf{为什么长期关系会"变淡"}（常见原因）：

\begin{itemize}[leftmargin=2em,itemsep=0.3em]
  \item \textbf{新鲜感消退}：多巴胺驱动的"热恋期"通常持续 6--24 个月，之后回落是生理规律。
  \item \textbf{安全感的副作用}：越亲密越怕被看穿、被拒绝，反而"不敢"主动表达欲望。
  \item \textbf{生活负担挤压}：工作、家务、育儿消耗精力，性成为"最后被牺牲"的那件事。
  \item \textbf{刻板化}：固定时间、固定体位、固定流程——大脑需要"变化"才能持续唤起。
  \item \textbf{沟通断裂}：以为"对方应该知道我要什么"——这是最常见的长期关系杀手。
\end{itemize}

\textbf{怎么破（可操作建议）}：

\begin{enumerate}[leftmargin=2em,itemsep=0.3em]
  \item \textbf{把"表达需求"当日常技能练}——"今天想试试 $\cdots$""你摸我 $\cdots$ 的时候我特别有感觉"；具体、简短、就事论事。
  \item \textbf{不要只在床上谈性}——在轻松场合（散步、睡前聊天）里聊比"做完之后复盘"压力小得多。
  \item \textbf{制造变化}：时间、地点、节奏、器具、角色——哪怕只是换一套床单，都能激活大脑的"新鲜感"回路。
  \item \textbf{重新引入"前戏"概念}：不只在床上，一条微信、一个眼神、一次不经意的触碰都是"唤起链"的起点。
  \item \textbf{别把"拒绝"当"否定"}：今天不想 $\ne$ 对你没感觉。把"好的，那你想怎样？"当默认回应。
  \item \textbf{允许自己"表演"一下}——不是虚伪，而是给关系加一点"仪式感"和"意外感"（既往案例提醒我们：先确认对方的接收频道再行动）。
  \item \textbf{性不平等是双向的}——女性主动要求性不是"放荡"，男性偶尔不想也不代表"不爱"；双方都有"发起"和"拒绝"的完整权利。
\end{enumerate}

\section{沟通风格与关系调适}

以下整理自 20 世纪 90 年代家庭性教育著作《夫妻间的心理差距》，保留其"沟通、尊重、体谅"的观念内核，
以现代关系心理学语言重新表述。

\textbf{伴侣间常见沟通差异}（倾向，非定论；个体差异远大于性别差异）：

\begin{tabularx}{\textwidth}{@{}l X X@{}}
\toprule
\textbf{维度} & \textbf{常见倾向 A} & \textbf{常见倾向 B} \\
\midrule
情感表达 & 更直接、就事论事、倾向给出明确判断 & 更委婉、含蓄、倾向用暗示和留白 \\
\midrule
冲突处理 & 先"想清楚"再回应（有时被感知为"冷战"） & 边说边理清思路（有时被感知为"唠叨"） \\
\midrule
信任与安全感 & 安全感来源偏重"空间独立" & 安全感来源偏重"情感确认和陪伴" \\
\midrule
亲密主动性 & 可能更习惯由"需要"驱动 & 可能更习惯由"情绪和连接"驱动 \\
\midrule
拒绝的表达 & 倾向直说"不" & 倾向用"改天吧""下次再说"来表达 \\
\bottomrule
\end{tabularx}

\vspace{0.5em}
\noindent\textbf{怎么用这张表}：不是给谁贴标签，而是当"对方没听懂我"或"我没听懂对方"时，
多考虑一步——"ta 的表达风格可能和我不同"，而不是"ta 不在乎我"。

\textbf{性生活中的平等与表达}（既往文献核心观点，现代语境改写）：

\begin{itemize}[leftmargin=2em,itemsep=0.3em]
  \item \textbf{性主动权是双向的}——任何一方（无论性别）都有主动表达性需求的权利；被拒绝不丢人，拒绝别人也不需要理由。
  \item \textbf{旧式"被动等待"模式}（"如果你真的爱我，你应该知道我要什么"）是长期关系最常见的沟通杀手；用"请求"代替"暗示"，用"我 $\cdots$"句式代替"你怎么 $\cdots$"。
  \item \textbf{伪装满意比被拒绝更伤}——假装高潮或假装不疼，短期省了尴尬，长期积累成"不被看见"的怨气。
  \item \textbf{性不平等常反映日常权力不平等}——谁决定时间、谁负责氛围、谁承担家务后的精力损耗——这些背景不解决，"床上的平等"很难实现。
\end{itemize}

\textbf{婚后心理负荷与调适}（既往文献"排除婚后心理负荷"要点）：

\begin{itemize}[leftmargin=2em,itemsep=0.3em]
  \item \textbf{理想 $\to$ 现实的落差}是正常的——婚前"展示美好一面" $\to$ 婚后"缺点一览无余"，这不是"对方变了"，而是"滤镜掉了"。
  \item \textbf{角色转换需要时间}——从"恋人"到"室友+搭档+经济共同体"，多重角色切换会消耗精力，性成为最先被压缩的维度之一。
  \item \textbf{不必苛求"永远如新婚"的性张力}——但需要主动"制造事件"打破日常惯性（一次短途旅行、一个不设手机之夜）。
  \item \textbf{性欲望的波动 $\ne$ 感情问题}——压力、疲劳、药物、季节、生理周期都会影响；把它当"天气"而非"态度"。
\end{itemize}

\textbf{关于"第三者"和关系破裂风险}（既往文献"婚外关系"部分核心）：

\begin{itemize}[leftmargin=2em,itemsep=0.3em]
  \item 婚外情的\textbf{多数根源在内部}：长期情感忽视、沟通断裂、性生活长期不和谐、角色固化。
  \item "在外面找到新鲜感"只是症状，不是病因——\textbf{不修复内部连接，换人也只是重复循环}。
  \item 发现外遇后的第一反应（愤怒、指责）可以理解，但\textbf{冷静分析 + 有效沟通 + 必要时寻求专业咨询}比"抓到小三打一顿"更能真正解决问题。
  \item 预防 > 修复：保持"配偶第一优先"的关系排序，定期创造二人独处时间，保持性生活的主动沟通。
\end{itemize}

\textbf{性格差异：互补而非对立}（既往文献"性格迥异的利弊"要点）：

\begin{itemize}[leftmargin=2em,itemsep=0.3em]
  \item \textbf{互补型组合}（外向×内向、冒险×稳定、行动型×思考型）在性生活和日常中都可以是优势——关键不是"消除差异"，而是\textbf{在差异中建立共识}。
  \item 磨合期的核心原则：\textbf{就事论事、不否定全部}——"你这次太快了" $\ne$ "你从来都只顾自己"。
  \item \textbf{己所不欲，勿施于人}——自己想要的，对方未必需要；先问再给，比"我觉得你需要 $\cdots$"有效得多。
  \item 正确的"弱点管理"：可改正的引导，不可改的接纳，\textbf{别把"性格"当"缺陷"来修}。
\end{itemize}

\vspace{1em}
\noindent\textbf{要点：}慢一点、问清楚、尊重彼此、保护好自己。
\section{婚姻危机期与性关系修复}

以下整理自家庭性教育著作《相互体谅才能让婚姻稳固》，聚焦对性关系影响最大的部分。

\textbf{婚后"危险期"}（既往文献引 1986 年上海数据：婚龄 1--5 年离婚占 41\%；现代数据同样显示前 5 年是关系不稳定期）：

\begin{tabularx}{\textwidth}{@{}l X@{}}
\toprule
\textbf{危机诱因} & \textbf{对性关系的具体冲击} \\
\midrule
角色切换（恋人 $\to$ 室友/搭档） & 精力被家务/工作/育儿挤压，性成为"最先被牺牲"的事 \\
光环效应消退（只看到优点 $\to$ 缺点全暴露） & 从"性吸引力滤镜"中醒来，身体吸引力感下降 \\
注意力再分配（夫妻 $\to$ 夫妻+孩子+父母+同事） & 共同话题减少，性交流也减少；双方"各忙各的" \\
外部压力（经济、工作压力、婆媳等） & 情绪低落时性欲自然下降，容易被误解为"不爱了" \\
第三者/外遇 & 性体验被"转移"；回归后重建信任需要时间和耐心 \\
\bottomrule
\end{tabularx}

\vspace{1em}
\noindent\textbf{感情五周期}（既往模型，现代改写）：
钟情幻想期 $\to$ 热情实现期 $\to$ 惆怅淡漠期 $\to$ 孤寂反感期 $\to$ 否定厌倦期。
\textbf{不是所有关系都会走到"冰点"——知道周期存在，才能在"淡漠期"主动介入，而不是等"否定期"才后悔。}

\textbf{危机期性生活怎么调（既往文献 11 条精选，已去时代化）}：

\begin{enumerate}[leftmargin=2em,itemsep=0.3em]
  \item \textbf{沟通 $\ne$ 批评}：在"非性时刻"聊感受；用"我 $\cdots$"而不是"你总是 $\cdots$"。
  \item \textbf{宽容 + 就事论事}：不要拿一次表现否定全部。
  \item \textbf{平等和尊重}：拒绝"夫唱妇随"或"妻管严"模式——性关系里的权力不平等会直接反映在质量上。
  \item \textbf{主动"感情投资"}：纪念日、生日、小惊喜——\textbf{仪式感和意外感是唤起的重要前驱}。
  \item \textbf{幽默化解紧张}：尴尬时用自嘲比用沉默好。
  \item \textbf{保持"性心理交流"}：不只是身体交合，还包括爱抚、亲吻、语言、目光——原文"多一点诗意和罗曼蒂克，少一点动物性本能"。
  \item \textbf{产后/育儿后的"再分配"}：别让亲子情完全替代夫妻情；留"二人时间"。
  \item \textbf{别拿"离婚"当日常武器}——反复威胁会摧毁安全感和性欲望。
  \item \textbf{背景压力（钱、家务、工作）不解决，"床上的问题"会反复出现}——性不是万能的，但性是关系的"晴雨表"。
\end{enumerate}

\textbf{性关系"修复路径"}（既往文献"情感危机可修复"改写）：

\begin{itemize}[leftmargin=2em,itemsep=0.3em]
  \item \textbf{第一步：恢复非性亲密}——拥抱、牵手、并肩散步、一起做饭；\textbf{先重建身体安全感，再谈性}。
  \item \textbf{第二步：重新建立"性对话"渠道}——不是"今晚想不想 $\cdots$"，而是"最近你感觉怎么样""你希望我们 $\cdots$"。
  \item \textbf{第三步：降低"表现焦虑"}——不追求高潮/时长/频率的"指标"，回到"彼此感受良好"。
  \item \textbf{第四步：引入变化}——时间、地点、节奏、方式；\textbf{大脑需要"新鲜感"才能重新唤起}。
  \item \textbf{第五步：必要时寻求专业帮助}——性治疗师、婚姻咨询师不是"最后手段"，而是"加速工具"。
\end{itemize}

\textbf{关于"平等"的底线}（既往文献"平等靠友爱互信"核心）：

\begin{itemize}[leftmargin=2em,itemsep=0.3em]
  \item 真正的性平等 $\ne$ "各做一半"，而是\textbf{双方都拥有完整的"发起权"和"拒绝权"}。
  \item "妻管严"和"大男子"是同一枚硬币的两面——都是\textbf{权力不对等}，都会侵蚀亲密感。
  \item 经济独立 $\ne$ 性独立——"我挣钱多所以我说了算"和"我做饭多所以我说了算"对性关系的影响一样大。
  \item \textbf{平等的前提是"看见"对方}：看见对方的需求、看见对方的疲惫、看见对方"不想"背后的真实原因。
\end{itemize}

\chapter{亲密过程中的身体反应：性反应周期与注意事项}

\begin{center}
\begin{tikzpicture}[font=\small]
  \draw[mygray,thin] (0,0) -- (9.5,0) node[midway,below,mygray]{时间 →};
  \draw[mygray,thin] (0,0) -- (0,3.6) node[midway,left,mygray]{唤起水平};
  \draw[myblue,thick] (0,0.4) to[out=45,in=170] (2.2,3.0) to[out=0,in=180] (4.6,2.7) to[out=0,in=110] (5.6,3.4) to[out=350,in=90] (9,0.5);
  \fill[myblue] (5.6,3.4) circle (0.06);
  \node[mygray] at (1.1,3.9){兴奋期};
  \node[mygray] at (3.4,3.9){平台期};
  \node[myblue] at (5.6,3.9){高潮};
  \node[mygray] at (7.4,3.9){消退期};
  \node[mygray,font=\footnotesize] at (1.1,-0.5){唤起·充血·润滑};
  \node[mygray,font=\footnotesize] at (3.4,-0.5){维持高位·张力升};
  \node[myblue,font=\footnotesize] at (5.6,-0.5){节律性收缩};
  \node[mygray,font=\footnotesize] at (7.6,-0.5){回落基线};
\end{tikzpicture}
\par\vspace{0.5em}\noindent\footnotesize 图：性反应周期（Masters \& Johnson）四阶段唤起水平曲线。
\end{center}
\label{chap:cycle}

本章讲"过程中身体到底会发生什么"——用经典的\textbf{性反应周期}四阶段模型，
从生理角度把过程说清楚。全程用医学科普的口吻，聚焦正常反应、个体差异和注意事项，
不涉及露骨的动作描写。先记住一句：这些反应大多是自动发生的，你不需要"表演"，
也不需要拿自己去对照任何标准。

\section{整体框架：四个阶段}

经典模型（Masters \& Johnson，1966 年《Human Sexual Response》）把一次亲密过程的身体反应
分为四个阶段：\textbf{兴奋期 $\to$ 平台期 $\to$ 高潮期 $\to$ 消退期}。
它不是机械流水线——各阶段长短因人而异，有时会跳过或反复，情绪、疲劳、压力、环境都会影响。
理解它，主要是为了知道"什么反应是正常的""什么时候该注意"。

\section{兴奋期：身体开始"准备"}

\textbf{身体会发生什么}（男女共同）：心跳加快、呼吸变快、血压上升、肌肉逐渐紧张、皮肤可能泛红；
男性阴茎勃起；女性阴道分泌润滑液、阴道扩张、外生殖器充血。

\begin{itemize}[leftmargin=2em]
  \item \textbf{唤起需要时间}，个体差异很大。女性往往需要更长时间，不要急着进入"正题"——这也是前戏的意义。
  \item \textbf{勃起受情绪和状态影响很正常}。紧张、疲劳、饮酒、压力大时勃起不理想很常见，不等于"有问题"。
  \item \textbf{干涩就补水基润滑剂}，这是正常生理差异（年龄、压力、药物都会影响），不是"不够投入"。
  \item 保持沟通："这样舒服吗？"比默默猜测可靠。
\end{itemize}

\section{平台期：反应持续加强}

\textbf{身体会发生什么}：心率、呼吸进一步加快，肌肉紧张加剧，部分人胸腹部出现红疹（"性红晕"）；
男性尿道口可能分泌少量液体；女性阴道下段收缩、相关组织进一步充血。

\begin{itemize}[leftmargin=2em]
  \item 这是"持续"阶段，节奏可根据双方感受调整——\textbf{快了就慢一点，累了就停一停}。
  \item 男性注意：\textbf{射精前的分泌物中可能已混有精子}，所以"体外排精"并不可靠，
        不能当作避孕方法（见第~\ref{chap:intimacy}~章）。
  \item 若一直"卡"在此阶段、许久没有进展，不必焦虑——放慢、换方式或顺其自然结束都可以。
        \textbf{亲密不等于必须走到某一步。}
\end{itemize}

\section{高潮期：短暂的峰值}

\textbf{身体会发生什么}：盆底及相关肌肉出现节律性收缩，心率、呼吸达到高峰，
主观上有强烈愉悦感；持续几秒到几十秒。男性射精通常与高潮同时发生，
女性是否高潮、强弱差异极大。

\begin{itemize}[leftmargin=2em]
  \item \textbf{女性高潮差异非常大}：有人容易、有人需要特定条件、也有人从未体验过——
        \textbf{都正常}，不必假装，更不要用"对方有没有高潮"评判这次亲密是否成功。
  \item 高潮不是必须的终点。\textbf{连接、愉悦、彼此舒服才是目的}，高峰只是可能的组成部分。
  \item 一次不理想很正常，别给自己下"不行"的定义。
\end{itemize}

\section{消退期：身体慢慢恢复}

\textbf{身体会发生什么}：心跳、呼吸、肌肉紧张在数分钟内逐渐恢复。
男性通常有\textbf{不应期}（refractory period）——此期间难以再次勃起，年轻时较短、随年龄增长变长；
女性一般没有明显不应期，但也可能因敏感而需要休息。

\begin{itemize}[leftmargin=2em]
  \item \textbf{不应期是正常的生理现象}，不是"能力下降"。别给自己或对方"马上再来"的压力。
  \item 消退期是\textbf{温存的好时机}：拥抱、说说话、聊聊感受。这段"事后时光"对关系的意义，
        往往比高峰本身更大。
  \item 若有无保护行为，记住 72 小时窗口（紧急避孕 / PEP），并按第~\ref{chap:sti}~章安排检测。
\end{itemize}

\section{贯穿全程的注意事项}

\begin{enumerate}[leftmargin=2em,itemsep=0.3em]
  \item \textbf{没有"标准"}：时长、反应、节奏人人不同，也随年龄、状态、情绪变化。
        别拿自己或伴侣和任何图像、视频比较。
  \item \textbf{安全套全程、正确使用}：这是全程不变的底线（用法见第~\ref{chap:intimacy}~章）。
  \item \textbf{保持清醒}：酒精会降低判断力，也让"同意"变得可疑。
  \item \textbf{持续沟通}：任何新动作前确认意愿，"还好吗？""要不要慢一点？"——沉默不是同意。
  \item \textbf{疼痛、出血、明显不适就停}，并考虑就医，不要忍。
  \item \textbf{情绪连接高于一切}：WHO 把良好的性健康定义为"积极、尊重、愉悦且安全的体验，
        免于胁迫、歧视和暴力"。技术是次要的，尊重与舒服才是核心。
\end{enumerate}

\section{几个常见困惑}

\begin{itemize}[leftmargin=2em]
  \item \textbf{"没有高潮是不是有问题？"}——多数情况下不是。关注愉悦和连接，而非某个生理终点；
        若长期困扰可咨询医生。
  \item \textbf{"中途软了/勃起不理想"}——偶发很常见；若持续出现并带来困扰，
        去正规医院（泌尿外科/男科）评估，别乱吃补药。
  \item \textbf{"总觉得太快了"}——先分清是不是"正常变异"或"主观感受"，别急着贴标签；
        持续困扰可就医（EAU 指南强调不能只看时间）。
  \item \textbf{"总觉得疼/干"}——先试水基润滑、放慢节奏；若持续疼痛，就医排查原因。
\end{itemize}

\chapter{避孕方法详解：怎么选才靠谱}

\begin{center}
\begin{tikzpicture}[font=\small]
  \node[anchor=east] at (0,0) {IUD / 埋植棒(LARC)};
  \fill[mygreen] (0.2,0) rectangle (0.247,0.3);
  \node[mygray] at (0.75,0.15) {\textless 1\%};
  \node[anchor=east] at (0,-0.7) {结扎};
  \fill[mygreen] (0.2,-0.7) rectangle (0.247,-0.4);
  \node[mygray] at (0.75,-0.55) {\textless 1\%};
  \node[anchor=east] at (0,-1.4) {口服药 / 贴片 / 环};
  \fill[myteal] (0.2,-1.4) rectangle (0.62,-1.1);
  \node[mygray] at (1.1,-1.25) {约 9\%};
  \node[anchor=east] at (0,-2.1) {安全套(男用)};
  \fill[myteal] (0.2,-2.1) rectangle (0.81,-1.8);
  \node[myteal] at (1.5,-1.95) {约 13\% · 唯一防 STI};
  \node[anchor=east] at (0,-2.8) {安全期 / 日历法};
  \fill[myamber] (0.2,-2.8) rectangle (1.05,-2.5);
  \node[mygray] at (1.55,-2.65) {约 18\%};
  \node[anchor=east] at (0,-3.5) {体外射精};
  \fill[myamber] (0.2,-3.5) rectangle (1.14,-3.2);
  \node[mygray] at (1.64,-3.35) {约 20\%};
  \node[anchor=east] at (0,-4.2) {安全套 + 杀精剂};
  \fill[myamber] (0.2,-4.2) rectangle (1.52,-3.9);
  \node[mygray] at (2.0,-4.05) {约 28\%};
  \node[anchor=east] at (0,-4.9) {不避孕};
  \fill[myred] (0.2,-4.9) rectangle (4.2,-4.6);
  \node[myred] at (4.7,-4.75) {约 85\%};
\end{tikzpicture}
\par\vspace{0.3em}\noindent\footnotesize 图：各避孕方法典型使用失败率（每年）。条形越长越易意外怀孕。
\end{center}

\label{chap:contraception}

前面多次提到安全套与紧急避孕。这一章把\textbf{全部避孕方法}放进同一张图里看，回答一个最实际的问题：\textbf{怎么选才既防怀孕、又防性病。}

\section{关键认知：两种"失败率"的差异}

\begin{itemize}[leftmargin=2em,itemsep=0.3em]
  \item \textbf{完美使用}：每次都严格按说明做（每天按时吃、每次都戴）。
  \item \textbf{典型使用}：真实生活里——会漏服、会忘、会中途出问题。
\end{itemize}

\noindent 同一个方法，两个数字可能差 3 倍以上。判断"靠不靠谱"，要看\textbf{典型使用失败率}（按真实使用，一年里约有多少人意外怀孕）。

\section{按"真实可靠度"分层（越靠上越可靠）}

\begin{tabularx}{\textwidth}{p{3.4cm}CCC}
\toprule
方法 & 典型使用约怀孕比例(每年) & 能否防性病 & 是否长效 \\
\midrule
宫内节育器/皮下埋植(LARC) & $<$1\% & 否 & 是(多年) \\
男/女结扎 & $<$1\% & 否 & 是(不可逆) \\
短效口服药/贴片/阴道环 & 约 7--12\% & 否 & 否(需每天) \\
安全套(男用) & 约 13\% & \textbf{是(唯一日常)} & 否(每次) \\
安全套(女用) & 约 21\% & 是 & 否(每次) \\
安全套+杀精剂 & 约 28\% & 部分 & 否(每次) \\
安全期/日历法 & 约 12--24\% & 否 & 否(每次) \\
体外射精 & 约 20\% & 部分(前液) & 否(每次) \\
不避孕 & 约 85\% & — & — \\
\bottomrule
\end{tabularx}

\vspace{0.5em}
\noindent\textbf{要点：} 要"几乎不担心怀孕"选 LARC（IUD/埋植）；要"顺便防性病"必用安全套。\textbf{两者不冲突，可以叠加}（安全套 + 口服药 = 既防孕又防 STI）。

\section{屏障法：安全套（唯一日常防 STI 的方法）}

\noindent\textbf{怎么用才防得住 STI}：\textbf{事前戴、全程戴、不翻转、戴在勃起上、射精后按住退出、用完即弃}。漏用/滑脱是中途出问题的主因；\textbf{不要与油基润滑同用}（会溶乳胶，见\ref{chap:practices}）。

\noindent\textbf{选型号与厚度}（国内常见，按勃起后直径选）：

\begin{tabularx}{\textwidth}{lX}
\toprule
型号(直径) & 说明 \\
\midrule
大号 35mm / 中号 33mm & 中号使用最普遍 \\
小号 31mm / 特小号 29mm & 过紧易破、过松易滑脱，都易失败 \\
\bottomrule
\end{tabularx}

\noindent 厚度：普通型 $0.04$--$0.07$mm、薄型介于其间、超薄型约 $0.03$mm；波纹/颗粒型增加摩擦，部分含延时或香氛成分。购买看有效期限（过期橡胶老化变脆），一次别囤太多；存放干燥阴凉、避光避油。

\subsection*{女用安全套}

女用安全套是一种由\textbf{使用者自主放置}在阴道内的屏障工具，不依赖男方配合。结构为一段聚氨酯管，\textbf{两端各有一个弹性环}：内环（较小、可弯曲）卡在阴道深处、覆盖宫颈，外环（较大）留在阴道口外、覆盖外阴。长度约 15--17\,cm、宽约 7--8\,cm，比男套更宽松。

\textbf{怎么用（关键步骤）：} 捏紧内环（呈 8 字形）推入阴道深处，直至\textbf{内环越过耻骨联合}、确认套体未扭转；外环自然留在体外。性交时阴茎进入套内而非直接进阴道。\textbf{性交后、起身前先捏住外环旋转拧住}再缓缓抽出（防精液漏出），用后即弃。\textbf{可提前放置，最多约 8 小时}——这是它独有的便利（经期、孕期、产后也可用，因\textbf{不含激素}）。

\begin{tabularx}{\textwidth}{l X}
\toprule
\textbf{项目} & \textbf{说明} \\
\midrule
避孕效果 & 完美使用约 95\% 有效；典型使用约 75--82\% 有效（\textbf{失败率约 21\%}，明显高于男套，主因放置不当/移位） \\
材质 & 聚氨酯（部分新版为丁腈），导热性好、\textbf{对乳胶过敏者更友好} \\
防 STI & \textbf{能防}；覆盖外阴，对 HPV/疱疹等皮肤接触型感染\textbf{覆盖面比男套更大}，但对 HIV 等体液传播的阻断效果\textbf{略逊于男套}（MedlinePlus） \\
\bottomrule
\end{tabularx}

\textbf{优点：} 女性自主掌控、不依赖男方配合；可经期/孕期/产后使用、无激素；乳胶过敏可用；保护外阴区。
\textbf{缺点：} 价格较高、需练习放置；套体可能\textbf{减少阴蒂直接刺激/发出摩擦声}（补水基润滑可缓解）；典型使用失败率偏高。

\textbf{注意：}\textbf{男用套与女用套不可同时用}（两层套互相摩擦会移位/破裂）；乳胶材质版\textbf{不与油基润滑同用}（聚氨酯/丁腈版相对宽松，仍以水基润滑最稳妥）；每次性行为\textbf{换新的}，不要重复使用；用后丢入垃圾、\textbf{勿冲入马桶}（易堵塞）。

\subsection*{屏障法失败后的补救}

套破/滑脱/遗留阴道内时：女方坐起让精液流出、温水冲洗外阴（灌洗阴道效果有限），\textbf{尽快按"紧急避孕"一节处理}（72h 内药片、或 5 天内铜 IUD），并留意后续月经是否如期。

\section{外用杀精剂与含药屏障}

\noindent 靠杀精剂直接灭活精子，常与"屏障"配合，单独用失败率偏高：

\begin{tabularx}{\textwidth}{lX}
\toprule
形式 & 要点 \\
\midrule
药膜/药片/药栓 & 性交前推入阴道深处，10--15 分钟溶化后再性交；每次都要放、间隔 1 小时再放 1 张 \\
杀精胶冻 & 与隔膜/安全套合用更佳，兼作润滑 \\
阴道隔膜(子宫帽) & 乳胶帽+弹簧圈，盖住宫颈；需医生测型号(50--80 号，国内常用 65/70/75)；性交后 8--12h 再取，勿超 24h；须配合杀精剂 \\
阴道避孕海绵 & 浸满杀精剂，可盖宫颈、24h 内可多次；失败率较高、少数过敏 \\
\bottomrule
\end{tabularx}

\noindent 共性提醒：杀精剂少数人局部有烧灼/过敏；患阴道炎、宫颈糜烂、子宫脱垂、阴道过紧者慎用或禁用；\textbf{杀精剂不能替代安全套防 STI}。

\section{长效可逆：LARC（IUD / 激素宫内环 / 埋植棒 / 避孕注射）}

\noindent LARC = \textbf{L}ong-acting \textbf{R}eversible \textbf{C}ontraception，一次放置/注射、多年有效，放置后基本"免维护"，典型失败率 $<$1\%，是\textbf{非绝育里最可靠}的，且\textbf{可逆}（取出后生育力通常很快恢复）。需到正规医院操作：

\begin{tabularx}{\textwidth}{lX}
\toprule
方式 & 要点 \\
\midrule
铜 IUD & 含铜、T 形，置入后可用约 5 年；少数有经量增多/经期腹痛，个别排斥 \\
激素宫内环(IUS) & 缓释孕激素，约 3 年；\textbf{经量过多/痛经者尤其合适}，前 3 个月可能点滴出血 \\
皮下埋植棒 & 火柴棒大小、植入上臂皮下，约 5 年；不适合有心脑血管/肝乳癌等禁忌者 \\
避孕注射 & 肌肉注射孕激素、按周期补打；可能增重、不规则出血，须按时复诊 \\
\bottomrule
\end{tabularx}

\noindent 破除一个常见误解：\textbf{"上环影响生育"多不成立}——取出后不影响；真正要注意的是放置/取出时机与个体适应证（听医生评估）。

\section{口服与经皮（短效口服药 / 贴片 / 阴道环）}

\noindent 核心就一条：\textbf{每天都吃/每周期按时换、别漏}——漏服/晚服是失败主因。\textbf{它不防任何性病}。刚起始的前 7 天建议配合安全套。常见健康收益（并非人人有）：月经规律、经量减少、痛经减轻、痤疮改善；需医生评估适应证（偏头痛伴先兆、吸烟且年龄较大、血栓风险等人群多不宜）。

\noindent 关于口服药几个高频误区：

\begin{itemize}[leftmargin=2em,itemsep=0.3em]
  \item \textbf{发胖}——多数研究提示体重变化与服药无明确因果；若首月后仍持续增重，可换配方。
  \item \textbf{停药周可有可无}——"21 天吃 + 7 天停"的停经周是设计遗留，可遵医嘱连服以避免每月出血。
  \item \textbf{漏服补救}——发现漏服一般\textbf{12 小时内补服}仍算数；超过则按说明书补 + 那几天加用安全套。
  \item \textbf{加倍剂量不加倍保护}——漏服按补救流程走，不是多吃一片了事。
\end{itemize}

\noindent 贴片（经皮吸收，7 天一片）、阴道环（置入阴道，3 周一换）同理：\textbf{按时换、不防 STI}。女性口服药的完整适应证、禁忌、漏服细节见 D 册“避孕药品”。

\section{自然法（安全期 / 日历法 / 基础体温 / 宫颈粘液 / 哺乳闭经）}

\noindent 靠"避开排卵期"使精卵不相遇，不干扰生理、无副作用、可"避孕与助孕两用"，但\textbf{需要规律月经 + 训练 + 双方配合}，失败率高于 LARC/口服药：

\begin{tabularx}{\textwidth}{lX}
\toprule
技术 & 要点与局限 \\
\midrule
日历法(安全期) & 依过去周期推算易孕期；周期不规律者误差大，单用不推荐 \\
基础体温(BBT) & 每日晨起静卧测口腔温，排卵后升高 $0.3$--$0.5^\circ$C；须整个排卵前禁欲、判读有滞后 \\
宫颈粘液法 & 观察白带：蛋清样、透明、拉丝长 = 接近排卵；需 3 个月训练 \\
症状-体温法 & 综合日历+粘液+体温，自然法里相对最可靠 \\
哺乳闭经法 & 仅限：完全母乳 + 产后 6 个月内 + 未复潮，三条件缺一即失效 \\
\bottomrule
\end{tabularx}

\noindent\textbf{要点：}自然法适合\textbf{愿意认真记录、月经规律、且能接受更高不确定性}的人；否则用屏障/LARC 更省心。

\section{紧急避孕（事后补救，不是日常方法）}

\begin{itemize}[leftmargin=2em,itemsep=0.3em]
  \item 左炔诺孕酮片：\textbf{越早越好，72 小时内}效果最佳（5 天内仍部分有效）。
  \item 乌利司坦：120 小时内。
  \item 铜 IUD：\textbf{事后 5 天内放置，是最可靠的紧急避孕}。
  \item \textbf{不能替代常规避孕}，频繁用会紊乱经期、失败率高于常规法。
\end{itemize}

\section{初学者常见误区}

\begin{enumerate}[leftmargin=2em,itemsep=0.3em]
  \item \textbf{"第一次/月经前后是安全期"}——排卵会提前、精子可存活数天，安全期并不可靠。
  \item \textbf{体外射精很安全}——射精前已有含精子/STI 的前液，典型失败率约 20\%。
  \item \textbf{"吃了药就不用套了"}——口服药只防孕，不防 STI；性伴不明确时两者都要。
  \item \textbf{靠"感觉"判断排卵期}——身体感受不准，别赌。
  \item \textbf{只在射精前才戴套}——前液已含精子与 STI，须全程戴。
  \item \textbf{把安全套放钱包/口袋久放}——高温摩擦易老化破裂，注意效期。
\end{enumerate}

\section{如何选择（决策清单）}

\begin{itemize}[leftmargin=2em,itemsep=0.3em]
  \item \textbf{要不要防 STI？} → 要：\textbf{安全套是底线}。
  \item \textbf{想长期、几乎不想操心？} → LARC（IUD/埋植），需就医。
  \item \textbf{暂时、短期、可逆、能按时吃药？} → 短效口服药 + 安全套。
  \item \textbf{一次性、不确定长期关系？} → 安全套为主。
  \item \textbf{已发生无保护行为？} → 尽快紧急避孕（72h 内）+ 考虑检测。
\end{itemize}

\vspace{1em}
\noindent\textbf{要点：} 最可靠的组合 = \textbf{LARC/口服药管怀孕 + 安全套管性病}；体外射精和安全期都不可靠；紧急避孕是补救不是日常。

\chapter{生命周期：不同年龄段的性健康}

前面各章讲的是"通用原则"，本章补上唯一缺的维度——\textbf{身体会随年龄变化，性和亲密需求也要跟着调整}。分四个阶段：青春期、产后、更年期、老年。核心观念先立住：\textbf{性健康不是某一个年龄的事，而是贯穿一生的；每个阶段都有该关注的不同重点。}

\section{青春期：建立认知和保护的起点}

\textbf{身体变化}：第二性征出现、月经初潮/遗精、性欲萌动——全是正常生理，不必羞耻或恐慌。

\textbf{这个阶段最该建立的三个认知}：

\begin{enumerate}[leftmargin=2em,itemsep=0.3em]
  \item \textbf{避孕与防病是两件事}——WHO 明确：\textbf{安全套是唯一既能避孕、又能防 STI（含 HIV）的方法}；口服药、埋植、IUD 都只避孕不防病。
  \item \textbf{"第一次"没有标准剧本}——不出血、不痛都很常见（呼应前面章节）；别被成人片或同伴话术带节奏。
  \item \textbf{避孕不会导致不孕}——WHO 澄清：现代避孕方法（含 IUD、埋植）\textbf{不会造成不孕}；"上环影响生育"是常见误解。
\end{enumerate}

\begin{itemize}[leftmargin=2em,itemsep=0.3em]
  \item 越早建立"全程安全套 + 同意 + 检测"的习惯，越受益；青少年同样可安全使用避孕药和 LARC（需医生评估）。
  \item 遇到困惑找\textbf{正规渠道}（校医、疾控、正规医院妇科/皮肤科），别靠网络偏方。
  \item 遇到任何让你不舒服的触碰——\textbf{同意可以随时撤回}，这不是"半途而废"。
\end{itemize}

\section{产后：恢复到重新亲密}

\textbf{什么时候可以恢复同房}（NHS/WHO 共识）：

\begin{itemize}[leftmargin=2em,itemsep=0.3em]
  \item 一般建议\textbf{等到恶露停止、会阴/剖宫产伤口愈合}，通常是\textbf{产后约 6 周}；顺产有撕裂或剖宫产者按医生具体指导。
  \item 过早（恶露未净时）恢复会\textbf{增加感染风险}。
  \item 恢复初期\textbf{动作宜慢、宜浅、用润滑}（激素变化常致干涩，尤其哺乳期），疼就停。
\end{itemize}

\textbf{避孕——这是最容易被忽略的点}：

\begin{tabularx}{\linewidth}{X X}
\toprule
\textbf{要点} & \textbf{说明} \\
\midrule
\textbf{产后可能很快恢复排卵} & 月经没来 $\ne$ 不会怀孕；先排卵后月经，可能"没来月经就怀上" \\
\textbf{哺乳闭经避孕法（LAM）} & 需\textbf{同时满足}三条件：(1) 纯母乳喂养 (2) 月经尚未恢复 (3) 产后 6 个月内；任一条件改变即失效，失败率随之上升 \\
\textbf{药物避孕时机} & WHO：纯母乳喂养者避孕药一般\textbf{产后 6 个月}开始；部分母乳喂养者\textbf{产后 6 周}；\textbf{仅孕激素迷你 pill 产后可立即用、不影响乳汁} \\
\textbf{LARC（IUD/埋植）} & 高效、可逆，产后由医生评估放置时机 \\
\bottomrule
\end{tabularx}

\vspace{0.5em}
\noindent\textbf{要点：}\textbf{别把"哺乳期"当天然避孕}——要么严格满足 LAM 三条件，要么用可靠的避孕方法。

\section{更年期（围绝经期）：变化是正常的，可处理}

\textbf{身体变化}（通常 45--55 岁，因雌激素下降）：

\begin{itemize}[leftmargin=2em,itemsep=0.3em]
  \item \textbf{阴道干涩、灼热、同房痛}——称\textbf{泌尿生殖综合征（GSM）}，是雌激素下降的直接结果，\textbf{不是"你不想了"或"感情问题"}。
  \item 性欲可能波动；潮热、情绪、睡眠也会受影响。
  \item 这一变化\textbf{很常见、且可处理}，不必默默忍受。
\end{itemize}

\textbf{怎么处理（NHS 明确）}：

\begin{tabularx}{\linewidth}{X X}
\toprule
\textbf{方法} & \textbf{要点} \\
\midrule
\textbf{润滑剂 / 阴道保湿剂} & 第一步、最安全；同房用润滑，日常用保湿剂 \\
\textbf{阴道局部雌激素（local HRT）} & 低剂量、只作用于阴道，\textbf{不增加乳腺癌风险、通常无需配孕激素、可按需长期使用}；最多约 3 个月起效；需医生评估 \\
\textbf{全身 HRT} & 针对潮热/情绪/睡眠等，需医生权衡利弊；局部阴道症状优先用局部雌激素 \\
\textbf{睾酮（低性欲）} & 非更年期常规适应证，仅专科医生在特定情况（绝经后、HRT 无效的低性欲）考虑 \\
\bottomrule
\end{tabularx}

\vspace{0.5em}
\textbf{两个易错点}：

\begin{itemize}[leftmargin=2em,itemsep=0.3em]
  \item \textbf{局部雌激素 $\ne$ 全身 HRT}：前者风险低得多，别因为害怕 HRT 而连局部治疗也拒绝。
  \item \textbf{更年期 $\ne$ 不用避孕}：NHS 建议——\textbf{50 岁以下}停经后仍需避孕 2 年，\textbf{50 岁以上}需 1 年；\textbf{55 岁后}自然怀孕极罕见（仍需注意 STI 防护）。
\end{itemize}

\section{老年阶段：需求不消失，方式可调整}

\textbf{核心认知}：性需求不会因年龄"消失"——\textbf{亲密、触摸、陪伴的需要贯穿一生}。
"老年人不该有性生活"既不科学也不人道。
Masters \& Johnson 的研究指出：\textbf{如果从早年起就维持活跃的性表达，且不出现急性/慢性生理残废，多数人在 70 甚至 80 岁仍保有活跃的性能力}。
调查数据：约 70\% 的男性在 68 岁时仍有规律性活动；15\% 的 60--94 岁老年人在 60 岁后还经历过性活动\textbf{增长}阶段。

\subsection*{男性生理变化（30--40 岁起出现，50--60 岁加速）}

\begin{enumerate}[leftmargin=2em,itemsep=0.3em]
  \item 性唤起和勃起速度减慢（需要更长时间、更多刺激）；
  \item 维持坚实勃起需要更直接的刺激（非"不行"，是"慢一些"）；
  \item 爱抚时间过长后勃起消退再恢复更困难；
  \item \textbf{不应期（refractory period）延长}（两次之间需要更长时间）；
  \item 高潮时生殖器区肌肉收缩减少、强度下降；
  \item 兴奋不再完全集中于生殖器（更"弥散"）；
  \item \textbf{勃起失败率随年龄上升}——但大部分是\textbf{心理/血管/药物}因素，并非不可逆。
\end{enumerate}

\noindent\textbf{关键区分}：偶尔勃起不好 $\ne$ ED。ED 的定义是\textbf{持续 3 个月以上}多数时候无法获得或维持满意勃起。偶尔因疲劳、压力、饮酒而表现不佳，是完全正常的。

\subsection*{女性生理变化（绝经后为主）}

\begin{enumerate}[leftmargin=2em,itemsep=0.3em]
  \item 阴道润滑速度和量下降（雌激素减少的直接结果）；
  \item 阴道壁弹性减低、变薄（萎缩性阴道炎/VVM）；
  \item \textbf{生殖器的敏感度反而可能增强}（不是"没感觉了"）；
  \item 高潮时阴道/子宫收缩次数可能减少，\textbf{但多重高潮能力不随年龄消失}；
  \item 需要更多生殖器\textbf{以外}的身体刺激来唤起（皮肤、亲吻、拥抱、言语）；
  \item 有时需要\textbf{人工润滑剂}辅助——这是工具，不是"病"。
\end{enumerate}

\noindent\textbf{处理优先级}：润滑剂（日常保湿 + 同房润滑）$\to$ 局部低剂量雌激素（医生评估）$\to$ 全身 HRT（针对潮热等其他症状，非专为性）。

\subsection*{观念与态度比生理变化更重要}

\begin{itemize}[leftmargin=2em,itemsep=0.3em]
  \item \textbf{把性当"欲望发泄" $\to$ 很快就"没意思了"；把性当"亲密连接" $\to$ 可以持续一生。}
  \item "性欲必随年龄下降"的信念本身就会让人\textbf{提前放弃}——这是自证预言。
  \item 老年夫妻的性活动"重心"自然从"频率和高潮"转向"情感质量和陪伴感"——这是\textbf{升级}不是"退化"。
  \item "少年夫妻老来伴"——老年性生活的核心价值是\textbf{情感交流 + 身体亲密}，而非"征服"。
\end{itemize}

\subsection*{心血管安全与性活动}

\textbf{心脏病 $\ne$ 禁欲}。当前医学共识：

\begin{itemize}[leftmargin=2em,itemsep=0.3em]
  \item 性活动的体力消耗 $\approx$ 慢速爬两层楼梯或中等步行 10 分钟（约 3--5 METs）——\textbf{对多数恢复期心脏病人在安全范围内}。
  \item 高潮时心率/血压会短暂上升（收缩压 +40--100 mmHg，心率 +50--100 bpm），\textbf{但持续时间极短}（数秒）。
  \item \textbf{需要避免/暂停的情况}：不稳定型心绞痛、未控制的心衰、严重心律失常、急性心梗后 6 周以内（需医生评估）。
  \item \textbf{实用建议}：避免饱食后立即进行；避免在极度疲劳/愤怒/寒冷环境下；准备硝酸甘油在床边（\textbf{但注意 PDE5 抑制剂禁忌——见 C 册}）；时间不要太长；\textbf{侧卧位}比俯卧位/仰卧位更省力、更保护心血管。
  \item 出现胸闷/气短/心悸 $\to$ 立即停止 $\to$ 休息 $\to$ 持续不缓解则就医。
\end{itemize}

\noindent\textbf{高血压}：血压控制稳定时可正常活动；控制不稳时应暂缓。不要因"怕血压高"而完全放弃——适度活动对心血管有益。

\subsection*{老年性生活的 11 条实践注意事项}（既往文献要点，已更新）

\begin{enumerate}[leftmargin=2em,itemsep=0.3em]
  \item \textbf{室温调节}——体脂减少、畏寒加重；冬天保暖，夏天避免大汗虚脱。
  \item \textbf{皮肤保护}——皮肤变薄、弹性下降；避免过度摩擦和高温（热毛巾温度要低一些）。
  \item \textbf{动作幅度减小}——肌肉萎缩、关节退化；避免过猛过大的动作，选择省力的体位。
  \item \textbf{颈部保护}——避免过度弯曲/伸展颈部（脑血流量减少，过度体位可能影响）。
  \item \textbf{耐心是关键}——男方等待女方充分唤起再开始；女方不因男方勃起慢而急躁。\textbf{双方都"慢一点"是正常需要，不是缺陷。}
  \item \textbf{呼吸平稳}——选择侧卧位/坐位，避免"费体力"的体位（如要求大幅用力的）。
  \item \textbf{避免额外心血管负担}——不过分用力、时间不过长、不过频。
  \item \textbf{不饱食后立即进行}——影响消化和舒适度。
  \item \textbf{排空小便}——避免不适；平时练习 Kegel 有助于功能维持。
  \item \textbf{频率减少但时间可延长}——不应期延长是正常生理；把"次数"指标放下，关注"质量"。
  \item \textbf{心理 $>$ 技巧}——把注意力放在双方的情感交流、拥抱、亲吻上；\textbf{不必每次以"插入 + 高潮"为目标}。
\end{enumerate}

\subsection*{疾病与性：不只是"能不能"，还有"怎么替代"}

\begin{itemize}[leftmargin=2em,itemsep=0.3em]
  \item \textbf{糖尿病}：神经病变 + 血管病变是 ED 主因；控制血糖 + 医生评估 PDE5 抑制剂或阴茎注射疗法。
  \item \textbf{前列腺术后}：暂时/永久性 ED 和尿失禁常见，但多数可恢复或部分恢复；别在恢复期就"判死刑"。
  \item \textbf{癌症治疗后}：体像变化（切除、造口）会影响自信，但\textbf{性需求依然存在}；伴侣的接纳和"重新发现身体"比任何技巧都重要。
  \item \textbf{残疾/运动障碍}：物理治疗 + 辅助器具 + 调整体位——性生活的范围远大于"标准姿势"。
  \item \textbf{丧偶后}：性需求不消失；再婚/新伴侣是社会可接受的，不必"为子女忍耐"。
  \item \textbf{慢性病与药物}：很多药（降压药、抗抑郁药等）会影响性功能——\textbf{可与医生沟通调整}，别自行停药。
  \item \textbf{安全套别停}：老年人 STI 感染率在部分人群上升，\textbf{新伴侣仍需防护}。
  \item \textbf{亲密形式可以更广}：拥抱、按摩、陪伴、相互爱抚——\textbf{"性"不等于"插入"}（呼应\ref{chap:outer}）。
\end{itemize}

\textbf{何时就医}：突发性功能变化、性交疼痛、出血、或因药物明显影响生活——都值得找医生，\textbf{不是"难为情"的事}。
医生应主动询问老年患者的性生活——但现实中 \textbf{90\% 以上从未被问过}（20 世纪 90 年代数据，至今未根本改善）。如果你没主动提，医生大概率不会问。

\section{本章速览}

\begin{tabularx}{\linewidth}{C X X}
\toprule
\textbf{阶段} & \textbf{身体关键变化} & \textbf{最该记住的一件事} \\
\midrule
\textbf{青春期} & 第二性征、初潮/遗精、性欲萌动 & 安全套是唯一既避孕又防病；避孕不致不孕 \\
\textbf{产后} & 恶露、伤口、激素波动、干涩 & 月经没来 $\ne$ 不排卵；LAM 三条件缺一不可 \\
\textbf{更年期} & 雌激素降、GSM 干涩、性欲波动 & 局部雌激素安全有效；停经后仍需避孕一段时间 \\
\textbf{老年} & 干涩加重、ED、药物影响 & 需求正常、可治疗；新伴侣仍要防护 \\
\bottomrule
\end{tabularx}

\vspace{1em}
\noindent\textbf{要点：} 生命周期视角 = \textbf{每个阶段身体不同、需求正常、都有对应办法}——青春期重"建立保护习惯"，产后重"恢复 + 别误判避孕"，更年期重"干涩可治、别怕局部雌激素"，老年重"需求不消失、可调整可治疗"。

\chapter{常见困惑深挖：很多"问题"其实是"正常"}
\label{chap:concerns}

很多让人尴尬的"情况"，其实处在\textbf{正常范围内}，或属于\textbf{可治疗、很常见}的问题——先别自我怀疑、也别羞耻。这一章把最常被问到的几类讲清楚。

\section{一、男性：勃起困难（ED）与早泄（PE）}

\subsection*{勃起功能障碍（ED）}
\begin{itemize}[leftmargin=2em,itemsep=0.3em]
  \item \textbf{原因常是两类叠加}：心理（焦虑、压力、关系紧张、表现焦虑）+ 生理（血管、神经、激素、某些药物）。
  \item \textbf{偶发很常见、属正常}；若\textbf{持续数周至数月且影响生活}，才值得就医。
  \item \textbf{可治疗}：治疗基础病（糖尿病/高血压/血脂）、在医生指导下使用 PDE5 抑制剂（如西地那非，\textbf{需处方、注意禁忌}）、心理与伴侣沟通。
  \item \textbf{红旗}：ED 有时是\textbf{心血管问题的早期信号}，医生常会顺带评估心血管风险。
\end{itemize}

\subsection*{早泄（PE）}
\begin{itemize}[leftmargin=2em,itemsep=0.3em]
  \item 分\textbf{原发}（一直如此）与\textbf{继发}（后来才出现，可能与其他因素有关）。
  \item \textbf{处理}：行为法（停--起始法、挤压法）、局部麻醉剂、心理（减压、与伴侣沟通）、在医生指导下使用处方药（部分 SSRI）。
  \item \textbf{关键}：别把它当"个人缺陷"——它是\textbf{常见、可管理}的，伴侣的理解与沟通本身就是治疗的一部分。
\end{itemize}

\section{二、女性：高潮困难与性疼痛}

\subsection*{高潮困难（anorgasmia）}
\begin{itemize}[leftmargin=2em,itemsep=0.3em]
  \item \textbf{多数女性需要直接阴蒂刺激}才到高潮（见\ref{chap:outer-pleasure}），"光插入不够"非常普遍，\textbf{不代表异常}。
  \item \textbf{常见原因}：焦虑/分心、关系因素、激素变化、药物（如部分抗抑郁药）、盆底紧张。
  \item \textbf{处理}：充分唤起与阴蒂刺激、与伴侣沟通偏好、必要时就医（评估激素/药物/盆底/心理）。
\end{itemize}

\subsection*{性疼痛（dyspareunia）/ 阴道痉挛（vaginismus）}
\begin{itemize}[leftmargin=2em,itemsep=0.3em]
  \item \textbf{疼痛原因多样}：感染、干涩、激素（围绝经期）、子宫内膜异位、盆底肌过度紧张、阴道痉挛。
  \item \textbf{处理}：润滑剂、局部激素（医生指导）、盆底治疗、心理、循序渐进的扩张训练、\textbf{找出根本原因}。
  \item \textbf{原则}：\textbf{疼不要忍}——反复疼痛应就医排查，而不是"习惯了就好"。
\end{itemize}

\section{三、双方：性欲差异（libido gap）}

\begin{itemize}[leftmargin=2em,itemsep=0.3em]
  \item \textbf{两人欲望不一样，很常见}，不等于"谁有问题"，更不等于谁不爱谁。
  \item 若\textbf{长期显著下降并造成困扰}，才考虑是否为\textbf{HSDD（女性性兴趣/欲望减退障碍）}，可就医评估（激素、药物、心理、关系）。
  \item \textbf{处理核心是沟通}（见\ref{chap:communication}）：谈感受、谈边界、找共同舒适的频率，\textbf{而非互相贬低或指责}。
\end{itemize}

\section{四、何时该去看医生（红旗信号）}

\begin{itemize}[leftmargin=2em,itemsep=0.3em]
  \item 反复/持续的\textbf{疼痛、出血、异常分泌物或体味}。
  \item \textbf{持续}的勃起困难、早泄并造成困扰。
  \item 性欲\textbf{长期大幅下降}并影响关系或情绪。
  \item 伴随\textbf{情绪困扰、创伤经历}——找正规心理/性医学科。
  \item 任何你"拿不准"的身体变化——\textbf{问医生比猜更靠谱}。
\end{itemize}

\section{本章小结}

\begin{tabularx}{\textwidth}{lCCX}
\toprule
情况 & 常见原因 & 可先做 & 何时就医 \\
\midrule
勃起困难 & 心理+血管/激素/药物 & 减压、沟通 & 持续数月、伴晨勃消失 \\
早泄 & 原发/继发、紧张 & 行为法、沟通 & 长期困扰 \\
高潮困难 & 需阴蒂刺激/心理/药物 & 前戏+阴蒂刺激 & 长期困扰、情绪受影响 \\
性疼痛 & 感染/干涩/激素/盆底 & 润滑、别忍 & 反复疼痛、出血 \\
性欲差异 & 关系/压力/激素/药物 & 沟通找原因 & 长期显著下降并困扰 \\
\bottomrule
\end{tabularx}

\vspace{1em}
\noindent\textbf{要点：} 绝大多数"性方面的困扰"都落在\textbf{正常范围}或\textbf{常见、可治疗}的区间——\textbf{别羞耻、别硬扛、先沟通，持续或疼痛就找正规医生}，这是最省时间也最安全的路。

\chapter{性治疗与性咨询：什么时候、找谁、怎么做}

\begin{center}
\myhei\large 找对"人"比"硬扛"或"乱试"有效得多——分清治疗师、咨询师、普通心理，走对路径
\end{center}

\section{三种专业帮助的界定}
\begin{tabularx}{\textwidth}{X X X}
\toprule
\textbf{角色} & \textbf{擅长} & \textbf{什么情况找} \\
\midrule
性治疗师（sex therapist） & 性功能的系统评估与循证干预（多含心理+行为+必要时生理协作） & 持续的勃起/早泄/高潮/疼痛/欲望问题，或关系中的性冲突 \\
性咨询师（sex counselor） & 关系、沟通、价值观、取向、创伤相关的性议题 & 性沟通困难、关系中的性张力、创伤后重建 \\
普通心理咨询师/心理科 & 抑郁、焦虑、PTSD 等影响性生活的心理状态 & 情绪障碍、创伤、压力主导的性功能变化 \\
\bottomrule
\end{tabularx}
\par\smallskip
\textbf{关键}：三者常\textbf{协作}而非"二选一"——严重的往往是\textbf{生理+心理+关系}交织，\textbf{先排查生理、再处理心理/关系}是稳妥顺序。

\section{值得就医的情形}
\begin{itemize}[leftmargin=2em,itemsep=0.3em]
  \item \textbf{持续 3 个月以上}的功能问题（勃起/早泄/高潮困难/疼痛/欲望明显变化），且带来困扰。
  \item \textbf{疼痛反复}：性交痛、干涩、阴道痉挛——\textbf{先看妇科/男科排查生理，再考虑盆底或心理}。
  \item \textbf{关系中的性冲突}：欲望不匹配（见 B 册）、沟通谈崩、亲密感下降。
  \item \textbf{创伤史}：性侵、早期负面性经历、宗教/家庭带来的性羞耻。
  \item \textbf{药物/疾病影响}：SSRI 类抗抑郁药、前列腺术后、糖尿病等导致的性变化——\textbf{先和主诊医生沟通}。
  \item \textbf{自我调整无效}：已经尝试过沟通、玩具、节奏调整等，仍无改善。
\end{itemize}
\textbf{不必找}：偶发的"状态不好"、单纯频率差异（多数是正常变异，见 B 册“欲望不匹配”）。

\section{常见循证疗法}
\begin{tabularx}{\textwidth}{X X}
\toprule
\textbf{疗法} & \textbf{适用 / 说明} \\
\midrule
认知行为（CBT） & 性焦虑、表现压力、错误信念（"必须每次都高潮"） \\
创伤聚焦认知行为（Tf-CBT） & 性侵/创伤后的性回避、闪回、信任下降 \\
EMDR & 创伤记忆的再加工，对 PTSD 相关性问题有较强证据 \\
盆底治疗 & 盆底肌紧张/无力导致的疼痛、控制问题（男女都适用） \\
伴侣咨询/性咨询 & 沟通、欲望差异、关系中的性张力 \\
药物（如 PDE5i、SSRI 调整） & 生理主导的功能问题，\textbf{需医生评估，别自行加减} \\
\bottomrule
\end{tabularx}
\par\smallskip
\textbf{提醒}：循证疗法有\textbf{适应证}，不是"万能"；\textbf{先明确问题类型}（生理/心理/关系），再选对应疗法。

\section{就医路径（科室与流程）}
\begin{enumerate}[leftmargin=2em,itemsep=0.3em]
  \item \textbf{先排查生理}：女性$\to$\textbf{妇科}；男性$\to$\textbf{泌尿外科/男科}。把症状、时间、诱因、已尝试过的办法说清楚。
  \item \textbf{心理/关系}：心理科/精神科、或正规医院的\textbf{性医学科/心身医学科}；也可找\textbf{有资质的性治疗师}（看学历、督导、受训体系）。
  \item \textbf{盆底康复}：\textbf{盆底康复科/康复医学科}，部分妇科/泌尿科也有。
  \item \textbf{怎么判断资质}：看\textbf{受训机构}（如 AASECT、ISSW 等认证体系）、\textbf{督导小时数}、\textbf{是否多学科协作}；\textbf{避"包治""秘方""一次见效"}。
\end{enumerate}

\section{初诊问诊内容}
\begin{itemize}[leftmargin=2em,itemsep=0.3em]
  \item \textbf{病史}：问题多久了、什么时候开始、是否伴随疼痛/干涩/出血/勃起变化。
  \item \textbf{生活}：压力、睡眠、饮酒、运动、用药（尤其 SSRI、降压药）。
  \item \textbf{关系}：伴侣情况、沟通模式、是否有创伤史（\textbf{可只说"有相关经历"，不必展开细节}）。
  \item \textbf{期望}：你想解决什么、能接受到什么程度——\textbf{"我想改善 XX"比"我不行了"更有用}。
\end{itemize}
\textbf{一句宽心的话}：多数性问题是\textbf{可改善}的；\textbf{说出来是第一步}，不是"丢人"。

\section{本章小结}
性治疗 = \textbf{循证、分型、协作}，不是"秘方""一次见效"。
\textbf{三种角色}：性治疗师（系统干预）/ 性咨询师（关系·创伤·沟通）/ 普通心理（情绪·创伤）——\textbf{常协作，非二选一}。
\textbf{何时找}：持续 3 个月+ 的功能问题、反复疼痛、关系性冲突、创伤史、药物/疾病影响、自我调整无效。
\textbf{顺序}：先排查生理（妇科/男科）$\to$再心理/关系$\to$盆底康复按需。
\textbf{循证疗法}：CBT、Tf-CBT、EMDR、盆底治疗、伴侣咨询、必要时药物——\textbf{有适应证，先分型再选}。
\textbf{资质}：看受训体系/督导/多学科协作，避"包治""一次见效"。


\chapter{多元取向与关系形态：尊重、安全与同意}

\begin{center}
\myhei\large 取向与关系形态都是"正常"的——核心始终是尊重、安全、可撤回的同意
\end{center}

\section{取向：尊重与语言}
\begin{itemize}[leftmargin=2em,itemsep=0.3em]
  \item \textbf{取向是个人体验}，不是"选择"也不是"阶段"——异性/双性/同性/无性/泛性等都成立。
  \item \textbf{用对方自称}：别"纠正"或"归类"（"你其实没那么喜欢"）——\textbf{用对方对自己的说法}。
  \item \textbf{非二元 / 无性 / 泛性}：不以"异性恋/同性恋"二分理解——\textbf{问"你如何描述自己"比"你是男/女"更尊重}。
  \item \textbf{安全不挑取向}：STI 预防、避孕、沟通在所有取向中都同样重要——\textbf{不要因为"看起来安全"就跳过屏障}。
\end{itemize}

\section{不同关系形态：同意与边界}
\begin{tabularx}{\textwidth}{X X X}
\toprule
\textbf{形态} & \textbf{核心} & \textbf{关键} \\
\midrule
一对一（monogamy） & 双方约定排他 & \textbf{"排他"要明说}，别默认 \\
开放式关系（open） & 允许多个性/情感对象 & \textbf{边界、知情、安全}三件事 \\
多角（polyamory） & 多段情感关系 & \textbf{沟通成本高}，定期对齐 \\
异地 & 物理分隔 & \textbf{信任+节奏}，见 B 册“远程与数字性爱” \\
\bottomrule
\end{tabularx}
\par\smallskip
\textbf{共同底线}：无论哪种形态，\textbf{知情 + 同意 + 可撤回}都成立；\textbf{边界是谈出来的}，不是"应该"的。

\section{安全（STI / 避孕 / 沟通）在多元关系中}
\begin{itemize}[leftmargin=2em,itemsep=0.3em]
  \item \textbf{屏障}：多对象情况下，\textbf{安全套/口交套}仍是第一道防线；\textbf{别因"信任"就跳过}。
  \item \textbf{检测节奏}：多对象情况下，\textbf{定期 STI 检测}（每 3--6 个月或换对象时）比"等出现症状"更稳。
  \item \textbf{避孕}：多对象情况下，\textbf{避孕责任要分工清楚}（谁用、用什么、备用方案）。
  \item \textbf{沟通}：\textbf{新对象前告知现有关系形态}（至少"我有伴侣"这一层），别"各藏各的"。
\end{itemize}

\section{语言与称呼}
\begin{itemize}[leftmargin=2em,itemsep=0.3em]
  \item \textbf{用对方自称的代词}：别"默认"。
  \item \textbf{别"纠正"取向}：对方说"我是 X"，就接受；\textbf{不必"分析"对方}。
  \item \textbf{亲密用语}：\textbf{问对方喜欢什么称呼}（"宝贝"有人喜欢有人不喜欢）。
  \item \textbf{身体部位/功能}：\textbf{别用"不像男人/女人"评价性表现}——\textbf{身体不是性别考题}。
\end{itemize}

\section{本章小结}
取向与关系形态都是\textbf{正常}——核心是\textbf{尊重、安全、可撤回的同意}，不因取向而变。
\textbf{语言}：用对方自称、别"纠正"、别"分析"、别用"性别考题"评价身体。
\textbf{关系形态}：一对一/开放/多角/异地——\textbf{边界是谈出来的}，知情 + 同意 + 可撤回。
\textbf{安全}：多对象情况下，屏障 + 定期检测 + 避孕分工清楚 + 新对象前告知现有关系。


\chapter{术语表与全书总览}
\label{chap:glossary}

把全书出现的关键术语整理成中英对照，方便日后检索；再把整本书串一遍，方便你按主题快速翻回对应章节。

\section{核心术语中英对照}

\begin{tabularx}{\textwidth}{l C X}
\toprule
术语 & 英文 & 要点 \\
\midrule
前戏 & foreplay & 唤起阶段的爱抚、亲吻、触摸 \\
边缘控制 & edging & 接近高潮时放慢/退回，再推上去 \\
同意 & consent & 持续、自愿、可撤回，不只是"没说" \\
偷摘安全套 & stealthing & 未经同意移除/不使用安全套 \\
长效可逆避孕 & LARC & 一次放置、多年有效的避孕（IUD/埋植） \\
宫内节育器 & IUD & 放入子宫的长效避孕装置 \\
性传播感染 & STI & 经性接触传播的感染总称 \\
疱疹 & HSV & 可口/生殖器互染，皮肤接触传播 \\
人乳头瘤病毒 & HPV & 与宫颈癌/口咽癌相关，疫苗可防 \\
牙科薄膜 & dental dam & 口交时的口腔屏障工具 \\
尿路感染 & UTI & 性行为后女性常见，事后排尿可降风险 \\
性疼痛 & dyspareunia & 性交疼痛，需找根因，别忍 \\
阴道痉挛 & vaginismus & 盆底肌不自主紧张致难以进入 \\
高潮困难 & anorgasmia & 长期难以达高潮，多与刺激/心理/药物相关 \\
勃起功能障碍 & ED & 持续勃起困难，常心理+生理，可治疗 \\
早泄 & PE & 射精过早泄且困扰，常见可管理 \\
性欲减退障碍 & HSDD & 性欲长期显著下降并困扰 \\
性欲 & libido & 对性的兴趣与欲望，因人而异 \\
暴露前预防 & PrEP & 高危者用药预防 HIV \\
暴露后预防 & PEP & 无保护暴露后 72h 内紧急用药 \\
紧急避孕 & emergency contraception & 事后 72h/120h 内的补救措施 \\
唤起伸展 & tenting & 唤起时阴道伸展增加空间 \\
性反应周期 & sexual response cycle & 兴奋/平台/高潮/消退四阶段 \\
\bottomrule
\end{tabularx}

\section{全书总览（按主题串一遍）}

\begin{enumerate}[leftmargin=2em,itemsep=0.3em]
  \item \textbf{为什么要读这本书}——定位与用法。
  \item \textbf{五个方面的准备}——认知、安全、沟通、健康、心态。
  \item \textbf{认知准备}——先懂基础，破除谣言。
  \item \textbf{安全防护}——最实际的一章（安全套/润滑/检测）。
  \item \textbf{沟通与同意}——不可逾越的底线。
  \item \textbf{健康维护}——把风险挡在前面。
  \item \textbf{心态准备}——从容比技巧更重要。
  \item \textbf{去哪里获取可靠信息}——权威渠道清单。
  \item \textbf{深入篇：高频疑问解答}。
  \item \textbf{性传播感染详解}——各 STI 的循证说明。
  \item \textbf{关系与现实}——小白最该知道的几件事。
  \item \textbf{见面了怎么办}——从邀约到在场的清单。
  \item \textbf{男性视角的补充}——同样重要的一半。
  \item \textbf{到了私密场所后}——以酒店为例的具体清单。
  \item \textbf{床上的事}——亲密过程的健康框架。
  \item \textbf{性反应周期与注意事项}——四阶段的生理过程。
  \item \textbf{亲密行为的类型与健康防护}——四类行为总览。
  \item \textbf{非插入类行为详解}——完整、合理、不遗漏。
  \item \textbf{非插入类行为的愉悦技巧}——从"不危险"到"很愉悦"。
  \item \textbf{亲密关系沟通}——从"会说话"到"更亲密"。
  \item \textbf{避孕方法详解}——怎么选才靠谱。
  \item \textbf{口交、阴道性交、肛交：逐类深入}——三类对照。
  \item \textbf{阴道性交详解}——单独深入。
  \item \textbf{口交详解}——咽部/黏膜感染与屏障。
  \item \textbf{肛交详解}——无自润滑+黏膜薄的防护铁律。
  \item \textbf{常见困惑深挖}——很多"问题"其实是"正常"。
  \item \textbf{术语表与全书总览}——本条。
  \item \textbf{出发前的清单}——逐项确认的实操清单。
\end{enumerate}

\section{去哪里查权威信息}

\begin{itemize}[leftmargin=2em,itemsep=0.3em]
  \item 国际：\textbf{WHO（who.int）}、\textbf{CDC（cdc.gov）} 的性健康/STI 页面。
  \item 国内：\textbf{国家卫生健康委员会}、\textbf{中国疾控中心} 的官方发布。
  \item 就医：正规医院的\textbf{性医学科、男科/泌尿外科、妇科、感染科、心理科}。
  \item 原则：\textbf{以权威发布为准，别把网络偏方当结论}；拿不准就问医生。
\end{itemize}

\vspace{1em}
\noindent\textbf{全书要点：} 性健康 = \textbf{知情 + 自愿 + 愉悦 + 安全}，四者缺一不可；这本书把"认知、防护、沟通、行为、方法、困惑"串成一条可执行的线，剩下的交给你的判断与医生的专业。

\chapter{附录：全书知识框架与补充索引}

前面各章是写给读者的内容，本附录换个视角：给出全书的\textbf{知识骨架与索引地图}，既方便按图索骥、也留作日后逐步补充的框架。每个节点标注了核心问题、关键概念、证据等级与覆盖状态。

\section{框架依据与图例}

\begin{tabularx}{\linewidth}{X X}
\toprule
\textbf{依据} & \textbf{说明} \\
\midrule
WHO 性健康定义（2002） & 与性有关的\textbf{身体、情感、心理、社会}四方面完好状态，不仅是无疾病；含"积极、尊重、愉悦、安全、免于胁迫歧视暴力" \\
生物-心理-社会模型 & 任何性问题都从\textbf{生物、心理、社会/关系}三层面看，避免单一归因 \\
性权利框架 & 性自由、平等、隐私、免于暴力、性健康权；对应同意、偷拍、影像传播等议题 \\
\bottomrule
\end{tabularx}

\vspace{0.5em}
\noindent\textbf{图例}：证据等级 \textbf{1} = 系统综述/权威指南（WHO/CDC/ACOG/EAU/NICE）；\textbf{2} = RCT/大型队列；\textbf{3} = 专家共识/权威教材；\textbf{4} = 科普共识/观察描述。覆盖状态 \textbf{已} = 已覆盖；\textbf{半} = 部分覆盖；\textbf{待} = 待补充。章号对应正文。

\section{A 基础认知}

\begin{tabularx}{\linewidth}{C X c c}
\toprule
\textbf{节点} & \textbf{核心问题 / 关键概念} & \textbf{证据} & \textbf{状态} \\
\midrule
A1 定义与权利 & 性健康不只是"没病"；性权利有哪些 & 1 & 已 \\
A2 女性解剖 & 内外生殖器；阴蒂/尿道/阴道位置；尿道短易 UTI & 3 & 已 \\
A3 男性解剖 & 阴茎/睾丸/附睾/前列腺；勃起射精生理 & 3 & 半 \\
A4 性反应周期 & 兴奋/平台/高潮/消退；唤起需时、阴蒂高潮 & 2 & 已 \\
A5 性取向与性别 & 取向与性别认同的区别；常见误解 & 3 & 待 \\
A6 性心理 & 性欲如何产生；心理唤起先于生理 & 3 & 半 \\
A7 生命周期总览 & 一生的性健康重点如何随年龄变化 & 3 & 已 \\
\bottomrule
\end{tabularx}

\section{B 健康与防护}

\begin{tabularx}{\linewidth}{C X c c}
\toprule
\textbf{节点} & \textbf{核心问题 / 关键概念} & \textbf{证据} & \textbf{状态} \\
\midrule
B1 STI 总论 & 病原谱、传播途径、无症状与筛查 & 1 & 已 \\
B2 重点 STI & HIV/HPV/疱疹/梅毒/淋病/衣原体/乙肝 & 1 & 已 \\
B3 避孕方法 & 完美 vs 典型失败率；按需求选 & 1 & 已 \\
B4 紧急避孕 & 72h/120h/铜 IUD 5 天 & 1 & 已 \\
B5 屏障工具 & 安全套用法、口用套/薄膜、润滑选择 & 1 & 已 \\
B6 暴露前后预防 & PrEP/PEP 适用与就医 & 1 & 半 \\
B7 疫苗 & HPV/乙肝/甲肝（肛交者） & 1 & 半 \\
B8 女性非 STI 问题 & BV/念珠菌/UTI 识别与就医 & 2 & 半 \\
B9 男性常见问题 & ED/早泄；处方药与就医指征；伴侣支持视角、自慰正常化 & 1 & 已 \\
B10 检测与筛查流程 & 查什么、窗口期、去哪查、阳性怎么办 & 1 & 待 \\
\bottomrule
\end{tabularx}

\section{C 行为与实践}

\begin{tabularx}{\linewidth}{C X c c}
\toprule
\textbf{节点} & \textbf{核心问题 / 关键概念} & \textbf{证据} & \textbf{状态} \\
\midrule
C1 行为分类 & 插入/非插入；体液 vs 皮肤接触风险 & 1 & 已 \\
C2 非插入类 & 阴蒂为核心；慢+润滑、边缘控制 & 2 & 已 \\
C3 阴道性交 & 唤起、润滑、阴蒂叠加、防 UTI & 2 & 已 \\
C4 口交 & 咽部感染、薄膜/口用套、HSV-1 & 1 & 已 \\
C5 肛交 & 无自润滑、黏膜薄、换套、水/硅基润滑 & 1 & 已 \\
C6 愉悦元技巧 & 慢、润滑、读反应、开口问、高潮非 KPI & 2 & 已 \\
C7 姿势与场景 & 按需求选；孕期/初次调整 & 4 & 已 \\
C8 常见困惑 & 高潮困难/性疼痛/性欲差异；红旗信号 & 2 & 已 \\
C9 BDSM 与 Kink & 去病理化(DSM-5)；SSC/RACK、安全词、aftercare；颈压/窒息风险 & 2 & 已 \\
\bottomrule
\end{tabularx}

\section{D 关系、同意与伦理}

\begin{tabularx}{\linewidth}{C X c c}
\toprule
\textbf{节点} & \textbf{核心问题 / 关键概念} & \textbf{证据} & \textbf{状态} \\
\midrule
D1 同意 & 有效同意要素、可撤回、stealthing & 1 & 已 \\
D2 沟通 & NVC 四步、时机、倾听 & 3 & 已 \\
D3 约会与现实 & 邀约到在场的清单、防偷拍 & 4 & 已 \\
D4 关系形态 & 长期/异地/开放关系的边界协商 & 3 & 待 \\
D5 法律伦理 & 隐私、偷拍、影像传播的法规边界 & 1 & 半 \\
D6 数字时代的同意 & 私密影像、AI 合成、网络骚扰 & 3 & 待 \\
\bottomrule
\end{tabularx}

\section{E 生命周期与特殊人群}

\begin{tabularx}{\linewidth}{C X c c}
\toprule
\textbf{节点} & \textbf{核心问题 / 关键概念} & \textbf{证据} & \textbf{状态} \\
\midrule
E1 青春期 & 身体变化、保护习惯、可靠信息 & 1 & 已 \\
E2 孕期 & 安全性、红线、姿势、干涩、引产 & 1 & 已 \\
E3 产后 & 恢复时机、避孕陷阱、LAM 三条件 & 1 & 已 \\
E4 更年期 & GSM 干涩、局部雌激素、避孕时限 & 1 & 已 \\
E5 老年 & 需求、ED、药物、防护 & 2 & 已 \\
E6 残障/慢性病 & 身体限制下的亲密与用药 & 3 & 待 \\
E7 LGBTQ+ 健康 & 针对性筛查、去污名 & 2 & 待 \\
\bottomrule
\end{tabularx}

\section{F 求助、资源与工具}

\begin{tabularx}{\linewidth}{C X c c}
\toprule
\textbf{节点} & \textbf{核心问题 / 关键概念} & \textbf{证据} & \textbf{状态} \\
\midrule
F1 就医路径 & 症状对应科室、如何开口 & 3 & 待 \\
F2 红旗信号 & 剧痛、大出血、发热、神经症状须就医 & 1 & 半 \\
F3 权威信息源 & WHO/CDC/ACOG/NHS/中国疾控/卫健委 & 1 & 已 \\
F4 术语表 & 中英术语对照 & --- & 已 \\
F5 检测与疫苗地图 & 本地化检测/疫苗渠道 & 3 & 待 \\
F6 更新维护 & 每次补充记录日期/节点/来源 & --- & 待 \\
\bottomrule
\end{tabularx}

\section{科学补充的方法}

\begin{enumerate}[leftmargin=2em,itemsep=0.3em]
  \item \textbf{一个节点一段}：先写结论，再写"为什么"（机制/证据），最后写"怎么做"。
  \item \textbf{标注证据等级}：新结论注明来源机构或文献；无权威来源者标"观点/共识"，不写成定论。
  \item \textbf{数字要核对}：失败率、时限、剂量等关键数值，写入前联网核对权威来源（同正文做法）。
  \item \textbf{中国语境}：涉及法律、就医、疫苗可及性时以中国法规与本地服务为准，注明地区差异。
  \item \textbf{避免重复}：补充前先看状态列，已覆盖的节点按需"深化"而非新建。
  \item \textbf{配图}：流程/分层/曲线类节点按全书 TikZ 风格补图，保持风格统一。
\end{enumerate}

\chapter{出发前的清单}

把前面的内容浓缩为一份可操作的清单。逐项确认，就能大幅降低风险。

\begin{enumerate}[leftmargin=2em,itemsep=0.4em]
  \item 我了解基本的生理与避孕常识，不靠谣言做决定。
  \item 我准备了合格、在有效期内、尺寸合适的安全套，并会用。
  \item 我备好水基润滑剂，知道油基会破坏乳胶。
  \item 我和伴侣已经谈过避孕、检测和各自的边界。
  \item 我确认双方都清醒、自愿，且可随时喊停。
  \item 我了解 HPV 疫苗，并考虑接种与定期筛查。
  \item 我知道无保护行为后的紧急避孕与就医途径。
  \item 我知道去哪里查权威信息，也知道遇到问题该找医生。
\end{enumerate}

\noindent\textbf{说明：} 以下为原《性学》一书内容，按现代 LaTeX 结构整理：原手写"篇/章"标题已转为标准章节命令，删除了与原正文重复的大纲段，其余文字尽量保留原貌。部分观点成书年代较早，请以正文各章的循证结论与医生建议为准。

\chapter{身体与性：男女构造、敏感点与荷尔蒙}
\label{chap:body}

这章把"身体"这件事一次讲清——男女生殖器官各是什么、哪里最敏感、荷尔蒙如何驱动性欲，以及几类流传很广的性观念误区。看懂这些，后面谈技巧、谈健康才有地基。

\section{女性生殖构造}

乳房

乳房的构造简单来看只是两团脂肪组织，整个区域唯一的肌肉只有乳头。当然，从性的角度来讲，乳房、乳头和乳晕（乳头周围的有色区域）对触摸会有强烈的性反应。这就是整个乳房区域在性刺激下会胀大的原因。

乳头是乳房最敏感的区域，因为它是由可勃起组织以及极为敏感的神经组织所构成。这也意味着在受到触摸或寒冷刺激时，整个乳房都会膨胀而且极为敏感。

外阴

从基本上来说，它是当你倾下身用镜子来窥探下体便可以看得见的性器官。它包括阴阜、大阴唇、小阴唇（也称为内、外“唇”），以及阴蒂。

阴阜：就是覆盖在耻骨上方、脂肪较厚的那部分皮肤，所有的阴毛都生在这个区域。可以在肚脐下方大约10厘米（4英寸）处找到这一区域。

在性爱中，一定要对这一区域善加利用。

不要把阴毛剃光，除非想让自己的下体上长满了令人刺氧的毛茬。

大阴唇：阴道的唇——就是遮盖阴道口的两层皮肤褶皱。

小阴唇：大阴唇内侧细小的皮肤褶皱。有些女性的小阴唇会下垂，或是从大阴唇中探出，而其他人的则是隐藏在外层的褶皱内。

如果注意到阴唇上有一些白色物质（这里要注意一下），那只是一种自然分泌物，是由分泌油脂的腺体产生的，跟男性包皮下的分泌物相似，很容易清洗掉。

阴蒂

阴蒂实际上只是一个较大器官的一个部分——我们所指的是延伸到你体内大约四英寸，并与阴道四周的可勃起组织相连的器官。女性的这种小型器官是与男性的阴茎相对应的。大致来看，它也有自己的柱状部分以及龟头，而且也包裹着一层类似于包皮的皮层，它被称为hood。此器官存在的作用就是给你带来性愉悦。

要找到它的位置，需要先找到小阴唇最内侧的两片唇闭合处的阴道口。要注意的是，这整个器官都是由跟阴茎相似的组织构成，所以，轻微的刺激就可以使它膨胀、发硬，而受到摩擦时会带来更多的美妙感受。阴蒂处于松弛状态时，只有一颗松子那么大，但受到刺激后，却能膨胀到一颗大花生米那样大。

不要对你的阴蒂大小大惊小怪。或大或小都没什么关系，不会影响你达到性高潮。

阴道

阴道实际上就是生殖器的体内部分，是一条有肌肉束缚而成的三至四英寸长的柔软管道。它具有很多种功能：首先，阴茎由此进入女性的身体；第二，它是婴儿出生的管道；第三，它还是经血的排出管道。

阴道在放松状态下，富有弹性的肌肉壁就会平贴在一起，但是在性事中，肌肉壁则会伸展开来，然后在“运动”结束后又会弹回原状。

其实，阴道各部分之间也有很大差别：大多数的奇妙感觉都是由阴道靠外的部分感受到的，所以阴茎进入此处时会带来愉悦感。而分布在中间部分以及靠近子宫颈部分的神经末梢就少得多了，所以这两处不会给你多少感觉。

阴道在每月的不同时期都会有一些正常的分泌物。

不要为了贪图“有趣”，就什么都往这里面塞。最好遵守这样一条经验法则：不能往嘴里放的东西，也不要往阴道里塞。

子宫颈

也被称为子宫的颈部，就位于体内阴道的上方，是一处大约2.5厘米长的由粗壮纤维组织构成的狭小部分。它的中央是一个小开口，精液可由此进入，而经血也是由此排出。这个开口在怀孕期间就会关闭，然后在生产时扩张开来以产出婴儿。总体来讲，子宫颈是由坚韧的肌肉组织以及其他组织构成的，但从细胞活力的角度来看，这里却是人体的薄弱环节。因为，在两性交合时，两种细胞在此相遇，而这有可能引发子宫颈细胞发生质变，还有可能导致宫颈癌。所以，定期的子宫颈涂片测试是必不可少的。

子宫颈很敏感，所以采用一些较深入的性体位时会有疼痛的感觉。即使触到了子宫颈，也不必中途停止性事，只要换个体位就行了。

子宫

在怀孕期间，这里就是婴儿的居所。子宫的长与宽大约都是7.5厘米（3英寸），而且具有强壮到不可思议的肌肉组织。在月经周期里，子宫内会生成一层子宫内膜。而在随后的一段时间内，如果没有受孕，子宫内膜就会自动脱落。

处女膜

部分覆盖阴道口的那层薄膜组织就叫做处女膜。因为它被认为是童贞的象征，在一些国家掀起轩然大波。实际上，处女膜经常会在剧烈活动中或使用卫生棉条时破裂——或许在你有性行为之前，它就已经破掉了。

如果处女膜是在你的“第一次”时破裂的，会有一点儿疼。处女膜在运动时破裂是很难察觉的。

\subsection*{女性常见敏感点}

乳房

一开始，你要让他把精力都集中在你乳房四周的区域。每边乳房下方和侧面的肌肤都是极为敏感的区域。让他由此入手，然后再慢慢继续后面的“工作”。

乳头

当他触及你的乳头区域时，不要让他用手捻弄，而要让他选择一些吸吮、亲吻或轻轻摩擦的方式，而且还要将乳晕包含在内。

G点

要找到这一点，要将一个手指探入阴道当中，然后向上移动5厘米（2英寸），触到阴道前壁，用手可以感觉到的一处小的肌里突起，这就是G点。如果他在你处于半兴奋状态时抚摸这里，你或许会幸运的体会到生平第一次G点高潮——一种超具诱惑力的高潮。

阴蒂

这是人体神经末梢分布最为集中的女性敏感点。如果能以正确的方式刺激这一点，你就立刻会有“哇~嘿”的反应。不过还是要给他一点建议：要在你下体完全湿润之后才能触摸此处，动作要轻柔，力道也要恰到好处。

会阴

这是位于阴道与肛门之间一小片稍稍隆起的区域，是极为敏感的性感地带，受到摩擦、舔吮或抚摸时，会有阵阵快感袭来，挑逗着你所有的敏感神经，并且有助于增强高潮。

\section{男性生殖构造}

阴茎

它是男性生理构造中对性刺激最为敏感的部位。阴茎是由可勃起的海绵组织、不随意肌以及血管组成的。别看它只是个挂在体外的简单柱状体，实际上，最有趣的还在于它的内部构造：

·阴茎的顶端（或称为龟头）极为敏感，龟头中央向外的开口就是尿道口，精液和尿液都是由此射出。

·阴茎在松弛状态时的一般长度为9厘米（3.5英寸）。

·阴茎无论大小都能勃起和达到高潮。

·大小无关紧要，实际上，虽然软下来的阴茎看起来不怎么相同，但是勃起时的阴茎一般都有13至18厘米（5至7英寸）长。

尿道

由膀胱通到阴茎顶端的管道。

·尿道是精液（含有精子的液体）和尿液共享的排出管道。

·当然精液和尿液是不会同时排出的。

尿道常识：性事之中，膀胱的一处肌肉会将出口紧闭，这样在射精时就不会有尿液排出。

包皮

覆盖阴茎并且包裹龟头的皮层。

·有些男性会出于宗教或健康的原因在童年时期就已经把包皮割除了。这种割除手术被称为包皮环切术。经过这种手术后，龟头就会永远显露在外了。

·包皮环切术不会影响到他或你的性快感。

·没有割除包皮的男性则要注意个人卫生，以预防非性病类传染病。只要经常洗澡并注意清洗包皮下面，即使没有割除包皮也会一样健康。

包皮常识：有人认为对女性来讲，跟割除了包皮的男性做爱会比较干净，其实是没有道理的。从芝加哥大学针对包皮环切术的一项研究来看，割过包皮的男性与没有割过的男性，在性健康方面没有任何明显差异。还有人认为包皮环切术可以预防宫颈癌，这种说法也已经被美国癌症协会否定。

睾丸

就是阴茎下方的两个小腺体。

·睾丸每天能产生多达几百万粒的精子。

·单个睾丸一般都有4厘米（1.5英寸）长，3厘米（1.25英寸）宽。

·通常左侧的睾丸要略微高于右侧的，以防止它们相互碰撞。

·每边的睾丸都是由一种索状结构悬于阴囊中的，这种索状结构被称为精索。在某些情况下，这两个球体会被拧扭或纠起，会引起极度的疼痛和反胃的感觉，这时可就需要紧急的医疗救护了。

·睾丸还会产生和分泌男性性荷尔蒙——睾丸素。

·千万别拿踢别人的睾丸开玩笑——一次直接的重击就能使它们破裂。

睾丸常识：睾丸癌是在20至40岁男性中最常见的癌症种类。80%的肿瘤最先都是以肿块的形式出现在睾丸中。但如果发现得足够及时，就会有95%的治愈可能。

阴囊

就是包裹睾丸的一层皮囊。

·阴囊就像一个袋子一样将睾丸悬在体外，因为产生精子需要的温度比人体正常体温（约 37℃）低约 2–3℃，体外位置正好提供这个"降温"。

阴囊常识：虽然阴囊看起来只有一个囊，但实际上分为两个部分，每边各装有一个睾丸。

前列腺

前列腺位于膀胱下方，负责产生精液中30%的液体。

·前列腺也被称为男性G点。

·它的大小跟核桃差不多。

前列腺常识：此腺体在直接刺激下会产生令人惊异的感觉。不幸的是要触到这一点唯一方法是要将手指伸入直肠中去。

精子

这种东西被称为精子或精虫，貌似蝌蚪而且也会游来游去，它们实际上是男性的卵细胞。精子的生理功能就是由男性体内向女性体内传递遗传信息。

·一般情况下，男性产生精子的速度可达每秒钟1000个。

·精子在睾丸中的成长需要两周时间。

·每位男性每月产生的精子多达三千万个，而每汤匙精液中就含有五千至一亿个精子（虽然其中40%的都是无用的）。

精子常识：女性能产生的卵子数量生来就是一定的，而男性与女性不同，他们在出生时还不能产生精子，要到青春期之后才会有这种功能。

精液
男性射精时，从阴茎中射出的体液与精子的混合物就叫做精液。

·精液是由10%的精子以及90%的体液（60%来自于精囊，30%来自前列腺）构成的。

·一次射精射出的精液只有一汤匙那么多。

精液常识：与一些常见的说法恰恰相反，精液中所含的并不都是热量，还有蛋白质和维生素，而且每次射出的精液中所含的热量很少（约 10–30 千卡，远不足以"补充营养"）。

\subsection*{男性常见敏感点}

阴茎

很多男性对他们的阴茎顶端总是紧张兮兮的，因为这里实在是太敏感了，而且有谁不怕疼呢？所以千万要小心，不要直接用牙齿去咬阴茎，也不要用舌头对它施加过多压力。口交功夫要做好，就要用你的嘴唇包住牙齿，吮吸的力道不能太大，更不能用力咬或抓。

男性乳头

尽管此处十分敏感，可不是所有的男人都喜欢你抚摸这里。如果你想要接近这一区域，一定要慢慢来，不要用手捻弄，而要用舌头。先观察一下他的反应，如果他露出笑容，那就说明你可以继续下去了。

系带

这是阴茎下侧一小块皮肤隆起，它可是极易点燃激情的敏感点，所以它是性爱中必不可少的一点。

男性G点

也被称作前列腺，要触摸到这一点，就要将手指探入他的肛门之中（这可不是所有男人都喜欢的举动，所以在这之前你先要征得他的同意）。

会阴

这是位于阴茎与肛门之间的一小块皮肤，有助于在性事中达到高潮。这里充满了快感神经，极为敏感，所以，男性们很喜欢这里受到轻柔的摩擦和爱抚。

\section{荷尔蒙如何驱动性欲}

性这东西，细想来真是有些滑稽，两个赤身裸体的人为了想象中的快感，大汗淋漓的交缠在一起。如果这件事不是所有人的必需，我们就能一道来谴责它是粗俗下流的行为。但事实上，总有一些个因素让我们就是无法对它嗤之以鼻。而且就是这些因素，总是不断怂恿我们想要跟人发生性行为。这些因素中最主要的就是我们的性荷尔蒙。这些偷偷摸摸的化学信号总是会像焰火一样在大脑中突然爆炸开来，唤起我们全部的性感官，并将其激发成性爱所需的一切感觉和情绪。如果你了解到它们是如何发挥作用的，你便可将它们为己所用。
\subsection*{男性：睾丸素}
睾丸素（testosterone）是主要的男性性激素，对维持性欲、勃起和肌肉量都重要。健康成年男性血液中的睾酮水平大约是女性的 10–20 倍。需要说明的是：\textbf{睾酮主要受生理状态（睡眠、运动、压力、体重）影响，而不是"看到异性就猛涨"这种即时反应}。以下六条建议中，符合现代证据的是"别太苛刻、选对时间、让他运动、表达渴望、保持规律性生活"；"让他自己解决（自慰）"属于偏好差异，因人而异。

对付男性荷尔蒙的六种方法：

a． 不要对他过于苛刻。
当他情绪低落时，他的睾丸素水平会一落千丈，你面前只是一个阴沉压抑、逆来顺受的可怜虫，那儿还会有什么性欲。

b． 小心挑选时间。
不要将性爱的激情都寄托在夜晚，他在那时会很疲劳，过多要求只会让你搬起石头砸了自己的脚。最好选在早上他“性致”高昂、体力充沛时，让他来尽一下“义务”。

c． 让他自己解决。
如果他在跟你做爱的同时还会自慰，你可不要感到厌恶。睾丸素不仅会激发他的性欲，也会增加他自慰的欲望。

d． 让他运动起来。
经常锻炼身体不仅可以促进人体全部性荷尔蒙的分泌，还能提升睾丸素的效率。

e．向他表明你对他的渴望。
当他"兴致"不高时，你表现出对他的渴望，本身就能提升他的情绪和性欲——积极的反馈是比任何激素名称都更有效的"启动器"。

f．有更多的性生活。
性爱可以促使大脑刺激睾丸素在睾丸以及脑部的分泌。因此，你们“做”得越多，他就越想“做”。
要注意在日常情况下，男性的睾丸素水平在早上较高（你可以分辨出来），到下午晚些时候下降到较低水平（这时他开始累了），而到了晚上就会非常低（这时他需要休息了）。其间，还可能有些上下波动。

\subsection*{女性：雌激素与孕激素（随月经周期波动）}
如果你觉得男性荷尔蒙是一篇“噩耗”，那再来看看女性荷尔蒙——雌激素以及孕激素。雌激素可是荷尔蒙中的一只热辣辣的小野猫，它性感、活泼，渴望他人的瞩目。而孕激素则是一张沉闷的湿毯子，渴望扑灭一切激情的火焰，毁掉你们的乐趣。这两种交替出现的荷尔蒙一直影响着你在月经周期每一天里的性情绪。这样，你该明白为何你的性欲总是起起伏伏了吧。
月经周期第一至七天
你会感到：
好像刚刚从一场恶梦中醒来，与前面七天不同，你会感到心情放松、愉快且平和。这都亏了脑垂体所分泌的一种荷尔蒙FSH(卵泡刺激素)。这种荷尔蒙还可以刺激卵巢产生女性性荷尔蒙——雌激素。
这时适合——
·开始社交。
这时不适合——
·“猎捕”男性。
·沉溺于高强度的性生活。
第七至十四天
你会感到：
性感。如果在第一至七天你已从风暴中心平静下来，那么在第七至十四天，你又会被卷入另一场风暴之中。而这几天也是你的月经周期中最重要的部分，这就是说是时候舒展一下筋骨了。因为雌激素的存在，让你的精力在此时最为旺盛。从性的角度来讲，你会欣喜若狂的去追逐你的心上人，心里只有一个坚定的念头——性。然而，虽然此时你的自信心很高，判断力却低的可怜，所以一定要小心。这时，你很容易会找上一个差劲的家伙，最惨的是，还要再过两周你才能意识到。
这时适合——
·性事。
·尽情的调情。
第十四至二十一天
你会感到：
厌烦、疲倦。这段时间可以说是在风暴中毁灭的日子。最初，大脑分泌的LH（黄体色素荷尔蒙）会引起排卵，而且会使你的性欲终日维持在较高水平。然而，排卵一经完成（在一至四天内）你的雌激素水平便会很快下降，而孕激素水平则开始上升。
孕激素可是荷尔蒙中的马丽·怀特豪斯 ，它会使你的性欲一落千丈。从身体上来看，你也许会觉得自己的体重有所增加，特别是胸部和上腹部（你的身体正在储存水分，以防在经期到来时发生脱水现象）。而且因为孕激素有一种镇定作用，你还可能在运动时觉得软弱无力、浑身懈怠。你的体力在这一时期也会呈下降趋势。
这时适合——
·呆在家里。
·跟他进行亲密的“爱的谈话”。
·早早上床休息。
这时不适合——
·连续熬夜
第二十一至二十八天
你会感到：
十分低落。你要做好充分的准备，因为在这段期间你最爱的两种荷尔蒙——雌性激素和孕激素都会消失。这意味着你可能会变得暴躁、易怒、伤感而且还会有些压抑。
在这段时间内，你可能会买和穿一些会令你终生后悔的东西（想象一下一双淡黄色的牛仔靴），而且这段时间还是你最容易感到沮丧的时候。这是因为此时你的精力都已经转移到身体的其他部分，注意力水平就会一直偏低。不过，好消息是：孕激素已经从你的周期中消失了，睾丸素（男性性荷尔蒙）会在你的体内初露端倪。它与处于下降趋势中的雌性激素相混合，会使你产生一阵突然的性冲动。这也就意味着你此时很容易获得高潮。
这时适合——
·待在床上。
·体会一次惊天动地的性高潮。
·吃些巧克力。
“我的男友说在我的月经周期里，每次月经和抑郁期过后，我都会有大约一天的时间心情特别好。”
——蒂娜，20岁
\section{几类最常见的性观念误区}
我们的荷尔蒙可以促使我们享受性生活或保持性感的感觉，但有时它们却受到我们某些性观念的压抑。所以一旦你感觉自己在性事中没有得到满足，或是感到你的性欲受到了压抑，那就要小心，说不定就是下面这些恶魔在作祟。
误区：性是肮脏的。
只有你想让它变得肮脏，它才会真的肮脏。这可不是开玩笑，你所接受的教养有时会阻止你享受性爱，而且会妨碍你在性事中发挥最大的潜力。或许，你所生长的家庭认为婚姻之外的性是“坏”的或“错”的，或者为了乐趣而做爱会违反你的宗教信仰。如果说这些约束都是合理的，那就值得考虑一下这样的事实：性，只要是发生在两个自愿的成年人之间，就如同呼吸一样自自然然、普普通通。
误区：女性要比男性寡欲。
好嘛，这就是某些男性想要我们相信的事情。实际上，女性的性欲跟男性不相上下。之所以有这种说法，是因为男性所接受的教育告诉他们，只有男性才有享受性爱的权利，而女性则没有。换一种角度，有些女性也会为自己的性欲感到愧疚，因此会竭力去隐藏、压抑甚至要忘记她们也有性的欲望。
误区：男人不可能对你忠诚。
有些男人是不忠，有些女人也一样是朝三暮四——这与“生理”以及所谓的“播种需要”毫无关联，而只是与心理伪装有关罢了。
误区：你不可能通过学习来精通性事。
你可以通过学习学好任何东西，尤其是性事。实际上，要想成为一个好情人，你要做的只是花些时间去学习动作。毕竟，最棒的情人们也不是因为有什么天赋才精通房术，他们只是很幸运的结交到一些可以给予他们“专业指导”的人士。
误区：你只能跟你喜欢的人做。
如果真能找到这么一位，也不失为一个好建议。但要注意的是：美好的爱情不等于美好的性，美好的性也不等于美好的爱情。尽管在性爱中，灵与肉的交汇是令人惊奇的，但“为性而性”可比你预想的更为普遍。而且，对很多人来说，性事也是个消磨时光的好方法。
误区：只有瘦子才能享受美妙的性爱。
性爱的美妙取决于胸部、大腿以及腹部（对男性则是阴茎以及肱二头肌）的大小。只有在你过于扭捏，或者性事半途中只想把肚子吸回去，还或是坚持要在一片漆黑中“办事”，体格的大小才会影响到性爱的品质。
误区：你不是一个性欲旺盛的人。
这是真的，还是你现在（或过去）的性生活让你有这样的感觉呢？事实是，如果你的性生活真的很糟糕，你是不会有什么“性致”的，而且实际上你可能根本就不想做。这样的话，就有充足的理由来检验和观察一下掩藏于表面下的真相了。
误区：每个人的性爱到最后都会变得平淡乏味。
性爱只有在你变得乏味的时候，才会变得乏味。这也就是说，如果你停止了试验新方法，停止了尝试新东西，停止了翻来覆去，而只会上床、下床，那你的性生活很快就会变得无聊透顶。好消息是，如果你肯多下点工夫，你未来的性生活便又会大放异彩了。
\section{性活动与健康的相关性}
性爱当然能为生活增添乐趣，也有研究显示它与一些健康指标存在正相关。以下是常被引用的几项研究——\textbf{请注意，多数是"相关性"而非"因果性"，不能把"多做爱"当成某种健康处方}：

\textbf{心血管}：英国布里斯托大学对 2500 名男性的观察发现，性生活频率较高者心血管疾病风险较低（约降低 45\% 左右）——这可能与情绪、运动和关系质量等混杂因素有关，而非"性行为本身治心脏病"。

\textbf{寿命}：英国威尔士大学的研究发现，规律性生活者的死亡率较低，同样受生活方式整体影响。

\textbf{激素与经期}：美国弗吉尼亚大学的研究提示，规律性刺激可能维持睾酮水平并轻度缓解经期不适，但证据有限。

\textbf{主观"显年轻"}：爱丁堡一项观察性研究发现，自评性生活满意度高者外表显得更年轻——更可能是整体健康、自信与皮肤状态的综合反映。
让你保持性感。
在年轻时就经常有性生活的男性与女性，在年纪大了之后还能在性事中保持生龙活虎的状态。想象一下，如果你几个月都没有使用过双腿，它们萎缩的会有多快？！你的性欲也是如此。如果你没有定期的“锻炼”，当你年纪大一些之后，你的性能力便会迅速的退化。
不用节食了！
规律性生活可能有助于维持性激素（如 DHEA、雌激素）水平，间接对代谢和情绪有正面作用；但"不用节食也能瘦"属于夸张说法，别把它当成减肥依据。
降低中风的危险。
定期的性生活还能使血压整体下降，这对你的健康是有长远利益的。血压下降，不仅可以让你的心跳更有效率，还能帮你避免因上了年纪而引起的中风。
有助于睡眠。
性高潮能使人们睡得很熟，因为在性高潮中，催产素荷尔蒙的水平会比平常突增五倍。而熟睡转而又能刺激大脑神经中枢在脑部释放褪黑激素，据说这种激素可以增强性欲、消除压抑情绪。最重要的是，褪黑激素是一种自由基抗氧化剂（自由基被证实有助于预防年龄性或与年龄有关的疾病，例如，关节炎，心血管疾病等等）。

\section{自慰：认识、方法与常见问答}

\textbf{正名}：自慰（masturbation）是最常见、也最健康的性表达之一。大多数男性、绝大多数女性都会自慰；它不会"伤身"、也不会"降低高潮几率"，反而帮你认识自己的身体、学会如何获得快感。很多人（尤其部分女性）对它有羞耻感，那多半来自观念而非事实。

\subsection*{关于自慰的常见疑问}
这个词听起来耳熟吗？这就是事实——大多数人，甚至是最懒的情人们都会自慰。这是因为，对很多女性来讲，当然男性也是一样，自娱自乐也是一种快乐。无论你有何偏好（让我们面对现实——是有一些人不喜欢这种行为），无论那只是在浮现你头脑中的幻象还是已经实践过的剧情，活跃的自慰性行为都是美妙性爱的关键。而干坐在那里等待一位专业情人来为你开启性的大门，这种做法根本就是性的死亡之吻（乍看是个好主意，实则害人不浅）。
那么我就来告诉你一个好消息——自慰、性幻想以及富于想象力的自慰行为，是不会让你为性而毁灭的。实事恰恰相反，它们不仅不会使你泥足深陷，也肯定不会降低你获得高潮的几率。
关于自慰的方法——有一些人做得很完整，全身赤裸，用整个身体来获取快感。而其他人则不会真的去触摸自己的身体，只是通过性幻想以及按压生殖器，高效率的进行自慰。有些人偏好性玩具（用法与选购见本系列“亲密实践与技巧”册“情趣用品与氛围”章），而还有些人则喜欢在办公室的文具橱上快速自行解决（如果你就是这么干的，可要小心电视台的摄像机）。
不管你喜欢哪种方式，这一章都是为你而写的。因为，如果你想拥有成功的性生活，你就需要有丰富的自慰经验。只要学会让你的思维与你的身体一起悸动、轻哼、呜鸣，那我保证你的情人们一定会经常跑来向你祈求更多的缠绵。
\subsection*{自慰的常见益处}
那么为什么自慰呢？顶级性学者认为，要想终生保持旺盛的性活力，定期自慰是最保险的方法。这是因为自慰对你的性生活有以下种种益处：
飘飘欲飞的自慰会带给你——
·一次完美的性发泄，无论你是不是单身都是如此。
·如果把性欲看作是一块肌肉，锻炼得越多，它就会越强，那你所谓的“爱的肌肉”锻炼得越多，也会保持越来越旺的劲头。
·除了性爱对健康的益处，自慰对健康还有一些独到的好处。它不仅可以帮你消除压力和焦虑，缓解痛经，还有助于睡眠激素的分泌。
·最重要的是，你可以把自慰当作一个投资银行，你在性爱的账户里有了投资，最终就能从性爱中获取利息。
·自慰还能帮你给你的情人进行正确的指导，你就可以告诉他，怎样做才能让你感到舒服。而且还能向他示范该如何改进爱抚的技巧。
如果你对此仍有兴趣，这里将向你介绍该如何尝试。
法则一：忽略你的法则。
自慰只是一种偏好，怎样进行都由你自己决定。然而，那些坚持要按部就班的朋友，还有那些觉得好笑的朋友，你们需要做的基本上就只是用手爱抚全身。
法则二：私密一些。
当然一开始就做到这一点对你还是很有帮助的。所以，记得要锁上房门，因为如果你一直担心会不会有人闯进来，就不会有性感的感觉了。还有，一定要选择合适的时间和地点，避开那些吵吵嚷嚷的室友，关掉手机和电话答录机（否则，一旦听到你妈的电话留言，你就什么情绪也没有了）。之后，你就要看看怎样做才会让你觉得舒服些——对某些人来说，就意味着关上灯，而对其他人来讲，也许就意味着赤身裸体了。
“触摸自己，我可受不了这种想法。我曾经试过一次，但是觉得实在是太恶心、太愚蠢了。我的性观念并不保守，只是自慰实在不适合我。”
——杰克，20岁
法则三：别光躺在那儿。
现在让你的双手开始四处游移，让它们轻轻擦过你身体的各个部分。提醒自己在第一章中所了解的那些敏感点，同时集中注意力在其它的一些性感地带上，例如嘴唇、面颊以及颈部。你或需早已了解哪些地方可以令你兴奋和颤抖，但要是能找到其它那些不明显的区域，不仅会改善你对自己身体的感觉，还能帮你开发做“热辣”情人的潜力。但首先要记住，这可不是什么科学实验，所以不要把自己的身体当作肉块似的对待。还要试一下怀着性感的念头爱抚自己，这会帮你体会到性感的感觉。
法则四：怀着性感的念头。
如果这是你第一次尝试自慰，或者是第一次怀着某种念头来自慰，这或许能帮你进入适当的情绪。对某些女性来说，她们需要的是一部性爱电影或一本色情小说。如果这些对你行不通，那就试着去想一些性感的东西——过去遇到过的人，前任好情人（或是现任的），甚至还可以回想一下以前看过的电影。
这样做的目的就是让你的情绪进入状态，不要觉得难为情，也不要想这看上去有多可笑，这些举动可都不利于获得高潮。问问自己：
·在按压自己的胸部时，力道是轻些好，还是重些好？
·你喜欢大腿内侧被爱抚或被轻捏的感觉吗？
·摩擦或爱抚乳头对你有效吗？
当你觉得自己已经做好准备，那就将重心转移到生殖器区域去吧。接下来，就介绍几种女性自慰的有效方法。
·用手、手掌以及手指直接摩擦阴蒂。
·用一样东西（例如枕头、毛巾等）来挤压阴蒂。
·对阴蒂使用振动按摩器。
·用淋浴头向阴蒂冲水。
·用手指插入阴道。
·爱抚并摩擦乳房、乳头和臀部。
·在身体的性感地带和生殖器区域尽情逗弄。
法则五：使其生效。
在试过以上方法之后，还是什么也没有发生吗？千万不要放弃。也许是你的期望过高或是太紧张了，还或是有什么东西干扰了你。实际上，有时自慰会带来高潮，而也有时只会让你不时的在边缘徘徊，却又根本发挥不了关键作用。如果有这种情况出现，试一下下面这些解决方法：
·看看你是否需要润滑剂的帮助吗？也许你不喜欢用这种东西，但如果你的下体像沙漠一样干涩，摩擦起来可是一点享受也没有。所以，有必要的话还是买些KY胶吧。
·接下来，想一下你用的体位。你是端坐着、躺下、站着还是在椅子上？初学者最好采取仰卧的姿势，将双腿分开。如果你觉得不适合你，或是觉得太过暴露，那就改用俯卧，把头放在枕头上。虽然这种方式“关照”不到你的乳房，但确实有助于你沉浸在此刻的激情中。
·那么现在，再来试着改变一下触摸自己的方式。你会讨厌插入吗？那就先别这样做。你会讨厌私处在摩擦时那种拖曳的感觉吗？那就试着在身体上爱抚或画圈。或者用手指抚过体侧，还或是想象你在做爱时多么渴望他的抚摸，然后就照你想象中的情形来做。
·不要担心性高潮。研究表明：大多数女性自慰时只需四分钟就能达到高潮，但是很多女性需要更长时间。
·改变你的节奏。既然你不想让你的情人来猛抓你的胸部或是猛撞你的私处，所以要想靠自己达到一次完美的享受，也要形成一定的步调和节奏。注意变换抚摸的速度，而且要伸直双腿或是将双脚的脚底相抵，来把身体绷紧。
法则六：到卧室之外的地方。
虽然你可能把床当作是实现你所有最佳设想的地方，但当你由足够的信心之后，就可以尝试一下床垫之外的地方了。当然，我可不是鼓动你在37路公共汽车上当众自慰，那可需要足够的勇气。你可以试一下在浴室里，或是淋浴时，或是坐在长镜对面的椅子上，或是穿着衣服站着做，还或者甚至在你的情人面前自慰。最后这种听起来好像很可怕，但它的确能延长做爱的时间，还帮助你在每次性生活中都能达到高潮。
\subsection*{男性的自慰}
这里只是简单介绍男性如何自慰。
这听起来好像有点性别歧视，不过男性对自慰可没有什么羞耻感，他们通常在十四岁之后对这件事就已驾轻就熟了，更别说在二十五岁的时候了。这就是说，一个男人也许在被你“钓”到之前或许就已经是个自慰的专家了。总体上来说，大多数男性进行自慰的速度都很快，而且很频繁，这与他们是已婚，单身还是有性伴侣都没什么关系。如果他经常自慰也就意味着他确切的知道该如何指导你该怎样做、不该怎样做。
（这也是你需要对男性的自慰行为稍作了解的一个原因）。
自慰能教会男性——
·如何延长性高潮。
·“无回复点”（the point of no returns）的感觉（这对你的性生活很重要）。
·怎样能使他们的高潮来的更快，还有怎样能使他们产生性冲动。
·这也是一种消除厌倦感、驱走性挫折以及减轻压力的好办法。
男性进行自慰的方法
男性在开始自慰时，通常会想象一些能让他们产生性冲动的东西。而且，男性对视觉刺激的态度要比我们女性开放、诚实得多，他们会利用很多种“工具”，像是黄色录像或色情杂志之类的，但更多人只用想的就成了。大致来讲，男性性高潮是由阴茎受刺激而引起的，因此，男性进行自慰可能采用以下这些方法：
·用手快速摩擦阴茎的柱状部分。
·专注于自己阴茎的顶端（龟头）——这是一个布满神经末梢的区域。他们或是只按压此处，或是前后摩擦阴茎的柱状部分与按压龟头交替进行。
·激烈的爱抚自己全身，并且摩擦阴囊和睾丸。
·刺激会阴、大腿以及肛门。
·向后拉扯包皮，并渐渐加快节奏和频率。
·一些男性还会使用润滑剂来帮助自己摩擦。
整个过程只需几分钟而已，这也是很多男性在做爱时容易出现早泄（PE）的原因——早泄的定义与治疗见 A 册“性医学基础”速查；控制射精的行为训练见 B 册“高潮引导与边缘控制”。
疑难解答
现在，你已经掌握了自慰的相关内容，尽管自慰对你的性生活意义重大，但如果你把它看得太重，结果只会适得其反，所以还要注意以下事宜：
·注意经常变换自慰的方法。如果你总是通过某种固定方式来达到高潮，你的身体会变得只对这一种方式有反应，这可会给你的性生活带来大麻烦。要记住，自慰是获得美妙性爱的关键，但不是获得美妙性爱的唯一方式。
·不要在情人面前假装高潮，再偷偷的靠自慰来为自己解决。虽然，经历一次美妙的自慰之旅是件好事，但也不值得在劣质的性爱后用自慰来补偿自己（还要小心他会翻身过来取笑你）。如果他没有满足你，就拿着他的手告诉他应该怎样做。
·不必过分担心他的自慰性高潮过于激烈。大多数男性在自慰时都会将自己的阴茎握得很紧，这样会带来强烈的性高潮，而这种感觉是无法靠女性的阴道来获得的。
·不必为他技巧不佳而大惊小怪。切记，自慰只是能帮你改善性生活的向导，不是标记详细的路线图。
·亲自实践之后还是讨厌自慰吗？那就不要再做了，自慰并不是每个人的“必修科目”。其实，有关于性爱的事情都是如此，只要觉得不舒服的，都不必去做。作为代替，你可以只是看着自己的私处，集中精力开始想象（这就是自慰行为中的身体幻想元素）。
·不要妄想他在跟你上床后就会停止自慰。如果他已经持续了几年，即使跟你在一起，也能毫无问题的继续下去。这并不是说你不能满足他，而且你也不必把这放在心上。除非他宁愿自慰也不愿跟你上床，那就真的有问题了。再说，你也不必为他停止自慰呀。正如前面所讲的，自慰，能够改善性生活、增进性欲、还能保持你性感的感觉，何乐而不为呢？
“我记得，我的一位前女友撞见我在自慰就大呼小叫的，她好像把这当作‘外遇’一样，一连几天都不理我，真是太奇怪了！”
——汤姆，25岁
\subsection*{一起自慰与交互自慰}
把你对自己身体的认识用到两人之间。要知道，除非你的男友会读心术，要不然他绝对读不到你的心思。要让你们互相了解如何点燃对方激情，最性感且最懒惰的方法就是两人一起自慰。在你们俩人赤身裸体、热情四溢时，尽管性交谈要容易百倍，撇开这个不说，交互自慰也不会像你想的那样难为情。实际上，你很可能已经做过了（记得你们在前戏中的动作吗？那就是我们所说的交互自慰）。这里有很多种供你尝试的方法，选一种最不会让你感到难为情的吧。
互相观赏
好处：这是让你了解你的伴侣喜欢哪种爱抚游戏的绝好机会，在激起性欲方面也是功效不凡。
总体而言，男性们似乎更加自信，他们会在你向他们展示你的性感地带之前，向你展示他们自己的身体。这听起来挺愚蠢，尤其是如果你已经拥有充满活力的性生活，但要鼓起勇气向他展示你全部的“自慰私密”，也的确是件难事。如果你实在怕羞，就让他先来吧。轮到你热情回报时，你就想象是在为他表演。如果想让自己感觉好些的话，你可以衣衫半蔽，不过还是要告诉他你在做什么，为什么这样做。这不仅能让他产生欲望，还能提醒他该如何改进前戏中的风格。以下内容可能会对你有所帮助：
·注意观察他怎样做的，以及他在达到高潮时的速度和节奏。
·向他表明你喜欢他怎样做，还有你不喜欢他技巧中的哪一部分。
·发出点声响。你也许已经体会到摩擦带来的确切感受，但他可能接受不到你内心的信号。那就用呻吟、叹息或大声的啜泣，让他感受到你的回应。
·仔细观察。不是不知道该如何用手替他解决吗，不用问，他会亲自向你示范的。
·不要被性高潮所困扰，要注意他是喜欢专注于阴茎还是偏向于转移到其它区域。
·不要只注重你的生殖器区域。向他展示你其它的敏感区域——乳房、颈部，你渴望被抚摸的所有区域。
·不要难为情，让他也加入你的行动中来，把他的手放在你的手上，告诉他你在想什么，让他来采取主动，并且让这种游戏成为你日常性生活中的一部分。
\textbf{交互自慰}：两人同时抚摸对方，目标是仅靠触摸完成整个过程。
交互自慰就是指两人同时抚摸对方。使用这种方法很快就能引燃两人的激情，推起共同的高潮。交互自慰的目标就是仅靠触摸来完成整个性事过程，所以要充分利用你们的手、唇、口以及声音去探索彼此的身躯。记住一定要把注意力投放在你身体上已被发掘出的全部的性感地带，还要再加上这些：
a． 体侧。
这是令人难以抗拒但又经常被忽略的两个区域，分别位于身体的两侧，确切位置在腋窝下5厘米（2英寸）处。当这两个区域被抚摸时，你会感到腹中有一种翻腾的感觉，兴奋的波浪沿着脊柱上下奔涌。这两个区域密布着神经末梢，有很少被触摸，所以极为敏感。要体会到这种销魂的感觉，就要他轻轻抚触这两个区域，时间要长一些、力道也要稍大些。
b． 耳垂。
耳垂可是非常非常敏感的区域，对这里的刺激也是极为有效。所以，无论磨蹭、轻咬或爱抚，都能向你的大脑传去“对了，就是这里”的信号。但仍需要警告一下，与人们的普遍想法相反，当你在性事中用到这一区域时， 切记口水浴对耳朵来讲，可不是什么性感的举动。
c． 手指与手。
十指连心，所以足见它们有多敏感了。试着轻抚或吮吸手指。要达到最好效果，你可以暗示他将你的一、两个手指慢慢放在口中做挑逗状的吮吸。这样，你就能了解到他“下面”的工夫是否了得，还有，与他四目相对，更会有一种用言辞难以表达的性感。
d． 双脚与小腿。
脚底以及脚趾间的空隙与性器官紧密相连，摩擦这些地方能比直接的摩擦性器官更具刺激性。这个柔软脆弱的区域渴望着关注与爱抚，让他将一只手指轻轻推进你的脚趾间，然后缓缓抽动，性感的激流就会向正确的位置直奔而去，而你立即就会兴奋不已。最后，一定要让他细细吻遍你的脚踝四周，然后从小腿向上一直爱抚到膝盖——这里也是一处敏感地带。
e． 大腿内侧
你也许早就了解了这一区域的功效，但是你是否知道除了用手对大腿的直接刺激，还有没有其他的方法能让你攀上“顶峰”呢？这次，又与神经末梢有很大关系，而且性参与度也有一定的影响。当然，还有什么能比炽热的气息与好奇的手指在大腿内侧上下游移的感觉更妙呢？让他用长时间按摩式的爱抚来挑逗你，如果你还受的住，还可以试一下搔痒式的轻抚。
f． 头发
有没有想过为什么你总喜欢用手梳理头发？那是因为轻拉头发可以起到按摩头皮、养精提神的作用。如果你想借此来引发性激情，需要注意两件事：一，你那位可要知道轻扯与猛拉之间的差别；下一条，确定你的头发不会打结，这样他的手才能轻松穿插而过。让他从两侧的发根开始，然后向内移动——一只手轻轻穿插而过，顺着向后滑到你的颈部，按摩着、亲吻着一路向下移去，轻拉头发也会有销魂的性感。
g． 背脊。
做过按摩吗？做过的话，就不需要我告诉你脊柱上也散布着几百个微小的性感点，而且大多数都在脊柱的根部。建议他从你的颈部开始，一路画着小圈圈下移到臀部以唤醒你性的激情。让他的手指沿着你的脊柱上下滑动，就如同弹钢琴一般，这样不仅能弹动你的欲望，甚至连你的背部神经也会为之奏起疯狂的旋律。还想要更刺激的，就要他一面按摩你的肩部和背部，一面将部分身体轻压在你身上。
h． 膝后及腰部。
这些都是容易被忽视的性感地带，真是遗憾，要知道这些区域可都是一些强劲的激情引爆点。换一种说法，只要轻轻几下亲吻、爱抚、吮吸就能使这些区域热血沸腾、温度上升、悸动不已，而你也会沐浴在愉悦的光辉中……还用我再讲下去吗？
i． 小臂内侧。
肢体预言专家指出小臂内侧也是一大性感地带，因为这里很少受到关注和爱抚。想一想，你有几回注意过身体的这一部分？卷起袖子，让他亲吻小臂下侧的柔软区域。你就会尝到火辣辣的性感滋味，还想要更来劲的，就让他用他的脸和唇，挑逗式的细细磨蹭。
j． 腹部
这个敏感点也是经常被错过的，通常是因为你总是忙着把它吸回去，不给他挑逗的机会。不过，让他的指和唇在这一带游移可是会带来别样滋味。如果你还受得住，就让他顺便“研究”一下你的肚脐（这会使你的身体轻轻颤栗）。不过，还是要提醒一句，如果他想要给你做个腹部按摩，一定要按顺时针方向（不然会影响到你的消化）。而且在他开始前要先确定一下，你应该不大可能在这种时候“肠气”外泻吧。
\section{性幻想}
自慰并不是让人快速达到高潮的唯一方法。美妙且传统的性幻想或许也会是你最佳的“启动工具”。如果你想在以后的性生活中轻触“开关”就能开启激情，那你的目标就应该是建造一座“图书馆”，在里面存满色情的经历、剧情以及纯粹的性故事。不管你是“独处”或与他人“相处”时，随时都能派上用场。
在最著名的一本女性性幻想书籍——南茜·芙莱德的《我的神秘园》中，来自世界各地的女性都谈到了对她们最为有效的各种性幻想。从中，你能找到从女奴幻想、假强奸，到普通遭遇的各式各样的性幻想。如果在电视上看到这些，绝对会让你笑掉大牙。而这就是幻想所能带来的纯粹的愉悦——这不是现实，也不需要忠实于现实，而你也更没有必要在现实生活中将其重演。
性幻想的意义就在于找到一个适合自己使用的“性欲启动工具”。一幕美妙的性剧情（如果它对你有效的话）可以让你仿佛置身于中心舞台之上，让你在刹那间产生性感的感觉。虽然性幻想与自慰经常是一对好搭档，但也要注意，制造一个火辣辣的性幻想对两人性爱也会有所帮助，特别是如果你因为某种问题而徘徊于高潮边缘，或者你的伴侣根本连想都懒得想。
\subsection*{性幻想从哪里来}
性学专家们认为，性幻想是随着我们不断成熟而形成于我们的头脑中的，它们来源于我们的所见、所读，也许还有亲身经历过的事情。除了这些，一些单凭其一点邪恶和禁忌就能使人感到愉悦与兴奋的事物，也可能成为幻想的原料。无论性幻想背后的原因是什么，如果你对它会有罪恶感，就不要沉溺于其中——毕竟那些白日梦都只是大脑中的幻象，只是用来改善性生活的工具。它们都不具真实性，既不能说明你是个变态，也肯定不能代表你真正喜欢做的事情。
女性在性幻想中想些什么？
·被一个陌生人侵犯。
·跟某个名人做爱。
·来个“三人行”。
·与其他女性做爱。
·被人观看。
·受到惩罚或鞭打。
某些人的性幻想始终只有一个固定的主题，而其他人则会不断变换主题。但是研究显示：男性与女性的性幻想中总会有一定相同的模式与主题，都会倾向于“跟某位明星”、“被迫”、“惟命是从”，还有“被主宰”或“主宰”之类。那么再来看一下两性共有的幻想主题吧。
尽管男性偏好更为视觉化的影像（想一下那些色情杂志），但在所有的性幻想中都有这样一个不争的、共同的事实——男性与女性都偏好于自己喜欢做的事，起码也要是喜欢在想像中“实践”的事。还在为猜不透他的而烦恼吗？下面就让我们一起来戳穿这些有关性幻想的谬论。
误区：我梦到了女人，我一定是同性恋了！
根本不是（当然除非你知道自己是）。同性幻想只是跟自己的想象力一起玩的游戏，除了说明你的好奇之心，再没有其它的含义。当然，如果在产生这种幻想之后，你觉得自己以前选错了方向，而且渴望有一次同性经历，这时，也只有在这种时候，同性幻想就有更多含义了。
误区：我做爱时幻想到其他男人，那我一定要离开现在的男友了。
问一下任何一对相处数年的情侣，他们两人都会承认自己会不时假装在跟某个电影明星做爱，而不是他们的男友或女友。这并不代表着他们真的想跟那个明星有什么，这只是一种
为夜间长期运作的性生活增添情趣的秘密方法。所以你不必为此感到愧疚。当然，除非你每次做爱都会有这种幻想，而且也确实想离开现在的男友，那就另当别论了。
误区：有非常性爱的幻想就意味着你是个变态。
男性在性幻想中想些什么？
·巨大的乳房。
·主宰他人。
·被惩罚。
·观看两个女人做爱。
·与某个知名女性做一次。
梦到或幻想到变态的事，与想要在现实生活中做变态的事是有天壤之别的。有些女人经常会幻想被海盗赤身裸体的绑在树上，然后被一个骑士“解救”——这也真有够变态的了，可还是不能说明有这种幻想的人就是变态。只有在你的幻想使你感到恐惧，或是你感到被迫要去做某些会伤害到自己或他人的事（以一种很恶劣的方式）的时候，那就需要当心了。
误区：男人总是幻想跟别的女人在一起。
这也是冤枉的。问一下大多数男性在性事中想的是什么，答案通常都是“性事”。对男性而言，性在那一刻意味着很多——跟女性不同，他们的思维不会真的跑去他们的“幻想图书馆”里挑选剧情。因为他们的精力全都集中在高潮上了，想要分神，也是乏术。即使他在前戏中想过别的女人，这也不能说明他对你不忠（正如大多数与幻象嬉戏的人，他们也有权为自己的性经历添些个特别风味）。这也只是性爱游戏的一部分，说明不了其它任何事情。
\subsection*{几类常见的"反性爱"因素}
生活中有很多事物会给性激情泼冷水，常见的是：尽管这也许对你来说是少有的事（你也有权保有或遵循那些反性爱观念），但如果你不去理会它们，它们也就不会妨碍到你的性生活了。
反性爱主义者
作为一个肉欲者，但愿你还没有碰到过这一部分人口。他们把“性”当作是一个肮脏的字眼——不该用于享受，也不该发生在婚姻之外。也许你就是在这种观念中成长的，如果真是这样，可就逮到令你谈性色变的罪魁祸首了。一旦有这种感觉，你就需要找个人来探讨一下你所接受的那些性信息。一位资优的精神治疗专家不但可以帮你纠正错误的想法，还能将你从那些狭隘的观念中解脱出来。
暴君情人
这种情人总是要你遵从他们的规则。有些男人自以为比你更有经验就有资格让你按他的方式做事，遇到这种男人经常就意味着你遇到了暴君情人。如遇这种情形，你要遵循的规则很简单——就是不要做你不想做的事。正如我们前面所提到的，自慰不是适合于所有人，某些动作、幻想以及某些性行为方式也是如此。而且，什么能引起你的冲动，什么会让你对性产生厌恶，这些都是你凭自己的感觉而得知的，与你是否缺乏性经验根本毫无干系。
没有激情却扮狂野
跟上面那些因素相同，被迫表演也会使你对性产生深远的厌恶感。一些杂志，甚至还有书籍的错误引导会加重缺乏热情和经验给你带来的自卑感。幸好我们的情况还没有那么糟。研究显示，大多数人都懒得去用什么“奥林匹克”式的体位，而是坚持采用传统的传教士式。
再说，我们大多数人还是不会选择在恐怖的公共场所冒险，而是安安分分的呆在卧室里“行事”。大多数女性最时髦的做法，也不过就是在MS买些色色的内衣罢了。这不能说明我们都是平淡乏味的人（想一想，谁会知道你穿那些性感内衣是要干嘛），而假扮狂野也不是对任何男人都能奏效。如果你真的有这种性情，那也最好是慢慢展现，特别是在新情人的面前，可不要过火到把人家吓跑（情趣内衣与情趣玩法的系统介绍见 B 册“情趣用品与氛围”）。
色情图片
一想到色情图片就会让你浑身上下直打寒战？还是有股欲望的冲动？实际上大家对待色情图片的态度，也是仁者见仁，智者见智。有人觉得这种东西实在可耻，只适合那些穿着雨衣的可怜男人。而其他人一看到这些就会产生激情的悸动。无论你对色情读物有何观感，那可是男性们自青少年时代起在性爱历程上的第一个港口。这就是说，没有几个你认识的男人会对顶架上的杂志或录影带不动声色。或许他们正在用这些来排遣欲望。如果你真的想知道，不妨去问问他。
“我八岁的时候跟我的朋友一起偷看色情杂志，我们觉得那些照片真是太恶心了。我还记得当时我想，‘妈呀，女人两腿之间的是什么鬼东西？我一辈子都不会接近那里’。”
——马克，20岁
你得到的答案会取决于以下几方面：（1）你的语调（2）你是否已经知道答案。如果你已经找到一个箱子，里面装满他在近十年内收藏的各种色情杂志，而你还要问他，那你的目的好像只是故意要他难堪。如果真是这样，你可要注意以下几点：
·大多数男性，也有不少女性会，会在两人性事或自慰时把色情图片用作视觉刺激。
·只要他不会强迫你看，那他还不算是过分沉溺，而且那些图片也不会有多猛烈、多恶心——这都取决于他。
·当然，如果你问他只是想跟他分享其中乐趣，那就痛痛快快的问吧。

\section{口交要点（男女都适用）}

口交（oral sex）是最被偏爱的非插入方式之一：接受方几乎"什么都不用做"，而它常比抽插更容易唤起高潮、也没有受孕风险。但"做得好不好"和"舒不舒服"，确实有讲究。

\subsection*{口交的基本认知}
口交的好处是，如果你是接受的那方，就什么也不用做，只要专心躺在那儿慢慢享受就成了。这也许就是大家如此偏好口交的一大原由吧。口交受到青睐的其它原因还包括：
·它经常会使人感觉比生殖器的抽插还要深入。
·更容易唤起高潮。
·舌头带来的感觉可是大有不同。
·有一种邪恶的感觉。
·不会有受孕的危险。
·这是引发性冲动的最快方法。
但是，如果你是施与者就需要下一番工夫了。你要注意，尽管是有很多人想以“追求性享受”为名进行口交，但不是所有人都能轻松的享受到这种乐趣。如果你不确定怎样才能把他“吹”好，以下是要注意的几点：
\textbf{你真的愿意吗？}
这是你要问自己的第一个问题，他也许期待着你来做，但如果你讨厌这种方式，就要好好找找原因了。很多女性都不喜欢口交，因为害怕做的不对。还有一些人则是担心会有什么怪味会让她们中途窒息，还有她们到底要不要把那些东西吞下去？如果这也是你的顾虑，切记，口交是没有什么规则可循，你怎样舒服就怎样做。
\textbf{卫生先解决}
即使是最放荡的情人，在可怕的个人卫生状况面前，也会立即“性致”全无。如果你一来到他的腿间就会捏着鼻子发出作呕的声音，那他的卫生状况也是有够糟糕了。不要因为尴尬就不向他说明，因为卫生状况是女性避免口交的第一大原因。如果你不敢说“你实在太臭了”，还可以采取以下策略：
·拉他去洗澡，假装这是做前戏。
·请他去健身房洗温泉。
·建议他来一次漫长的“鸳鸯戏水”。
警告:
如果他洗过之后还是臭味难挡，那他带给你的可就不只是性高潮了。（小心会感染STI或细菌类传染病。）如果真有这种情况，那你最好是使用保险套，恐怕还要为你们两人安排一次在当地泌尿生殖科诊所的约会。
\textbf{不必追求“深入喉咙”}
这只是七十年代色情电影中的可悲想法，根本不能采用（除非你是个吞剑者）。因为所有人的喉咙后侧被碰触时都会有呕吐反应，所以一旦他这样做了，能得到的唯一回报也只能是让你吐个翻江倒海。
担心会有怪味
精液是由精子、蛋白质、糖类以及其它液体混合而成，与那些传言相反，它没有多少热量。精液的味道是有点怪：有时会发咸，但也有时也会受到饮食的影响。所以，没什么值得大惊小怪的。
\textbf{口交的基本要点}
只是想象一下你低下头去的样子，就能让他为你而疯狂。为什么呢？这是因为阴茎作为男性的性器官，上面布满了神经末梢，因此才能保证他可以达到性高潮。而且，男性们通常都会说喜欢看到你给他们的阴茎施加直接刺激。如果你想跟他一样沉浸其中，不妨试一下以下方法：
做的舒服些。
没有什么是比用一种令人尴尬的体位来进行口交更糟糕的事了。所以，你一定要处于正确的位置上。他可以坐着或站着，而你则可以跪在他前方。或是让他躺在床上，那么你就可以把头放在他的大腿之间。
不要让他扶着你的头
这是女性们对口交抱怨最多的事。你需要一直控制阴茎进出的节奏和深度，所以不要让他扶着你的头来控制你。
用手开始。
用你的手来爱抚他的阴茎，他很快就会“性致”勃发了。在用嘴环住他之前，先用一只手成拳状握住他阴茎的根部，这样就能防止他向你的嘴巴猛冲，还能帮你将他的阴茎引向你所希望的位置上。如果他喜欢猛冲，就让他躺下来，你就会有更多控制权了。
舔吮。
把他当作一根棒棒糖，不只是要吮吸还要舔吮和亲吻。要充分运用你的舌尖和舌中央平坦的部分（这会给他两种截然不同的感受），试着从根部一直舔到顶端。卷曲你的舌头环住他，然后再专注于顶端。要特别注意顶端下侧那一片隆起的区域，这个被称为系带的地方也是极为敏感的性感区域。对阴茎采取的最佳动作就是用你的嘴上下套弄并且轻轻的吮吸。
“有些女孩子很怕口交和手交，她们只会敷衍了事，根本起不了什么作用，正确的做法是要抓的紧一些。女孩们，来做吧！我们真的很喜欢这样的感觉。”
——麦克斯，24岁
再加上些花样。
用些气味芬芳的润滑剂或是味道强劲的薄荷糖，都能为此增添乐趣。这不仅有助于增加口腔的湿润度，还会使他的阴茎在被你的嘴唇环住时感到阵阵暖意。
用上你的手。
担心自己的下巴会疼？那就把手也用上。当你的嘴专注于他的阴茎顶端时，用手握住阴茎根部，并且同时爱抚和轻握睾丸（不过要事先确定他是否喜欢这样）。再试着去摩擦他的会阴处（再阴茎与肛门之间一小片隆起的区域），这样会使他更加兴奋。
不要忽略他身体的其它部位。
你在给他做口交的同时，还可以去爱抚和触摸他身体的其它部分。主要集中在他的大腿和臀部，甚至对肛门也要下一番额外的工夫。口交的目的是“以口做爱”，所以，要有一次完整的经历，就必须让你的手和唇在他全身上下持续的移动。
吐出来还是吞下去。
如果你不想就不必吞下去——这不会使你背上劣质情人的恶名。记得在附近放几张纸巾，一旦他射在你口中就吐出来。试着做得有礼貌些，毕竟看到自己的情人起身跑向浴室可不是件性感的事。如果你不希望发生这种情形，就要注意观察他即将射精时的征兆。在他射精之前（他也可以及时警告你）他的臀部会出现一些反应，阴茎也会膨胀和扭曲，睾丸也会更贴近身体。如果你看不出任何生理征兆，就注意他的呼吸以及他在冲刺中速度增加的幅度。
\textbf{接受口交的常见疑虑}
很多女性都乐于施与口交，但对她们男友的“回报” ——舔阴，却会觉得怪异。一些普遍的疑虑包括：
·担心自己的生殖器有异味。
·担心他不是真的想做。
·担心自己不会达到高潮。
·担心自己的性器官外观奇异.。
好消息是，以上没有一样会是真的。尽管阴道天生有一种霉味，但并不恶心也不会有腥味（除非你是得了轻度妇科感染，那就需要检查一下了——妇科症状的识别见 D 册“妇女保健及病症治疗”，性传播感染见 A 册“性传播感染详解”）。还要注意每个月的不同时期，阴道的润滑程度也有所不同——你只需保持好个人卫生，就能保证这些都不会成为问题。而关于生殖器外观，也不必大惊小怪，每个人的生殖器，跟眼睛、鼻子、耳朵一样，外观都不尽相同。
关于舔阴，值得注意的是：对大多数女性来说，以这种方式获得高潮要比用其它方式要容易得多。原因很简单，因为这是对阴蒂的直接刺激。
对女性施与口交应包括：
·刺激阴蒂、阴道以及会阴。
·刺激大腿内侧和乳房。
·将注意力集中在阴蒂上。
·用手插入刺激。
·舌头的动作要有稳定的速度和规律性的节奏。
·需要你给他一定的指导（如果你害羞到讲不出话来，就用呻吟声来表示赞赏，或用身体的移动来帮他找到确切的位置）
·还要有一位心甘情愿的情人。
\textbf{口交常见问答}
“69”式体位不适合你。
出于各种不同的原因，很多人都不喜欢使用“69”式体位：一些人觉得这样不舒服，还有人觉得这样很难获得性高潮。女性们对此最多的抱怨是，她们发现同时施与和接受会分散精力，无法专注于高潮。而男性们也抱怨说因为他们的女友无法集中注意力，他们也享受不到正确的口交，。原因就在于，当你接近性高潮的时候根本无法顾及自己的“职责”；而如果你想专注于自己的“职责”，你就不能尽情的享受高潮。这是一种恶性循环。如果你碰到这种情况，在保留“69”式的基础上，最好每次只有一人施与、一人接受。这样，你既能随意施与，也可以尽情享受，无论如何都不会分心了。
“手上工夫”要优于“口头工夫”。
这是口交技巧中的上上之选。成功的秘诀就在于——
·不要抓得太用力或是拉扯——这可是男士们的头号抱怨。因为阴茎看似坚挺，在勃起时却极为敏感。
·问他想要你怎样做。（如果一定会擦到你手上的伤口可就惨了）。每个男人的喜好都有所不同，有一些喜欢抓得紧些，还有些喜欢轻且快的动作。让他们先自己开始，然后把你的手放在他的手上去感受一下正确的做法。
·一旦你的动作上了轨道，掌握了他所喜好的强度和节奏，就试着用手从阴茎的根部一直上移到顶端的隆起。如果还想加点花样，就在顶端和根部稍加力道。
·用些润滑剂。这样既能给阴茎增加湿度，又能保持动作的顺滑。可以用KY胶或唾液。
·一定要沉浸其中——如果做的时候你摆出一副心不在焉、呵欠连篇的样子，不仅不会加快他的高潮的到来，连你自己的情欲恐怕也要烟消云散。
·在他临近高潮的时候就要放开阴茎，因为这时阴茎的顶端会跟阴蒂一样敏感。
你讨厌在口交时阴蒂受到过多“关注”。
这也是很常见的问题。因为高潮前，阴唇的肿胀有时会使阴蒂收缩不见，有些男性就会过分专注于寻找阴蒂。而这时阴蒂上所有的神经末梢都会处于红色警戒状态，所以阴蒂就会变得极为敏感，也就极易受到伤害（有些男性的阴茎顶端也是如此）。碰到这样的情况，要让你的情人退后点、温柔些。
不知道该拿包皮怎么办。
那就什么也不要做。然而，将包皮拉紧包住阴茎的柱状部分会给男性带来额外的愉悦感。最好的方法就是把手放在阴茎下距根部2.5厘米（1英寸）的地方，将手箍成环状像下拉，这样阴茎的上部就会稍稍露出。还可以试着舔吮包皮的四周然后再专注于阴茎的顶端。
“大多数男人的口交方法都不正确，他们只会一头扎下去猛流口水，就好像要舔出一杯牛奶来。真是讨厌！而且他们似乎认为，不管他们在下面做的怎样，你都要对他们感激不尽。”
——劳拉，24岁
你的男友拒绝口交。
他拒绝为你做口交的借口有很多种，像是“下面太暗了，我看不到自己在做什么”，“这太奇怪了”，还有最可笑的——“会有异味”。其实，男性拒绝口交的真正原因多半是他们害怕自己做的不对，除此之外根本没有别的缘由。
这也就是说，如果你想要你的情人为你做口交，最好的办法就是直截了当的告诉他你喜欢这种方式，还要告诉他应该怎样做。
这会立即奏效，因为（1）他不可能曲解你的要求；（2）你一字一句的指导定能让他心领神会。尽量多用一些视觉性的描述，男性往往要通过这种方式才能搞清应该怎么做。一定要让他保持你所需要的持续且有规律的节奏，动作要柔，力道要轻，还要告诉他要用怎样的性刺激来将你推向高潮。
\section{前戏与初次性爱}
\textbf{前戏}
不需要我来教你怎样做爱吧？你只要把衣服脱掉，让他进入然后在你体内冲刺几下不就完事了，不是吗？
但愿你的反应是大叫一声“不对”。
性行为看似简单，但要做好也实属不易。因此现在最主要的就是先要综合一下你现在对性的全部了解。回想一下那些生理基础知识、有关自慰、精神挑逗、前戏、性感地带以及口交的内容。把这些全部综合就能指引你走向更美满的性生活，甚至是体会高潮中的更多美妙滋味。如果你太痴迷于抽插，即使有性生活也不会有持久的兴奋感。想要来些给懒女孩们的速成建议吗？那就往下看。
做好前戏需要有——
·在插入前需要至少多于五分钟的前戏。
·前戏中要有大量的拥抱和亲吻。
·互相爱抚、摸索彼此的身体。
·要覆盖大部分的性感地带（上文已逐一列举体侧、耳垂、手指、脚、大腿、背脊、腹部等敏感区；系统版见 B 册“非插入类行为的愉悦技巧”与“前戏完整剧本”）。
·要有某种形式的口交。
·聊几句（但不宜过多，全程解说可不是性感的举动）。
·用上你的手、指、舌、唇以及你能想到的任何东西。
\textbf{让前戏更好}
将你花在这一阶段的时间延长一倍。
如果你的男友是个喜欢“速战速决”的人，那你就要让他专注于他在进行的事情，而不是进行到哪里了。大多数男性急着从开始一下子猛冲到顶点，这并不是因为他们急切的想要进入，而是因为他们面对前戏有些无助，捏了几下你的乳房之后就毫无头绪了。要帮助他，就要告诉他前戏不只有接吻、逗弄乳头和笨拙的摸索。要开始并延长这个阶段，靠简单的亲吻就成了。
a． 脱衣动作要性感。
宽衣解带是绝对惹火的举动。让他抚摸你的胸罩十分钟，而你把头罩在上衣内，用手隔着他的衣服捉住他的阴茎，这些都是挑逗激情的“杀手锏”。所以，不要急着脱光，衣服要比你想象的性感多了。来段脱衣舞，让他观赏你轻解罗衫，他立刻就会“性致”昂扬。或者换你穿着内衣躺在床上，静静观赏他的表演，那就更棒了。
b． 问他喜欢怎样做。
男性在性沟通方面可是臭名昭著的无能，所以如果你不问他，他就会用他那些恶劣的方法来让你做他想要你做的事。不能让他随意就把你的头向他下身推去，更不能让他抽打你的身体，这样实在很杀风景。相反，你必须让他自己去要求/争取他想要的。这不仅能提升他的“性致”，还能增加这一时刻的性感价值。
c． 不要期待他会读懂你的心思。
与以上相同，“问”是能让你获得满足的关键所在。是不是需要更多的口交？少一些亲吻？多一些舔吮？那么就告诉他你有多么渴望他。如果不能用“言传”，那就用你整个身体来进行“身教”，挑动他的欲望并且给他正确的指引。用上他的手指、唇舌，让他试着用手指进入你，这样很快就能把你带向高潮。
d． 拒绝作假。
如果你觉得有什么不适之处说出来，否则你一生恐怕都要默默忍受这种烂技术了。有些男人总是会遵守“你不说，我不改”的原则，这就是说，如果你不跟他说明你讨厌他舔你的耳朵，他就会一直舔下去直到让你抓狂。
e． 前戏不可太烂。
最后，要记着，前戏固然重要，但如果做的太烂，那就另当别论了。软弱无力的口交、拍打阴茎或是半烘烤式的性按摩都不是合格的前戏技巧。又没主意了？那就再看一下前戏与口交的系统方法（见 B 册“前戏完整剧本”、“口交详解”），看看怎样才能使你的性生活活力四射。
\textbf{初次性爱}
童贞是一个惯用的侮辱性的词汇，除了在那些传言之中，童贞根本不能说明你是个怎样的人，只能说明你还没有过实质性的性行为。处女之身并不意味着你会是个无能的情人，也不能说明你不会获得高潮（让我们面对现实，你也许已经不是处女，但在其他性领域毕竟还是个初学者），或是比那些有多年床上经验的人更为出色。因为大家都在做（据英国国内数据统计署的数据，大约有四分之一的英国女孩在十六岁之前就不是处女了，而大多数择会等到19岁左右）。那么对那些迫不及待要摆脱处女身份的朋友，以下是你需要了解的一些内容：
接受现实也许要比想象中困难。
很多第一次做的情侣们都会对性爱的严酷现实产生种种疑问，这只是因为他们对处女膜的破裂还有诸多疑问。处女膜，就是部分覆盖住阴道口的薄膜组织。有些女性的处女膜在有性行为之前就已经破裂了，儿时的剧烈运动和使用卫生棉条很容易就能使处女膜破裂，而在很多情况下，处女膜在第一次性爱尝试中却能完好无损。
如果你的男友不能进入或是感觉有什么东西阻挡了他的前进，那他就是触到处女膜了。如果你感到一阵疼痛，而且看到有血流出，这就意味着你的处女膜破裂了。
有助于性事的最佳方法：
·开始时要采用传教士体位，这对“第一次”来讲是最佳体位。传教士式大致上就是俩个人躺下来，让他位于你的上方。
·当他第一次进入你的时候，用你的阴道肌肉将他向下推挤，就像要小便的感觉一样。
·接着，他会慢慢插入，然后在你们俩人都放松下来时，更加深入，直到完全进入。
·这会使你的阴道不断扩张以适应他的存在，而处女膜的破裂使你最终得尝禁果。
记得要提醒他这是你的第一次。
童贞或性爱都不会在你的头上挂个霓虹灯招牌来宣告你是否做过。“他看不出来吗？”如果你的处女膜已经破裂（见上），答案是——他的确看不出。话说回来，你又为什么要说谎呢？性爱可是件大事。如果你不想它成为一次美妙的经历，硬装出一副事事明白的样子，你就不能体会到个中乐趣了。
如果他是处子，要对他好一些。
每个人都会有第一次，如果你要跟一个处男上床，切记，男性跟我们女性不同，他们在第一次时前戏只需一笔带过。如果你花了太长的时间来挑逗他，再把他带上高潮边缘，他就会达到无回复点，恐怕在进行到实质性部分之前他就射精了。所以，你最好——
·选择一个能让你掌控全局的体位，例如男下女上式，采用这种体位的话，在他过早射精时，你就能及时起身。
·告诉他你喜欢怎样的触摸，怎样做会让你有美妙的感觉。因为他对这些了解不多（当然，除非你也只是处于尝试阶段，那你们两个可就是半斤八量了）。
小心会痛。
处女膜破裂时会有点疼，但疼痛很快就会消失。如果性事中还是很痛，那就需要有更多的前戏，还要注意以下几点：
·不要痴迷于插入。如果你做不到，想靠蛮力继续下去是很容易，但也会让你疼痛难忍，要再深入就更不可能了。你最好先暂停一下，在前戏、口交和亲吻上花上更多的工夫，完全放松后再试一次。
·尽量放松，当我们事先认定什么东西会很痛苦或很艰难的时候，就会潜意识的紧张起来，这样会导致阴道肌肉紧闭形成阻碍。一旦出现这种情况，无论用什么润滑剂也无济于事了。
时间就是一切。
一旦你们俩人都学会该如何挑起对方的激情，你不仅会在头脑中有所领会，你的骨盆区域也会有一种膨胀的感觉，这是你已经做好准备的信号。时间的积累还能在其它方面对你有所帮助，例如高潮（如何引导与达成高潮见 B 册“享受高潮”、“高潮引导与边缘控制”），所以如果第一次不能带来焰火或爆炸般的性感感受也没什么好奇怪的。就像第一次约会一样，有一点尴尬和难为情都是在所难免。而且这也不能代表以后的情形。性爱这种事，也是做得越多，就能做得越好。
选择避孕法。
民间有一种说法，“第一次”的时候，即使他采取“及时撤退”的方法都不能避免受孕。事实恰恰相反，如果你不想成为英国每年十一万意外受孕的女孩之一，就一定要采取某种保护措施。无保护性爱还可能使你染上性传播类传染病，所以一定要选择一种能起到屏障作用的避孕法，即使用保险套。（关于避孕的详细内容见 A 册“避孕方法详解”）。
“我在某些地方读到男性每六十秒钟就要想一次性。如果你跟我和我的朋友们待上一晚，你就会见识到，我们绝对能打破这个纪录。”
——露西，24岁
不要急于追求刺激。
在你开始像火堆上旋转着的狂舞者一般翻滚、扭曲你的身体之前，还是先掌握好全部的基本体位吧。你也许渴望着去体验性爱狂野的一面，但如果连最简单的都还不能做好，就更不用提那些高难度动作了。然而，也不必整晚都停留在传教士式——试着起身，站着做，向后倾斜，亲吻他，逗弄他，抚摸他。当你心满意足之后稍作休息，过一会儿再把他叫醒，继续“第二回合”。
不要期望过高。
美妙的性爱是需要时间的，每次你跟一位新男友开始一段新的“关系”，通常需要做过几次才能得其要领。所以，即使没有获得高潮，还有一些不满之处，甚至没有丝毫享受的感觉，你也不必大惊小怪。这并不能说明你是性无能或冷感者（骇人的虚假字眼，没有任何意义）。这只能说明你需要更多时间来学习性爱的基本内容。先回到性爱的插入前阶段，专注于亲吻和前戏，感觉对头了，再来发起第二次冲击吧。
\textbf{保持性生活的新鲜感}
做爱、性交、交和、交媾——这种游戏的名字就如同它的做法一样举不胜举。幸好我可不是在提倡倒吊在房顶上或是去Trantric，也不在怂恿你在邻居面前做爱。我的建议只是，不管你有过多少经验，技巧的多样化才是性生活美满的关键。
美满性生活
性生活的多样化是指：
·前戏技巧的变化。
·变换体位。
·尝试各种抚摸方式。
·为你的性生活增添乐趣。
·考虑一下“速战速决”的做法。
·考虑一下“持久战”的做法。
·告诉他哪里做得好。
·告诉他那里做得不好。
·寻求改进个人技巧的建议。
·到卧室之外的地方去。
性爱的麻烦在于，我们中的大多数人只要遇到了某人，坠入爱河，经历了一段极度兴奋的性冒险，之后就会急流勇退，停在一条固定的套路上，这时的性爱就会成了“在下雨的午后”才做的事。然后，我们中的大多数人就会依靠性的回忆生活着，不断的提醒着我们的伴侣我们曾有过那样非比寻常的性爱，或是还想要再次尝试。就这样，直到有一人觉得无聊，转身去寻找另一个会停下来跟他聊天、上床的人。
正如你的上班路线一样，如果你每次性事都按同一条路线，那你就会进入一种自动驾驶状态，因此就会上演上面的那出悲剧。这就是说，在他开始用鼻子磨蹭你的颈部时，你就会开始想起要清洗厨房地板，或是开始想象你是跟布莱德·皮特在一起。
所以，即使你的“基本功”已经上升到艺术的层次，但为了你和你的伴侣获得更多的享受，你仍需要重新修炼并扩展你的性技能。好消息是，这很容易做到，比你想象的要省力多了。

\section{常见性体位}
你知道多少种性体位呢？三种？三十种？还是三百种？不管你知道多少，其中有几种是你做得到的呢？而你又为什么要用这些做法呢？如果你都有所了解，那么最好是从中选出五种体位，并且要固定在这五种上。因为（1）既然你选择了这五种，就说明它们功效不凡（2）这些也是你能做得到（3）将双腿绑在头两侧这种高难度动作，你也做不了。这些都是显而易见的——但或许你还想了解一下，仅用一些小动作或身体的简单移动至多能变换出多少种体位，还有用平淡无奇的方法又怎样才能变换出全新的方式。下面就让我来告诉你吧：
传教士式
这是一种最为传统的体位，如果那些调查结果还算可信的话，现在仍是人们最常采用的体位。它虽然不是最能令人兴奋的种类，但如果动作准确，不仅能使你享受到性爱的美妙，还能燃起你们全部的激情。
做法：
呈仰卧姿势，双腿分开，让他躺在你身上，好了——这样就可以让他进入了。
如何变化：
·还照上面那样做，但是在他进入时要将双腿伸直、夹紧，这样会给你的阴蒂带来强烈刺激。
·改变一下角度：用双腿环住他的腿。这也有助于你控制他的冲刺。
·在你的臀部下方垫一个枕头以使你的身体后倾，他就可以进入得更深。
·让他躺在你上面，不过要移动到你身体上方的较高位置上，这样，他就要用双臂来支撑自己，也就不会把你压扁了。
有利于：亲密的“第一次接触”，对维持勃起有困难的男性。而且，对那些不能控制射精、容易“乱开枪”的男性尤其有利。
男下女上式
恰如其名——就是指你位于他的上方。这可是大多数女性的至爱，因为采用这种姿势，她们就可以掌控全局而且能保持较深的插入。男性们也喜欢这种体位，因为他们几乎什么都不用做。如果你的伴侣平日占尽主导地位，用这种方法就可以帮你控制他的插入深度。如果他容易早泄，这种体位也会对你有所帮助——当他看起来像要到达高潮时，你可以及时放慢速度。
做法：
让他呈仰卧姿势，双腿略微分开，你爬到他身上，跨坐在他的臀部两侧，让他的阴茎进入你的阴道。
如何变化：
a． 当你进入律动的节奏时，身体前倾，用他的骨盆区域摩擦你的阴蒂。这样你很快就能获得高潮。
b． 倾向他的胸部，翘起你的骨盆以获得G点高潮。
c． 如果你坐在他身上向下落去，你会使他更加深入。
d． 不要只是让他来冲刺，而要靠你自己上下移动来达到这种体位的最大功效。
e． 最后，可以试着以这种体位背对他，然后身体前倾，使他更加深入。
有利于：想要控制插入深度和冲刺速度的女性。还有利于那些懒男人们。而且，在抽插时会感受到绝妙的阴蒂刺激。
小狗式体位
这是“从后方进入的性爱方式”，与“进入后方的性爱方式”相对（见84页“肛交”）。因为这是犬类“恩爱”的方式，所以得此名称。男性们很喜欢这种体位，因为会有一点点邪恶的感觉，同时给予他们对进入与冲刺的控制权，能使他们更为深入。女性们喜欢这种体位的原因却在于它经常能刺激到G点，而且男性们也更容易接近她们的乳房、臀部和阴蒂。但这种体位也有一种缺陷，它会使很多女性感到疼痛。因为插入过深也就意味着阴茎有时会触到子宫颈，因此会产生疼痛。
做法：
用双手、双膝支撑身体，让他在后面抱着你的臀部，由后方进入。
如何变化：
a． 放低身子，用手肘支撑，这样可以改变他进入的深度。
b． 跪在床边，将你的躯干放在床上。这样会使你体会到使用这一体位的轻松之处——不必在他冲刺时支撑自己的身体。
c． 对那些似乎额外强壮、想要有一次边缘体验的朋友，可以尝试将身体向床上倾斜，双臂曲起来支撑身体的重量（把头放在枕头上）。让你的男友站立着接近你，将你的双腿抬起环在他的腰上，然后从后方进入。
有利于：更深的插入；使性爱有一种更加邪恶的感觉；能刺激到G点。
“各种各样的体位和动作的确挺棒，不过，没有什么能比得上跟自己的梦中情人用‘一上一下’的简单方式来做一次。”
——丹，25岁
站立式体位
有利于“速战速决”式的性爱，但如果你们两个身高、体重相差太大，那就不太适用了。尽管站立式体位在电影中看起来很简单，真正实践起来却很困难。因为你要站着找到正确的进入角度，经常就会经受不住重力的考验，破坏这一激情的时刻。
做法：
你需要处于适合插入的高度上，这往往就意味着你要将一条腿盘在他的腰上。
如何变化：
a． 让你的男友抱起你并且要扶住你。
b． 如果这样太困难了，想想找些什么东西能让你站在上面（不是很性感的举动，但的确有效）。
c． 穿上高跟鞋——性感，而且还能将那些不能穿着走路的鞋子善加 利用。
有利于：喜欢尝试不同事物的情侣们；还有利于增进两人的亲密感；以及“速战速决”。
并行式体位
如果你追求的只是亲密的而不是激烈的性爱，这个体位就很适合你了，因为这种体位的运动幅度较小，激烈程度有限。这种体位的优点在于，整个过程需要很多“手工”。
做法：
两人并行躺着，将你的大腿放在他的腿上，他就能从更好的角度进入了。
如何变化：
a． 将身体弯曲成勺子的形状（转过身去，让他贴住你的背），就会变成从后方进入的侧体位。
b． 弯曲膝盖环住他，让他以面对面的侧体位进入得更深。
有利于：那些避免过度深入并且追求亲密的慢节奏性爱的朋友。
肛交
肛交在绝大部分女性的愿望列单上排名都比较落后，因为大多数女性仍把它看作是禁忌的体位，而且觉得使用这种体位做起来会很痛/会染上些恶心的东西/最糟糕的是她们会受到粗暴的对待。但这并不意味着没有人去尝试。根据著名的美国金赛机构的调查，大约有43%的女性曾尝试过肛交，但只有不到30%的人觉得有愉悦感。肛交也是一种高风险行为（想想HIV吧），所以做的时候一定要使用强力的保险套来保护自己，而且一定要他慢慢来。肛门区域的组织都十分脆弱且没有弹性（不像阴道），因而很容易撕裂。
做法：
肛交对大多数女性来说都是很痛苦的，因为要放松肛门肌肉（要同时放松两片括约肌）是很困难的。要克服这一点，最好的方法就是多涂些润滑剂，先将一只手指插入静止一分钟，然后向内推挤。这样有助于放松括约肌，阴茎进入时也就比较轻松了。我还要建议你使用以下这些最佳体位：
a． 仰卧，在臀部下方垫一个枕头。你可以把手放在大腿下面把大腿向后抬起以加大身体弧度，使他更容易进入。
b． 或者坐在他身上，向下将阴茎压入。
c． 或者采用从后方进入的体位。
有利于：那些喜欢大胆行事的人。
\section{改善性生活的小建议}
亲吻
听起来的确显而易见，而女性们对性爱最大的不满之处就是她们的男友们总是忘记亲吻她们，对她们的性感地带也是蹭蹭了事。所以让我们从基础做起，回到热吻这一阶段——你不会后悔的，因为研究者们发现接吻不仅会使性爱持续较长时间，还能增强高潮、增加温情，以及增进性欲。
不要拿你的性生活当谈资。
大多数男人把性披露当作最大的出卖，（例如，“天呐！我那位的‘家伙’不像香肠像什么？”）。如果你还想受益于他的床上魅力，就要对你们的性事守口如瓶。
尽量不要分心。
忘记你的同伴们在做什么，也不要去想你牺牲了多么宝贵的饮酒时间。跟人们的普遍想法恰恰相反，“春霄一刻值千金”，忘掉那些恼人的烛光音乐也是值得的。
关掉电视。
电视最近被评为最大的“性欲杀手”。在床上看电视显然会将你催眠，减少你对性的敏感度继而消灭你的性欲。
掌控自己
用上你骨盆底部肌肉（小便时使用的肌肉）。每天做二十次收放练习，你很快就能控制好自己的身体，来将他夹紧。
捕猎G点游戏
如果你告诉他，他要寻找的是什么，那么难以捉摸的G点也就不需要再难以捉摸下去了。只需将一只手指插入阴道内，再向上移动直到触摸到阴道前壁，就能找到这一点。
身体对身体的按摩
让他趴着而你则躺在他身上，然后用你的身体在他的身体上上下摩擦。如果想要更具挑逗性的动作，就把手滑向他前侧。
不要太过激
特别是你在下面探索时，如果你又拉又拽，他可不会感激你。
对他进行口交
如果你要对他进行口交就要把功夫做足。当你抬起头呼吸时，不要露出奇怪的表情，不要屏住呼吸，而且千万不要表现出一副帮了他大忙的样子。
口交时速度要慢。
要表现的像一个沙漠绿洲中饥渴难耐的女人，你就永远不会让他的欲望熄灭。开始时还要用上舌头扁平的部位，而不是只用舌尖儿。
上床时带只香蕉。
可不是为了你所想象的原因。香蕉富含维生素B, 所以吃香蕉可以促使血液流向性器官，从而有助于增进性事和高潮。
用冰块来“增温”。
专家建议，在接吻或口交时可以在口中含一块冰。会不会太凉了？那就试一下在插入前将一小块冰块放入阴道中。
避开家居用品
相信除了把胡佛电动吸尘器插入你的（或他的）肛门之外，还有更好的方法来让你的父母为你感到自豪。
注意选择食物
你的阴道可不像你的嘴巴一样，不是所有你想塞的东西都能往里塞，（希望你不是这样）。某些食物（特别是富含糖分的），会使阴道肌肉恶化。这也就是说，要小心那些恶心的肿块和疥疮了。
注意他阴茎上的斑点
阴囊区域就如同人的面部，很容易感染上一些小斑点。（特别是因为阴毛比较卷曲，就很容易向后弯扎入毛囊，就像是在比基尼裤边儿上那些向内生长的毛发）。有时受感染的毛囊还会在腹股沟处的温暖潮湿的环境中长成疖子。
小心对待他
勃起的阴茎很容易受伤，所以把自己丢到床上时一定要小心。如果以你听到了一声惨叫，那你可就“有幸”目睹男人阴茎剧痛的情景了，这时要立即带他去治疗。如果他的阴茎破裂，可是需要用夹板来固定的。
不要对“失约”大惊小怪
与一些普遍的说法相反，勃起并不受他的主观控制。所以疲劳和压力往往会导致勃起发生问题。一些数据也确实显示出有七分之一的男性（16岁以上）一年之中至少有四次勃起问题。
让他把你绑起来
有利于控制那些不肯放弃主导权的怪人。但切记，如果你是被绑的那个，你可以在他向你做一些下流的事情时装出一副被动与无助的样子，或许可以幸免于难。
忘记那些平常的性数据
不要在意那些调查，或是你朋友讲她一周X次的性事。性爱并不是竞赛，关键在于你能从中得到享受，而不在于你做了多少次，这一点很重要。
试一下拉他的阴囊
这可不是什么“橄榄球”式体位，而是一种延长做爱的技巧，就是轻拉睾丸周围的皮囊。通常这样可以延长人的性欲，因为这意味着男性要等到阴囊收缩时才能射精，他也就会更加接近你的身体。拉的方法——要用大拇指和食指成环状，轻轻捉住阴囊的囊袋（就在睾丸上方），向下轻拉。
禁区试验。
尽管有时会很可怕，但这种试验可以增强你的性自信，帮你克服性障碍，从而改善你的性生活。如果你有勇气（而且想要）去做，而且想了解是否适合自己，那就大胆尝试一下吧。
穿着衣服做爱
很怪异但也很真实。一项性研究表明，64%的女性喜欢与衣衫半解的男性相互爱抚。
找到你的阴道盲管
不管你信不信，就在骨盆下方有一条布满神经末梢的极为敏感的通道。这就是所谓“盲管”。专家提示，如果你在男友开始冲刺时，收紧腹部和PC 肌（控制小便流出的肌肉），就会触到这一区域。
千万不要提起你的前男友
特别是在你们两个都赤身裸体的时候。
不要把性爱看得太重
你要明白，它只是一种享受！

\section{本章小结}

\begin{itemize}[leftmargin=2em,itemsep=0.3em]
  \item 男女生殖器官的结构决定了各自最敏感的"开关"不同——\textbf{女性主要是阴蒂，男性主要是龟头/系带}，但两者的会阴、乳头、颈部都同样敏感。
  \item 荷尔蒙驱动性欲，但"看到就涨"是夸张说法——\textbf{睡眠、运动、压力、关系质量}才是睾酮水平的主要变量。
  \item 月经周期中雌激素/孕激素的波动会影响女性性欲，但\textbf{个体差异巨大}，不能"对号入座"。
  \item 自慰是最常见、也最健康的性表达之一；它\textbf{不会伤身、不会降低高潮几率}，反而帮你认识自己的身体。
  \item 性幻想是大脑的正常"启动器"，\textbf{不代表你是变态}，也不代表你真的想去做。
  \item 口交是安全、容易唤起高潮的方式；\textbf{卫生先解决、节奏要稳、不必追求"深入喉咙"}。
  \item 前戏是多数女性达到高潮的\textbf{必要条件}，至少 5 分钟起，覆盖多个敏感区域。
  \item 性爱对健康的好处多为\textbf{相关性证据}，别把它当处方；真正影响健康的还是整体生活方式。
  \item 体位不是竞赛——\textbf{你能做的、舒服的、对方也舒服的}，就是最好的。
\end{itemize}

\chapter*{结语}
\addcontentsline{toc}{chapter}{结语}

性健康不是一门靠"顿悟"的学问，而是可以系统学习、理性准备的知识。
愿这本小册子帮你建立起安全、自愿、知性的观念：既保护自己，也尊重他人。
当你需要更深入的专业内容时，请参阅本书“性医学基础：医学维度的速查”章，
或直接向正规医疗机构的专业人员求助。

\vspace{2em}

\noindent\textbf{主要资料来源：}
\begin{itemize}[leftmargin=2em]
  \item World Health Organization (WHO). \emph{Family planning/contraception methods}; \emph{Medical eligibility criteria for contraceptive use}, 6th ed.
  \item U.S. Centers for Disease Control and Prevention (CDC). \emph{Primary Prevention Methods — STI Treatment Guidelines}; \emph{Condom Use: An Overview}; \emph{The Lowdown on How to Prevent STDs}.
  \item 国家卫生健康委员会及中国疾病预防控制中心关于 HPV 疫苗接种的相关公告（以最新发布为准）。
  \item European Association of Urology (EAU). \emph{Guidelines on Male Sexual Dysfunction}; ICD-11"早期射精"定义（第14章）。
  \item 科技日报、中国新闻网：酒店针孔摄像头识别与防范专家建议（藏在镜子背面、插座、空调出风口、烟雾报警器等处；手机"红点法"局限；手电查反光、查 WiFi 陌生设备、物理遮挡）（第15章）。
  \item 全国人大《中华人民共和国民法典》第 1033 条、《治安管理处罚法》、《刑法》第 283 条关于隐私空间与窃照器材的规定（第15章）。
  \item Non-consensual condom removal ("stealthing"): PMC\,6306234（墨尔本性健康诊所研究）；Drouin 等，\emph{Psychology of Violence} (2024)（第15章）。
  \item CDC. \emph{Clinical Overview of Urinary Tract Infections}（性行为引入细菌、女性尿道较短风险更高）；\emph{Condom Use: An Overview}（第16章）。
  \item WHO. \emph{Sexual health}（2006 工作定义：积极、尊重、愉悦且安全、免于胁迫）；\emph{Redefining sexual health for benefits throughout life} (2022)（第16章）。
  \item W.\,H. Masters \& V.\,E. Johnson. \emph{Human Sexual Response} (1966)：兴奋期/平台期/高潮期/消退期四阶段模型（第17章）。
  \item WHO. \emph{Comprehensive sexuality education}（科学准确、适龄的身体与关系知识，帮助保护健康、尊重他人）（第17章）。
  \item CDC. \emph{How to Prevent STIs}; \emph{About STI Risk and Oral Sex}; \emph{How to Use a Dental Dam}; \emph{Condom Use: An Overview}（屏障工具、换套、口交膜与杀精剂禁忌）（第18章）。
  \item CDC. \emph{MMWR Recommendations—Opportunistic Pathogens}（避免口--粪暴露、换套、手部清洁）；WHO. \emph{Condoms}（直肠无自然润滑、肛交需润滑）（第18章）。
  \item CDC. \emph{Appendix A: Known or Possible HIV Exposure Scenarios}（接吻、相互手刺激属可忽略风险，nPEP 不常规推荐，除非可见血 + 黏膜破损）；\emph{How HIV Spreads}（HIV 仅经特定体液经黏膜/破损皮肤传播）（第19章）。
  \item CDC. \emph{Preventing HIV with Condoms}（安全套用于共用器具高效防 HIV 及部分 STI）；\emph{How to Prevent STIs}（皮肤接触传播）（第19章）。
  \item 女性高潮与阴蒂刺激依赖：大样本与综述研究（多数女性需直接阴蒂刺激达高潮；仅阴道插入缺少阴蒂刺激时高潮率明显更低）——综合 Masters \& Johnson 性反应模型与后续大样本调查（第19章）。
  \item CDC. \emph{How to Use a Dental Dam}; \emph{About STI Risk and Oral Sex}（口--口/口--生殖器黏膜传播、口腔屏障工具、有活动性溃疡/疱疹时避免）（第19章）。
  \item M.\,B.\,Rosenberg. \emph{Nonviolent Communication: A Language of Life}（观察--感受--需要--请求四步框架，请求与指责的区分）（第20章）。
  \item WHO. \emph{Sexual health}（积极、尊重、愉悦且安全的性体验，免于胁迫与暴力）（第20章）。
  \item American College of Obstetricians and Gynecologists (ACOG); J.\,Trussell 等. \emph{Typical-Use and Perfect-Use Contraceptive Failure Rates}（各方法典型/完美使用失败率分层）（第21章）。
  \item WHO. \emph{Medical Eligibility Criteria for Contraceptive Use}, 6th ed.（口服药适应证/禁忌、避孕方法选择）（第21章）。
  \item CDC. \emph{Emergency Contraception}（左炔诺孕酮 72h 内最佳、乌利司坦 120h、铜 IUD 事后 5 天内最可靠）（第21章）。
  \item CDC. \emph{Pharyngeal Gonorrhea/Chlamydia}; \emph{About Hepatitis A and B}（有肛交者建议接种甲肝/乙肝疫苗）（第22章）。
  \item WHO. \emph{Condoms}（直肠无自然润滑、肛交需充分水基/硅基润滑；油基破坏乳胶）（第22章）。
  \item CDC. \emph{Clinical Overview of Urinary Tract Infections}（女性尿道短、性行为引入细菌、事后排尿/多饮水降低 UTI 风险）；\emph{Vaginal Douching}（阴道冲洗破坏菌群，不推荐）（第22章）。
  \item W.\,H.\,Masters \& V.\,E.\,Johnson. \emph{Human Sexual Response} (1966)（唤起期"tenting"伸展、高潮非必然）；ACOG 关于 \emph{Vaginismus / dyspareunia} 的临床说明（剧痛就医指征）（第23章）。
  \item WHO. \emph{Condoms}; CDC. \emph{Lubricants and Condoms}（水基/硅基可配乳胶套、油基破坏乳胶）（第23章）。
  \item WHO. \emph{Human papillomavirus (HPV)}（HPV 与宫颈癌/口咽癌关联、疫苗可防；HIV 经口传播极低的判断与\ref{chap:sti}一致）（第24章）。
  \item CDC. \emph{PrEP (Pre-Exposure Prophylaxis) for HIV Prevention}（高危者 PrEP 显著降低 HIV 感染风险）（第25章）。
  \item EAU. \emph{Guidelines on Erectile Dysfunction / Premature Ejaculation}（ED/PE 循证处理、PDE5 抑制剂需处方、行为法）；NICE / AUA 关于 \emph{dyspareunia、vaginismus} 与 HSDD 的说明（第26章）。
  \item NHS. \emph{Is it safe to have sex in pregnancy?}；Tommy's \emph{Sex in pregnancy}（孕期性生活一般安全、需避免的情况：前置/低位胎盘·破水·多次流产·不明原因出血；推荐并排侧卧、水基润滑；口唇疱疹避免口--生殖器；无证据表明性交会引产、假宫缩非临产；28 周后避免长时间仰卧）（第25章）。
  \item 阴道瓣（hymen）解剖与初次性行为：阴道口黏膜褶皱、中间有孔、形态多样；初次不一定出血或疼痛；疼痛/出血多源于紧张·润滑不足·动作过快（科普共识）（第25章）。
  \item NHS. \emph{Vaginal oestrogen / HRT}（更年期局部雌激素低剂量、不增加乳腺癌风险、无需配孕激素、可按需长期使用、约 3 个月起效；睾酮仅专科特定情况；停经后避孕时限 50 岁以下 2 年/以上 1 年、55 岁后极罕见）（第26章）。
  \item WHO. \emph{Family planning/contraception methods}（安全套唯一既避孕又防 STI、现代避孕不致不孕、青少年可安全使用）；\emph{Medical eligibility criteria / Oral contraceptives}（产后避孕时机：纯母乳 6 个月/部分母乳 6 周、迷你 pill 产后即可）。LAM 哺乳闭经避孕法三条件（纯母乳·月经未恢复·6 个月内）（第26章）。
  \item DSM-5（APA, 2013）关于 \emph{paraphilia 与 paraphilic disorder} 的区分（性偏好不再默认病理化；仅显著痛苦/功能受损方为障碍）；Indiana University kink 研究（拍打 30\%、dom/sub 22\%、束缚 20\%、鞭打 13\%）与 41 国 31.7 万人调查（约 20\% 用过眼罩/束缚）（第31章）。
  \item 颈动脉窦综合征/反射（心脏抑制型 50--70\%、血管抑制型 5--10\%、中枢型；60 岁以上晕厥患者 17.6\% 相关、70+ 达 25\%；触发：扼颈·勾肩·亲吻颈部·衣领过紧·颈部按摩）（腾讯医典/央广网）；ScienceDirect 系统综述（2010，48 例 CiRCA 案例：纯反射单独致死罕见，多合并脑缺氧）（第31章）。
  \item 催产素/多巴胺机制综述（OT 峰值后回落；OT 释放被触摸与生理支持增强，为 aftercare 提供依据）（PubMed 39827220 / 20584551）；WHO \emph{Sexual health}（积极、尊重、愉悦且安全的性体验，免于胁迫与暴力）（第31章）。
\end{itemize}

\begin{center}
\itshape 本文为健康教育科普，不构成医疗建议；具体诊疗请遵医嘱。
\end{center}

\clearpage

\end{document}
